\documentclass[11pt]{article}
% Shared preamble for manuscripts-v2 papers. Load with % Shared preamble for manuscripts-v2 papers. Load with % Shared preamble for manuscripts-v2 papers. Load with \input{../common/preamble.tex}
% immediately after \documentclass[11pt]{article}.
\ifdefined\pdftrailerid\pdftrailerid{}\fi
\usepackage[T1]{fontenc}
\usepackage{lmodern}
\usepackage[letterpaper,margin=1.02in]{geometry}
\usepackage{amsmath,amssymb,amsthm,mathtools,booktabs,array,microtype,enumitem}
\usepackage{tikz}
\usepackage{xurl}
\usepackage[colorlinks=true,allcolors=bluegreen]{hyperref}
\definecolor{bluegreen}{RGB}{22,83,91}
\usepackage[capitalize,nameinlink]{cleveref}

\theoremstyle{plain}
\newtheorem{theorem}{Theorem}[section]
\newtheorem{lemma}[theorem]{Lemma}
\newtheorem{proposition}[theorem]{Proposition}
\newtheorem{corollary}[theorem]{Corollary}
\newtheorem{conjecture}[theorem]{Conjecture}
\newtheorem*{theorem*}{Theorem}
\newtheorem{maintheorem}{Theorem}
\renewcommand{\themaintheorem}{\Alph{maintheorem}}
\theoremstyle{definition}
\newtheorem{definition}[theorem]{Definition}
\newtheorem{example}[theorem]{Example}
\newtheorem{question}[theorem]{Question}
\theoremstyle{remark}
\newtheorem{remark}[theorem]{Remark}
\crefname{maintheorem}{Theorem}{Theorems}
\crefname{conjecture}{Conjecture}{Conjectures}
\crefname{question}{Question}{Questions}
\crefname{example}{Example}{Examples}

\newcommand{\R}{\mathbb R}
\newcommand{\Q}{\mathbb Q}
\newcommand{\Z}{\mathbb Z}
\newcommand{\N}{\mathbb N}
\newcommand{\T}{\mathbb T}
\newcommand{\C}{\mathbb C}
\DeclareMathOperator{\den}{den}
\newcommand{\e}{\mathrm e}

% Immutable Lean source links at research-checkpoint-18.
\newcommand{\leancommit}{e4809cf561573aef588e6ffc491a35cdb79fd9ec}
\newcommand{\leanbase}{https://github.com/eldernos/MathematicsResearch/blob/\leancommit}
% \lean{File}{line}{label}: link to NormalityChecker/File.lean#Lline at the pinned commit.
\newcommand{\lean}[3]{\href{\leanbase/NormalityChecker/#1.lean\#L#2}{#3}}
% \repo{path}{label}: link to any repository path at the pinned commit.
\newcommand{\repo}[2]{\href{https://github.com/eldernos/MathematicsResearch/tree/\leancommit/#1}{#2}}

% Evidence labels (see papers/WAVE2-CONTRACT.md).
\newcommand{\evidence}[1]{\textup{[\emph{#1}]}}

\setlength{\emergencystretch}{1.5em}
\raggedbottom

% immediately after \documentclass[11pt]{article}.
\ifdefined\pdftrailerid\pdftrailerid{}\fi
\usepackage[T1]{fontenc}
\usepackage{lmodern}
\usepackage[letterpaper,margin=1.02in]{geometry}
\usepackage{amsmath,amssymb,amsthm,mathtools,booktabs,array,microtype,enumitem}
\usepackage{tikz}
\usepackage{xurl}
\usepackage[colorlinks=true,allcolors=bluegreen]{hyperref}
\definecolor{bluegreen}{RGB}{22,83,91}
\usepackage[capitalize,nameinlink]{cleveref}

\theoremstyle{plain}
\newtheorem{theorem}{Theorem}[section]
\newtheorem{lemma}[theorem]{Lemma}
\newtheorem{proposition}[theorem]{Proposition}
\newtheorem{corollary}[theorem]{Corollary}
\newtheorem{conjecture}[theorem]{Conjecture}
\newtheorem*{theorem*}{Theorem}
\newtheorem{maintheorem}{Theorem}
\renewcommand{\themaintheorem}{\Alph{maintheorem}}
\theoremstyle{definition}
\newtheorem{definition}[theorem]{Definition}
\newtheorem{example}[theorem]{Example}
\newtheorem{question}[theorem]{Question}
\theoremstyle{remark}
\newtheorem{remark}[theorem]{Remark}
\crefname{maintheorem}{Theorem}{Theorems}
\crefname{conjecture}{Conjecture}{Conjectures}
\crefname{question}{Question}{Questions}
\crefname{example}{Example}{Examples}

\newcommand{\R}{\mathbb R}
\newcommand{\Q}{\mathbb Q}
\newcommand{\Z}{\mathbb Z}
\newcommand{\N}{\mathbb N}
\newcommand{\T}{\mathbb T}
\newcommand{\C}{\mathbb C}
\DeclareMathOperator{\den}{den}
\newcommand{\e}{\mathrm e}

% Immutable Lean source links at research-checkpoint-18.
\newcommand{\leancommit}{e4809cf561573aef588e6ffc491a35cdb79fd9ec}
\newcommand{\leanbase}{https://github.com/eldernos/MathematicsResearch/blob/\leancommit}
% \lean{File}{line}{label}: link to NormalityChecker/File.lean#Lline at the pinned commit.
\newcommand{\lean}[3]{\href{\leanbase/NormalityChecker/#1.lean\#L#2}{#3}}
% \repo{path}{label}: link to any repository path at the pinned commit.
\newcommand{\repo}[2]{\href{https://github.com/eldernos/MathematicsResearch/tree/\leancommit/#1}{#2}}

% Evidence labels (see papers/WAVE2-CONTRACT.md).
\newcommand{\evidence}[1]{\textup{[\emph{#1}]}}

\setlength{\emergencystretch}{1.5em}
\raggedbottom

% immediately after \documentclass[11pt]{article}.
\ifdefined\pdftrailerid\pdftrailerid{}\fi
\usepackage[T1]{fontenc}
\usepackage{lmodern}
\usepackage[letterpaper,margin=1.02in]{geometry}
\usepackage{amsmath,amssymb,amsthm,mathtools,booktabs,array,microtype,enumitem}
\usepackage{tikz}
\usepackage{xurl}
\usepackage[colorlinks=true,allcolors=bluegreen]{hyperref}
\definecolor{bluegreen}{RGB}{22,83,91}
\usepackage[capitalize,nameinlink]{cleveref}

\theoremstyle{plain}
\newtheorem{theorem}{Theorem}[section]
\newtheorem{lemma}[theorem]{Lemma}
\newtheorem{proposition}[theorem]{Proposition}
\newtheorem{corollary}[theorem]{Corollary}
\newtheorem{conjecture}[theorem]{Conjecture}
\newtheorem*{theorem*}{Theorem}
\newtheorem{maintheorem}{Theorem}
\renewcommand{\themaintheorem}{\Alph{maintheorem}}
\theoremstyle{definition}
\newtheorem{definition}[theorem]{Definition}
\newtheorem{example}[theorem]{Example}
\newtheorem{question}[theorem]{Question}
\theoremstyle{remark}
\newtheorem{remark}[theorem]{Remark}
\crefname{maintheorem}{Theorem}{Theorems}
\crefname{conjecture}{Conjecture}{Conjectures}
\crefname{question}{Question}{Questions}
\crefname{example}{Example}{Examples}

\newcommand{\R}{\mathbb R}
\newcommand{\Q}{\mathbb Q}
\newcommand{\Z}{\mathbb Z}
\newcommand{\N}{\mathbb N}
\newcommand{\T}{\mathbb T}
\newcommand{\C}{\mathbb C}
\DeclareMathOperator{\den}{den}
\newcommand{\e}{\mathrm e}

% Immutable Lean source links at research-checkpoint-18.
\newcommand{\leancommit}{e4809cf561573aef588e6ffc491a35cdb79fd9ec}
\newcommand{\leanbase}{https://github.com/eldernos/MathematicsResearch/blob/\leancommit}
% \lean{File}{line}{label}: link to NormalityChecker/File.lean#Lline at the pinned commit.
\newcommand{\lean}[3]{\href{\leanbase/NormalityChecker/#1.lean\#L#2}{#3}}
% \repo{path}{label}: link to any repository path at the pinned commit.
\newcommand{\repo}[2]{\href{https://github.com/eldernos/MathematicsResearch/tree/\leancommit/#1}{#2}}

% Evidence labels (see papers/WAVE2-CONTRACT.md).
\newcommand{\evidence}[1]{\textup{[\emph{#1}]}}

\setlength{\emergencystretch}{1.5em}
\raggedbottom

% Author information approved by the repository owner on 2026-11-25:
%   name: Christopher Anderson; designation: Independent Researcher.
% No academic affiliation, ORCID or personal contact details have been supplied.
% Do not add any that the owner has not approved.
\newcommand{\OwnerAuthor}{Christopher Anderson}
\newcommand{\OwnerAffiliation}{Independent Researcher}
\newcommand{\OwnerContact}{Corrections and comments: \url{https://github.com/eldernos/MathematicsResearch/issues}}
\newcommand{\ReleaseLabel}{manuscripts-v2.0.0}
\newcommand{\LicenseLine}{\copyright\ 2026 Christopher Anderson. Licensed under \href{https://creativecommons.org/licenses/by/4.0/}{CC BY 4.0}.}
% \makebyline{version}{date}: centred author block used under the title.
\newcommand{\makebyline}[2]{%
  {\large \OwnerAuthor\par}%
  \ifx\OwnerAffiliation\empty\else{\small \OwnerAffiliation\par}\fi
  \vspace{0.3em}
  {\small Version #1 \enspace$\cdot$\enspace #2 \enspace$\cdot$\enspace \ReleaseLabel\par}%
  {\footnotesize \OwnerContact\par}%
  {\footnotesize \LicenseLine\par}%
  \hypersetup{pdfauthor={\OwnerAuthor}}}

\hypersetup{pdftitle={Synthetic Arithmetic Expansions and Prescribed Discrepancy}}
\newcommand{\dist}[1]{\lVert #1\rVert}
\newcommand{\fr}[1]{\{#1\}}
\newcommand{\Prob}{\mathcal P}
\newcommand{\Minv}{\mathcal M_T}
\DeclareMathOperator{\ord}{ord}
\DeclareMathOperator{\lcm}{lcm}
\newcommand{\ttt}[1]{\texttt{#1}}
\newcommand{\leanP}[3]{\leanat{\phaseXIXcommit}{#1}{#2}{#3}}
% Lean modules ProtectedPole and PowerSumPeeling, pinned at the later commit below.
\newcommand{\phaseXXIcommit}{347ec3fe631a3e7ab2ae25d53508a3a4c9ea9fa8}
\newcommand{\leanR}[3]{\leanat{\phaseXXIcommit}{#1}{#2}{#3}}
\renewcommand{\ReleaseLabel}{candidate draft, not yet released}
\setlength{\emergencystretch}{3em}

\begin{document}
\begin{center}
{\LARGE\bfseries Synthetic Arithmetic Expansions\\[3pt] and Prescribed Discrepancy\par}
\vspace{10pt}
\makebyline{2.0 (preprint draft)}{27 November 2026}
\end{center}

\begin{abstract}
Fix an integer base $b\ge2$, a prime $p\nmid b$, and let $t$ be the order of $b$ modulo $p$.
For an integer $d\ge2$, insert the coefficient $p^{-s v_p(k)}$ at base-$b$ position $tk^d$ for every
$k\ge1$, and call the resulting number $C(b,p,s,d)$. For $s=0$ the extreme discrepancy tends to one,
while for every $s\ge1$ the number is normal to base $b$; this switch holds on every schedule of
multiples of $t$ dominating $k^2$ and has been checked by the Lean~4 kernel. For odd $p$ we prove
that every nonconstant rational polynomial of total degree $r<d$ in numbers $C(b,p,s_a,d)$ with
distinct orders $s_a$ is normal, with discrepancy $O((\log\log N)^{-d/(d-1)})$ for every $d\ge r+1$,
and a negative power of $N$ when the maximal weighted degree is carried by a single linear term. In
particular $\xi=\sum_k3^{-v_3(k)}2^{-2k^3}$ satisfies $D_N(\xi)\le200N^{-1/3}$ and
$D_N(\xi^2)\ll(\log\log N)^{-3/2}$. The rate rests on a small-ball estimate for the limiting $p$-adic
mass of the shifted prefixes and on a short-orbit bound on prime-power moduli that is uniform in the
base. Choosing the rational state after each cubic insertion realizes every limiting discrepancy with
error $O(N^{-1/3})$ uniformly in the limit, and every admissible pair of extreme limits and set of
invariant measures. A positive binary series with every odd prime as a denominator prime is normal
and Liouville with $D_N=O(1/\log N)$. At the boundary $d=r$ the $p$-adic depth
of the prefix denominators is tight in density, but we prove that denominator data alone cannot
determine the digits on the runs that occur, even over packets of $\lfloor\beta\log_bN\rfloor$ shifts
with exact jumps. Using a theorem of Dietmann, Elsholtz, Kalmynin, Konyagin and Maynard on long gaps
between sums of two squares, $C(b,p,s,2)^2$ is irrational whenever $s<(390/449)\,t\log b$, in particular
for $s=1$; we also determine the limiting law of its fractional innovations on an explicit set of
arrivals. Normality of $\xi^3$ and of $C(b,p,s,2)^2$ remains open. All-prime polynomial normality and
irrationality exponents, treated in the previous draft of this version, are now the subject of a
companion manuscript.
\end{abstract}

\section{Introduction}\label{sec:intro}

Throughout, $b\ge2$ is an integer, $\T=\R/\Z$, $T(y)=by\bmod 1$, $\fr{y}$ is the fractional part
and $\dist{y}$ the distance to the nearest integer. For real $x$ and $N\ge1$ the \emph{extreme
discrepancy} of the first $N$ points of the radix orbit is
\[
D_N(x)=D_N(x;b)=\sup_{0\le a<c\le 1}\left|\frac{\#\{0\le n<N:\fr{b^nx}\in[a,c)\}}{N}-(c-a)\right|.
\]
A real number is normal to base $b$ if and only if $D_N(x)\to0$ \cite[Chapter~1, Theorem~8.1,
p.~70 and Chapter~2, Theorem~1.1, p.~89]{KN}. For a prime $p$ we write $v_p$ for the $p$-adic
valuation and $\ord_m(b)$ for the multiplicative order of $b$ modulo $m$.

The numbers studied here are \emph{synthetic expansions}: sums $\sum_k c_kb^{-n_k}$ with
positive rational coefficients placed at a sparse arithmetic schedule $n_k$, where the coefficients
are chosen so that the rational part $b^n\sum_{n_k\le n}c_kb^{-n_k}$ seen after each shift has a
prescribed reduced denominator or a prescribed residue. The first family is defined by a single
valuation weight.

\begin{maintheorem}[Normality switch]\label{thm:A}
Let $p\nmid b$ be prime, $t=\ord_p(b)$, $s\ge0$ an integer, and let $(n_k)_{k\ge1}$ be strictly
increasing positive multiples of $t$ with $n_k\ge k^2$. Put
\[
X_s=\sum_{k\ge1}\frac{1}{p^{s v_p(k)}\,b^{n_k}}.
\]
Then $D_N(X_0)\to1$, and $D_N(X_s)\to0$ for every $s\ge1$, as $N\to\infty$ through all integers.
\end{maintheorem}

\Cref{thm:A} is kernel-verified (\Cref{thm:switch}); so is the exact reduced denominator of every
shifted prefix (\Cref{thm:den}), which is the arithmetic mechanism: a positive order $s$ keeps a
pole of exact order $s\lfloor\log_pK\rfloor$ after the $K$th insertion, while $s=0$ retires all
denominators. We give a written proof of the following rate.

\begin{maintheorem}[Complete-period rate]\label{thm:B}
Let $p\nmid b$ be an odd prime, $t=\ord_p(b)$, and $s\ge1$, $d\ge s+2$ integers. For
$X=\sum_{k\ge1}p^{-sv_p(k)}b^{-tk^d}$ and $N\ge2$,
\[
D_N(X)\ll_{b,p,s,d} N^{-s/d}+N^{-(d-s-1)/d}+N^{-(d-1)/d}\log N .
\]
For the binary number $\xi=\sum_{k\ge1}3^{-v_3(k)}2^{-2k^3}$ one has $D_N(\xi)\le 200N^{-1/3}$
for every $N\ge1$.
\end{maintheorem}

The second half is proved first, in \Cref{sec:example}, by an argument that uses only the fact
that $2$ is a primitive root modulo every power of $3$. The general estimate is \Cref{thm:rate}.

Changing the coefficient after each insertion lets one prescribe the rational state, and hence the
orbit, on each gap. Retiring a prescribed weighted proportion of gaps gives every discrepancy limit
with an explicit error.

\begin{maintheorem}[Prescribed discrepancy]\label{thm:C}
Let $p\nmid b$ be an odd prime and $t=\ord_p(b)$. For every $\delta\in[0,1]$ there are rationals
$c_k\in(0,1]$ whose reduced denominators divide $p^{1+\lfloor\log_pk\rfloor}\le pk$ such that
$x_\delta=\sum_{k\ge1}c_kb^{-tk^3}$ satisfies
\[
|D_N(x_\delta)-\delta|\le C_{b,p}N^{-1/3}\qquad(N\ge1),
\]
with $C_{b,p}$ independent of $\delta$. On the same schedule $tk^3$ and with the same denominator bound,
every pair $0\le\alpha\le\beta\le1$ occurs as $(\liminf D_N,\limsup D_N)$, with
$D_N=\Lambda_N/N+O(N^{-1/3})$ for an explicit weighted count $\Lambda_N$. Every nonempty compact
connected set of $T$-invariant Borel probability measures on $\T$ occurs, again on the schedule $tk^3$,
as the set of weak limit points of the orbit measures of an expansion whose state after the $k$th
insertion is a purely periodic rational of period at most $k$.
\end{maintheorem}

Qualitatively, \Cref{thm:C} adds nothing to classical realization. Every radix orbit is, after
finitely many shifts, the orbit of a sum $\sum_kc_kb^{-n_k}$ with rationals $c_k\in(0,1)$ whose
denominators are powers of $p$, on any prescribed increasing schedule (\Cref{rem:qualitative}). The
realization of invariant-measure limit sets by digit sequences \cite{Colebrook,Sigmund74,DGS}, combined
with weak continuity of $D$ (\Cref{lem:Dcont}), therefore already yields every limit $\delta$, every pair
$(\alpha,\beta)$ and every limit set by such expansions. What \Cref{thm:C} adds, in the sources we
examined, is quantitative and arithmetic: an error $O(N^{-1/3})$ at every prefix, uniform in $\delta$;
coefficient denominators bounded by $pk$, with the target state chosen from two explicit values; and an
explicit arithmetic schedule of targets, given by the Beatty-type mask
$\lfloor k\delta\rfloor-\lfloor(k-1)\delta\rfloor$. The proofs are in \Cref{sec:program,sec:limitsets}.

The linear switch of \Cref{thm:A} rests on a single surviving pole. For squares that pole can
vanish: in the shifted prefixes of $\xi^2$ the top power of $3$ cancels on a set of shifts of positive
lower density, so the method of \Cref{thm:A} alone does not decide normality. A finer analytic argument
does, in every degree below the degree of the schedule. Write
\[
C(b,p,s,d)=\sum_{i\ge1}p^{-sv_p(i)}b^{-ti^d},
\]
so that $X_s=C(b,p,s,d)$ in \Cref{thm:B} and $\xi=C(2,3,1,3)$.

\begin{maintheorem}[Polynomial images below the schedule degree]\label{thm:D}
Let $p\nmid b$ be an odd prime, $t=\ord_p(b)$, $d\ge2$ an integer, and for integers $a\ge1$ put
$X_a=C(b,p,a,d)$. For every finite set $S$ of positive integers and every nonconstant polynomial
$P\in\Q[z_a:a\in S]$ of total degree $r\le d-1$, the number $P((X_a)_{a\in S})$ is normal to base $b$.
Moreover
\[
D_N(P((X_a)_{a\in S}))\ll_P(\log\log N)^{-d/(d-1)},
\]
and $D_N(P((X_a)_{a\in S}))\ll_PN^{-\eta_P}$ with some $\eta_P>0$ if the largest weighted degree
$\max\sum_a a\nu_a$ over the monomials $\prod_az_a^{\nu_a}$ of $P$ is attained by a single monomial of degree
one; in particular every nonconstant linear form, and every polynomial of degree one in a single $X_s$,
has a power saving. Consequently no $X_s$ is algebraic of degree less than $d$; the distinct
nonconstant monomials of total degree less than $d$ in the $X_a$ are jointly normal; $\xi^2$ is normal
to base $2$ with $D_N(\xi^2)\ll(\log\log N)^{-3/2}$, and the binary shifts of $(\xi,\xi^2)$ are
uniformly distributed in $\T^2$; and $C(2,3,1,4)$, its square and its cube are jointly normal to base $2$.
\end{maintheorem}

The case of total degree at most two and $d\ge3$ is \Cref{thm:quadratic}, whose proof without a rate
is kept in \Cref{sec:quadratic} as the model case. The general qualitative statement combines
\Cref{thm:higher} (one order) and \Cref{thm:several} (several orders); the rates are
\Cref{thm:effective} (one order) and \Cref{thm:severalrate} (several orders). The arithmetic input is the exact all-prefix criterion of \Cref{prop:criterion},
which reduces normality to showing that small reduced denominators occur on a set of shifts of density
zero. After normalization the relevant quantity converges $p$-adically to a function $M_r(u)$ of the
real position $u$ of the shift in its $p$-adic shell, with values in $\Z_p$, whose reduction modulo
each power of $p$ is a finite step function. We show that every level set of $M_r$ is Lebesgue-null.
In degree two, primitive jumps on the grids $p^{-H}\Z$ are isolated with positive proportion by the
third-derivative test of van der Corput, and a second-moment argument places almost every point in
infinitely many isolated neighbourhoods. In degree $r\ge3$ a single family of centres covers too little
measure, and we isolate instead a positive proportion of centres built on dominant anchors, using the
$q$-th derivative test with $q\le r+1$ and a divisor bound for exact coincidences of sums of $d$-th
powers. Uniform isolation at a fixed small scale is false (\Cref{prop:nouniform}); averaging over unions
of intervals gives the small-ball bound $|\{u:M_r(u)\equiv w\bmod p^L\}|\ll(\log L)^{-d/(d-1)}$
(\Cref{thm:smallball}), which is the source of the rate. The remaining analytic input is a short-orbit
estimate on prime-power moduli that is uniform in the base (\Cref{lem:uniformorbit}).

Finally we test synthesis with infinitely many denominator primes.

\begin{maintheorem}[Infinite sums]\label{thm:E}
For an odd prime $\ell$ let $s_3=2$, $s_\ell=1$ for $\ell\ge5$, and
$Y_\ell=\sum_{j\ge1}\ell^{-s_\ell j}2^{-16^{\ell^j}}$. Then $Y=\sum_\ell Y_\ell$, summed over all odd
primes, is normal to base $2$ and a Liouville number, and $D_N(Y)=O(1/\log N)$. Every finite
subfamily of the $Y_\ell$ is jointly normal, and there are positive rationals $c_\ell$ with
$\sum_\ell c_\ell Y_\ell=1$.
\end{maintheorem}

Thus finite joint normality of positive terms does not survive infinite summation, while a
uniform arithmetic grid (\Cref{lem:grid}) at the prime $3$ survives the introduction of every other
odd prime (\Cref{sec:sums}).

The series $Y$ retains every odd prime but is only shown to be normal itself. Protecting one pole
while auxiliary primes accumulate gives polynomial normality as well; those constructions, together
with their irrationality exponents, form the companion manuscript \cite{PaperXII} (see the paragraph on
companion manuscripts below).

At the boundary $d=r$ the insertions are no longer sparse, and the method stops. We locate where.

\begin{maintheorem}[The boundary $d=r$]\label{thm:boundary}
Let $p\nmid b$ be an odd prime, $s\ge1$, $d\ge2$, $C=C(b,p,s,d)$, and let $q_n$ be the reduced
denominator of the shifted prefix of $C^d$ at the shift $n=tp^{dJ(n)}u$, $1\le u<p^d$.
\begin{enumerate}[label=\textup{(\roman*)},nosep]
\item For every $L\ge3$ the shifts with $v_p(q_n)<dsJ(n)-L$ have upper density at most
$C_{p,s,d}(\log L)^{-d/(d-1)}$.
\item If $q=p^e$ and $b^T\le q/2$, every word of $T$ base-$b$ digits is the leading word of $a/q$ for
some $a$ prime to $p$; and for $b^T\le\epsilon q$ no bound for
$\sup_{p\nmid a}|\sum_{j<T}e(hab^j/q)|$ improves on $T$ by more than $2\pi|h|\epsilon/(b-1)$.
\item For $d=2$ and every $H\ge0$ one has $p^{sH}C^2=X_H+Z_H$, where $X_H$ has integer coefficients and
$1-D_N(X_H)\ll(\log\log N+H)/\sqrt{\log N}$ uniformly in $H$; $C^2$ is normal if and only if $Z_H$ is;
and for fixed $H$ and every $\epsilon>0$ the shifts $n<N$ lying in runs of length at least
$\epsilon\log N$ between consecutive insertions of $Z_H$ form a proportion tending to $0$, so that
almost every shift lies in a run below the threshold of \textup{(ii)}.
\item For $d=3$ the same decomposition holds, but the sparse-series transfer used for $d=2$ shows
$D_N(X_H)\to1$ only if the sums of three positive cubes have density zero, which is open.
\item For $d=2$ and $s<(390/449)\,t\log b$, in particular for $s=1$, the number $C^2$ is irrational.
\item For $d=2$, every $H\ge0$ and $k=\lceil(H+1)/2\rceil$, the fractional innovations of $Z_H$ at the
positions $p^{2k}n$, with $n>1$ a squarefree sum of two squares and $(n/p)=-1$, are asymptotically
uniformly distributed on the $\ord_D(2)$ residues $2^j/D\bmod1$, $D=p^{s(2k-H)}$, for both residue classes
of $p$ modulo $4$.
\item For $0<\beta<s$, at all but $o(N)$ starts $n\in[N/2,N]$ the exact jumps and all reduced
prefix denominators through the packet $(n,n+\lfloor\beta\log_bN\rfloor]$ admit every base-$b$ word of
that length when the incoming numerator varies.
\end{enumerate}
\end{maintheorem}

This is \Cref{sec:boundary}. Item (i) says that the $p$-adic half of the method survives at $d=r$,
and items (ii)--(iii) and (vii) that the real half does not: any proof must use the numerators of the
shifted prefixes at resolution $b^{-T}$, which is the distribution of the orbit itself. Item (v) uses
a theorem of Dietmann, Elsholtz, Kalmynin, Konyagin and Maynard \cite[Theorem~1]{DEKKM} as an external
input. Normality of $\xi^3$ and of $C(b,p,s,2)^2$, and even $\limsup_ND_N<1$, remain open, and so does
the irrationality of $\xi^3$ (\Cref{sec:questions}).

\paragraph{Companion manuscripts.} The perturbation estimates used here are finite comparisons
(\Cref{lem:compare}). Their sharp form, the identity $\sup_x|D_N(x+z)-D_N(x)|=1-D_N(z)$, together with
the classification of translations preserving discrepancy, the exact addition spectrum of $D_N$ and
the behaviour of $D_N$ under rational scalars, is the subject of the companion manuscript
\cite{CompanionH}. We quote the identity once (\Cref{thm:sharpcited}) and use it in \Cref{prop:tails}
and \Cref{cor:retire}. A second companion manuscript, \cite{PaperXII}, a candidate draft like this one,
treats series whose shifted prefix denominators eventually contain every prime not dividing $b$ and all of
whose nonconstant rational polynomial images are normal: a Liouville example in all degrees, polynomial
schedules with growing auxiliary support, prescribed finite irrationality exponents, and exponent two in
all degrees. It uses \Cref{lem:uniformorbit}, \Cref{cor:split}, \Cref{thm:smallball} and the weighted
isolation estimate of \Cref{sec:rates} as inputs, restating the first two.

\paragraph{Relation to earlier work.} Arithmetic constructions of normal numbers by sparse series
go back to Stoneham \cite{Stoneham73} and Korobov \cite{Korobov72}; Kano \cite[Theorem~1]{Kano93} gives
a general criterion for series $\sum A_j/(a^{\lambda_j}b^{\mu_j})$, with strictly increasing exponents
and size conditions on the numerators, and records in its introduction Stoneham's constructions with
prescribed numerators. Bailey and Crandall \cite[Theorem~4.8]{BC02} prove normality for such series
with numerator one under growth conditions on the two exponent sequences, and Vandehey
\cite[Corollary~1.2 and Theorem~7.4]{Vandehey19} gives power-saving estimates for incomplete Korobov
sums on moduli with fixed prime support, together with a normality criterion for series
$\sum 1/(c_kb^{m_k})$ with numerator one in which $c_k\mid c_{k+1}$ strictly. Two denominator sequences
should be distinguished in \Cref{thm:A}. The insertion denominators $p^{sv_p(k)}$ return to $1$ whenever
$p\nmid k$, so the series is not term by term an instance of these criteria; the reduced denominators
$p^{s\lfloor\log_pK\rfloor}$ of the shifted prefixes do form a divisibility chain. Grouping consecutive
insertions into single terms does not help either: on the cubic schedule a group of two or more terms,
written over its last radix position, has size times that radix power tending to infinity, which is
incompatible with the bounded normalized numerators of these criteria (the test is recorded in the
ledger, \Cref{sec:formal}). We therefore do not apply a published series criterion. The positive side
of \Cref{thm:A} combines the exact prefix denominators (\Cref{thm:den}) with a fixed-support short-orbit
estimate (\Cref{lem:short}), the qualitative consequence of Vandehey's Corollary~1.2; the side $s=0$
is elementary, and the switch holds on every schedule dominating $k^2$. Rings of numbers normal to one
base and not to another were constructed by Wagner \cite{Wagner95}, an uncountable ring all of whose
nonzero elements are normal to base $b$ and not normal to base $ab$ for a prime $a\nmid b$ (as described
in \cite[\S1]{Kano93} and in the introduction of \cite{KS93}), and by Kano and Shiokawa
\cite[Theorem~5]{KS93} for $a$ not necessarily prime. Since nonzero polynomial expressions in the
generators lie in the ring, these are direct precedents for the normality of polynomial images of
suitable generators. Polynomial images of the standard Stoneham numbers are treated in \cite{PaperI},
and Mira Cult's public note \cite[Theorem~1.1]{MiraPoly} proves the normality of every nonconstant
rational polynomial in $\sum_kq^{-k}b^{-c^k}$ when $q>2$ and some prime divides $c$ and $q$ but not $b$
(see the comparison after \Cref{cor:flagship}); the same statement for $c>2$ coprime to $b$ is
kernel-verified in \cite[Theorem~A]{PaperIv2}. The ring generators use rapidly separated exponents,
and \cite{MiraPoly} uses the exponential schedule $c^k$. \Cref{thm:D} concerns specific generators on a polynomial
schedule, whose consecutive positions have ratio tending to one and whose weight exponents $sv_p(k)$
return to zero infinitely often; we have not verified that the generator hypotheses of the ring
constructions hold for them. Its conclusion is limited to total degree below $d$, and it does not
produce a ring. The comparison of the all-prime constructions with the work of Manai and with the
literature on irrationality exponents is made in \cite{PaperXII}.
Normal Liouville numbers are known explicitly \cite{KS93,BC02}, absolutely normal Liouville numbers
exist \cite{Bugeaud02}, and a computable one was constructed in \cite{BHS15}; \Cref{cor:liouville} is a
further explicit example inside the switch family. For a family of signed Stoneham expansions the
attainable upper discrepancies form an explicit proper subset of $[0,1]$ with a gap above zero
\cite{PaperIII}, whereas \Cref{thm:C} shows that positive coefficients on a cubic schedule reach every
value with a uniform rate. Recent public work proves the normality of sums of two Stoneham numbers \cite{SudeSum} and the
joint normality of Stoneham series on finitely many distinct primes and orders \cite{MiraJoint};
\Cref{thm:combos} is a result of the same kind for the different series of \Cref{thm:E}. Normal numbers
with very small discrepancy are constructed in \cite{Levin99,ABSS17}; the rates of \Cref{thm:B,thm:C}
are much weaker and are included because they are elementary and hold for the specific arithmetic
families; the same applies to the rates of \Cref{thm:D}. The analytic inputs of \Cref{sec:quadratic}
are van der Corput's method in the forms of Graham--Kolesnik \cite{GK91} and Titchmarsh's $q$-th
derivative test \cite[Theorem~5.13]{Titchmarsh86}, and the Erd\H{o}s--Tur\'an inequality \cite{KN}. The
short-orbit bound of \Cref{lem:uniformorbit} is an elementary stride-differencing estimate in the
spirit of Korobov \cite{Korobov72} and Vandehey \cite{Vandehey19}; its point is uniformity in the base.
In \Cref{sec:boundary} we use Landau's theorem on sums of two squares \cite{Landau08}, Hooley's moment
bound for their gaps \cite{Hooley71}, the Bambah--Chowla gap bound \cite{BambahChowla47} and the
Selberg--Delange method \cite{Tenenbaum15}, for the irrationality of squares the long gaps between sums
of two squares of Dietmann, Elsholtz, Kalmynin, Konyagin and Maynard \cite[Theorem~1]{DEKKM}, and for the
cubic boundary Wooley's lower bound for sums of three cubes \cite{Wooley15}, the probabilistic model of
Deshouillers, Hennecart and Landreau \cite{DHL06} and Ghidelli's gaps between values of diagonal cubic
forms \cite[Theorem~1.1]{Ghidelli20}. We have not located \Cref{thm:A,thm:D,thm:E} in the sources examined; a detailed
comparison, including recent public repositories, is kept in the ledger accompanying the source files
(\Cref{sec:formal}).

\paragraph{Evidence labels.} Each numbered statement carries one label:
\emph{Kernel-verified} (the complete statement is a Lean declaration checked by the kernel, linked
at a pinned commit), \emph{Lean-supported paper proof} (named steps are Lean declarations),
\emph{Paper proof}, or \emph{External input}. A statement whose proof includes an exact computation
over a stated finite range carries in addition the label \emph{finite certificate}, and a proved
implication whose hypothesis is open carries in addition the label \emph{conditional}; questions that
remain unresolved are called \emph{open}. Kernel-verified statements also receive written proofs or proof outlines here.
A Lean-supported label means only that the named finite step is formal; the analytic estimates,
the measure theory and every rate are written proofs. Finite computations quoted as illustrations
are measurements, not evidence for asymptotic statements.

\section{Discrepancy on the circle and orbits of rationals}\label{sec:prelim}

For a finite sample $u=(u_0,\ldots,u_{N-1})$ in $\T$, identified with $[0,1)$, let $\mu_u$ be its
empirical probability measure and $D(u)=\sup_{0\le a<c\le1}|\mu_u([a,c))-(c-a)|$. A
\emph{half-open arc} is the image in $\T$ of $[\alpha,\alpha+l)$ with $0\le l\le1$; its length is
$|I|=l$, and its \emph{excess} is $\mu_u(I)-|I|$.

\begin{lemma}[Arcs]\label{lem:arcs}\evidence{Paper proof}
For every half-open arc $I$, $|\mu_u(I)-|I||\le D(u)$, and $D(u)=\sup_I(\mu_u(I)-|I|)$, the supremum
of the excess over half-open arcs. Consequently $D(u)$ is unchanged if every $u_n$ is rotated by
the same $\theta\in\T$.
\end{lemma}

\begin{proof}
A non-wrapping arc is an interval $[a,c)$. A wrapping arc is the complement of an interval
$[c,a)$, and its excess is the negative of the excess of that interval; the full circle has
excess zero. This proves the first claim and that the supremum of excesses is at most $D(u)$.
Conversely an interval with negative excess $-e$ has a complementary arc with excess $e$.
Rotation permutes half-open arcs and preserves length.
\end{proof}

\begin{lemma}[Finite comparison]\label{lem:compare}\evidence{Paper proof}
Let $u,v$ be samples of length $N$, let $\epsilon\ge0$ and $E=\#\{n<N:\dist{u_n-v_n}>\epsilon\}$.
Then $|D(u)-D(v)|\le E/N+2\epsilon$.
\end{lemma}

\begin{proof}
Let $I=[\alpha,\alpha+l)$. If $l>1-2\epsilon$, then $\mu_u(I)-l<2\epsilon$. Otherwise each $n$
with $u_n\in I$ and $\dist{u_n-v_n}\le\epsilon$ has $v_n$ in the half-open arc
$I'=[\alpha-\epsilon,\alpha+l+\epsilon)$ of length $l+2\epsilon$. Hence
$\mu_u(I)\le\mu_v(I')+E/N\le l+2\epsilon+D(v)+E/N$ by \Cref{lem:arcs}. Taking the supremum
over $I$ gives $D(u)\le D(v)+E/N+2\epsilon$; exchange $u$ and $v$.
\end{proof}

When $v_n=\fr{b^n(x+z)}$ and $u_n=\fr{b^nx}$, the exceptional set depends only on $z$; this case
is the kernel-verified declaration
\lean{Phase18Synthetic}{583}{\ttt{discrepancy\_transfer}}. We shall often use the following
form, in which part of the orbit is left uncovered.

\begin{lemma}[Block certificate]\label{lem:blocks}\evidence{Paper proof}
Let $x\in\R$, $N\ge1$, and let $B_1,\ldots,B_m\subset[0,N)$ be disjoint sets of integers with
$U=N-\sum_i|B_i|$. Suppose that $y_n\in\T$ is given for $n\in\bigcup B_i$ and that for every $i$
and every half-open arc $I$,
$|\#\{n\in B_i:y_n\in I\}-|B_i||I||\le H_i$. If $\dist{b^nx-y_n}\le\epsilon$ for all covered $n$
except $E$ of them, then
\[
D_N(x)\le\frac{U+E+\sum_iH_i}{N}+2\epsilon .
\]
\end{lemma}

\begin{proof}
Put $y_n=\fr{b^nx}$ at uncovered $n$. For every half-open arc the count of $(y_n)_{n<N}$ differs
from $N|I|$ by at most $\sum_iH_i+U$, so $D(y)\le(U+\sum_iH_i)/N$ by \Cref{lem:arcs}. Apply
\Cref{lem:compare}.
\end{proof}

The rational models will be orbits of rationals with denominator prime to $b$. We use the
standard lifting of exponents: if $\ell$ is an odd prime and $\ell\mid y-1$, then
$v_\ell(y^m-1)=v_\ell(y-1)+v_\ell(m)$ for all $m\ge1$. (For $\ell\nmid m$ the quotient
$(y^m-1)/(y-1)\equiv m\pmod\ell$; and if $y=1+\ell^ju$ with $j\ge1$, then
$y^\ell=1+\ell^{j+1}u+\binom{\ell}{2}\ell^{2j}u^2+\cdots\equiv1+\ell^{j+1}u\pmod{\ell^{j+2}}$.)

\begin{lemma}[Grids]\label{lem:grid}\evidence{Lean-supported paper proof}
Let $p$ be an odd prime, $q=p^eQ$ with $\gcd(Q,pb)=1$ and $p\nmid b$. Let $\tau\ge1$ satisfy
$b^\tau\equiv1\pmod{pQ}$, and put $c=v_p(b^\tau-1)$. If $e\ge c$, then for every integer $a$ prime to
$p$ (not necessarily prime to $Q$) and every $j_0\ge0$, the $S=\tau p^{e-c}$ points $\fr{ab^j/q}$, $j_0\le j<j_0+S$, form $\tau$
arithmetic progressions in $\T$, each of $p^{e-c}$ points with common difference $p^{c-e}$.
Consequently, for every $L\ge1$, every $L$ consecutive points $\fr{ab^j/q}$, and every half-open arc
$I$,
\[
\bigl|\#\{j:\fr{ab^j/q}\in I\}-L|I|\bigr|\le Lp^{c-e}+\tau p^{e-c}.
\]
\end{lemma}

\begin{proof}
Put $y=b^\tau$. By lifting the exponent, $y^m\equiv1\pmod{p^e}$ if and only if
$p^{e-c}\mid m$, so $y$ generates a subgroup of order $p^{e-c}$ of the units modulo $p^e$, contained
in the set of residues congruent to $1$ modulo $p^c$; that set has $p^{e-c}$ elements, so the two
coincide. Since $a$ is a unit modulo $p^e$, the residues $ay^m$ modulo $p^e$ are exactly those
congruent to $a$ modulo $p^c$; since $y\equiv1\pmod Q$, $ay^m\equiv a\pmod Q$ for every integer $a$.
The Chinese remainder theorem therefore gives
$\{ay^m\bmod q\}=\{r\bmod q:r\equiv a\pmod{p^cQ}\}$, and these residues divided by $q$ form a
progression of $p^{e-c}$ points with difference $p^cQ/q$. Writing $j=j_0+r+\tau m$ with $0\le r<\tau$
splits $S$ consecutive points into $\tau$ such cosets. A progression of $M$ points with difference
$1/M$ meets every half-open arc $I$ in $M|I|+\theta$ points with $|\theta|\le1$. For general $L$
write $L=fS+g$ with $0\le g<S$; the error is at most $f\tau+g\le L\tau/S+S$. The order computation
and the count over one stretch of length $S$ starting at $j_0=0$ (equivalently, at any numerator)
are the Lean declarations \leanP{SurvivingGrid}{43}{\ttt{stride\_order}} and
\leanP{PrimePowerGrid}{337}{\ttt{crt\_period\_count}}, the latter also for numerators prime to $p$
only; the passage to arbitrary $L$ is written.
\end{proof}

The lemma counts along one chosen admissible stride $\tau$, which is all that \Cref{thm:Y} needs. For
an arbitrary admissible $\tau$ the $\tau$ progressions need not be distinct, and $p^{e-c}$ need not be
the largest size of an equally spaced subset of the orbit; we do not use the exact period or the
maximal grid.

\begin{corollary}[Prime-power orbits]\label{cor:period}\evidence{Paper proof}
Let $p\nmid b$ be an odd prime, $t=\ord_p(b)$ and $v=v_p(b^t-1)$. For $e\ge0$, $a$ prime to $p$,
$L\ge1$ and every half-open arc $I$, any $L$ consecutive points $\fr{ab^j/p^e}$ satisfy
\[
\bigl|\#\{j:\fr{ab^j/p^e}\in I\}-L|I|\bigr|\le p^vL/p^e+p^e .
\]
\end{corollary}

\begin{proof}
If $e\ge v$ apply \Cref{lem:grid} with $Q=1$, $\tau=t$, $c=v$, noting $tp^{e-v}<p^e$. If $e<v$ the
left side is at most $L\le p^vL/p^e$.
\end{proof}

The sharp form of \Cref{lem:compare} for translations is proved in the companion manuscript. We quote
it here and use it only in \Cref{prop:tails}.

\begin{theorem}[Sharp translation identity {\cite[Theorem~3.1]{CompanionH}}]\label{thm:sharpcited}\evidence{Kernel-verified}
For all real $x,z$ and every $N\ge1$, $|D_N(x+z)-D_N(x)|\le1-D_N(z)$, with equality at $x=0$.
\end{theorem}

This is Theorem~3.1 of \emph{The Arithmetic of Extreme Discrepancy: Translation, Addition Spectra and
Rational Operations} \cite{CompanionH}, in the same release, where it is proved; it is the Lean
declaration \lean{SharpTranslation}{343}{\ttt{sharp\_translation}}, with the equality case
\lean{SharpTranslation}{386}{\ttt{sharp\_translation\_isGreatest}}. We do not reprove it. Taking $x=-z$
gives $D_N(-z)=D_N(z)$.

\subsection{Short orbits uniform in the base}\label{sec:uniformorbit}

The fixed-support estimate \Cref{lem:short} below is qualitative, and its support is fixed. The rates
of \Cref{thm:D} need an explicit estimate for orbits of a principal unit on a prime-power
modulus, uniform in the unit. For an integer $G\equiv1\pmod p$ with $G\ne1$ put $A(G)=v_p(G-1)$, and
write $e(\theta)=\exp(2\pi i\theta)$.

\begin{lemma}[Uniform short orbits]\label{lem:uniformorbit}\evidence{Paper proof}
Let $p$ be an odd prime and $\theta\in(0,1]$. Put $n_\theta=\lceil2/\theta\rceil+1$,
$\delta=2^{-2-n_\theta}$ and $M_*=p^{n_\theta+8}$. Let $M_0\ge M_*$, let $G\equiv1\pmod p$ be an integer
with $G\ne1$ and $A=A(G)\le\frac3{16}\log_pM_0-1$, and let $e,c$ be integers with
$A<e\le\theta^{-1}\log_pM_0$ and $p\nmid c$. Then for every integer $0\le M\le M_0$,
\[
\Bigl|S(G,e,c;M)\Bigr|\le9M_0^{1-\delta},\qquad S(G,e,c;M)=\sum_{0\le m<M}e\Bigl(\frac{cG^m}{p^e}\Bigr).
\]
\end{lemma}

\begin{proof}
Three facts are used. (i) By lifting the exponent (\Cref{sec:prelim}), $v_p(G^n-1)=A+v_p(n)$ for $n\ge1$.
(ii) Hence $G$ has order $p^{e-A}$ modulo $p^e$ and generates the subgroup $1+p^A\Z/p^e\Z$, which has
the same order. For $X=e-A\ge1$ a complete period therefore gives
$\sum_{z\bmod p^X}e(c(1+p^Az)/p^e)=e(c/p^e)\sum_{z\bmod p^X}e(cz/p^X)=0$, so $|S|<p^X$ for every length
$M$. (iii) For integers $H,w\ge1$, extending the summands by zero and applying the Cauchy--Schwarz
inequality to $H\cdot S=\sum_n\sum_{i<H}a_{n+iw}$ gives
\[
|S|^2\le\frac{M+(H-1)w}{H}\Bigl(M+2\sum_{h=1}^{H-1}|C_h|\Bigr),\qquad
C_h=\sum_{0\le m<\max(0,M-hw)}e\Bigl(\frac{c(G^{m+hw}-G^m)}{p^e}\Bigr).
\]
By (i), $c(G^{hw}-1)=c_hp^{A+v_p(hw)}$ with $p\nmid c_h$, so $C_h=S(G,e_h,c_h;\max(0,M-hw))$ with the
\emph{same base} $G$ and $e_h=e-A-v_p(hw)$.

Put $L_0=\log_pM_0\ge n_\theta+8\ge11$, $W=\lfloor\frac34L_0\rfloor$, $\kappa=\lfloor\frac12L_0\rfloor$ and
$H=\lfloor M_0^{1/4}\rfloor$. Let $\mathcal F$ be the set of triples $(G,e,c)$ with $A(G)\le W/4$,
$e-A(G)\ge1$ and $p\nmid c$. We prove by induction on $n\ge0$: \emph{if $(G,e,c)\in\mathcal F$ and
$X=e-A\le W+n\kappa$, then $|S(G,e,c;M)|\le C_nM_0^{1-\delta_n}$ for $0\le M\le M_0$}, where $C_0=1$,
$C_{n+1}=(4+8C_n)^{1/2}$ and $\delta_n=2^{-2-n}$. Since $4+2\sqrt5<9$ is the fixed point of
$x\mapsto(4+8x)^{1/2}$, every $C_n<9$.

If $X\le W$, (ii) gives $|S|<p^W\le M_0^{3/4}$. Let $W+n\kappa<X\le W+(n+1)\kappa$. Choose $k\in[0,\kappa]$
as follows: $k=0$ if $X-A\le W$; $k=X-A-W$ if $W<X-A\le W+\kappa$; $k=\kappa$ if $X-A>W+\kappa$. Then
$X_1=X-A-k$ satisfies $\frac34W\le X_1\le W+n\kappa$ (the first case uses $A\le W/4$, the last uses
$X\le W+(n+1)\kappa$ and $A\ge1$). Take $w=p^k\le M_0^{1/2}$. For $1\le h<H$ put $j=v_p(h)\le\frac14L_0$;
the differenced triple has $X_h=e_h-A=X_1-j$ with
$1\le\frac5{16}L_0-\frac34\le\frac34W-\frac14L_0\le X_h\le W+n\kappa$, so every $C_h$ falls under the case
$n$. Since $M+(H-1)w\le2M_0$ and $H\ge\frac12M_0^{1/4}$,
\[
|S|^2\le4M_0^{3/4}\bigl(M_0+2HC_nM_0^{1-\delta_n}\bigr)\le(4+8C_n)M_0^{2-\min(1/4,\delta_n)},
\]
which is the case $n+1$ with $\delta_{n+1}=\delta_n/2$. Finally the hypotheses give $A\le W/4$, so
$(G,e,c)\in\mathcal F$, and $e-A<e\le L_0/\theta\le W+n_\theta\kappa$, because
$W+n_\theta\kappa\ge\frac34L_0-1+(\frac2\theta+1)(\frac{L_0}2-1)\ge\frac{L_0}\theta+\frac54L_0-n_\theta-1$
and $L_0\ge n_\theta+8$. Apply the case $n=n_\theta$.
\end{proof}

Some bound on $A(G)$ of order $\log_pM_0$ is necessary. Let $G=1+p^a$ with $p^a\ge10$, $e=2a$ and
$c=1$. Then $G^m\equiv1+mp^a\pmod{p^{2a}}$, so $S=e(p^{-2a})\sum_{m<M}e(m/p^a)$. For
$M=\lfloor p^a/10\rfloor\ge1$ we have $0<M/p^a\le\frac1{10}$, and the geometric sum gives the exact value
\[
\frac{|S|}M=\frac{\sin(\pi M/p^a)}{M\sin(\pi/p^a)}\ge\frac{\sin(\pi M/p^a)}{\pi M/p^a}\ge
\frac{\sin(\pi/10)}{\pi/10}>0.983,
\]
using $\sin x\le x$ and the fact that $\sin x/x$ decreases on $(0,\pi)$. Here $M\ge p^{e/2-3}$ and
$A=a=\log_p(10M)+O(1/M)$ for fixed $p$. This exhibits a bad family at the logarithmic scale; it does not
show that every choice of exponent or length at that scale fails.

\begin{corollary}[Splitting off a coprime modulus]\label{cor:split}\evidence{Paper proof}
Let $p\nmid b$ be an odd prime, $t=\ord_p(b)$, $a_0=v_p(b^t-1)$, and let $y=a/(p^eQ)$ with
$\gcd(a,pQ)=1$ and $\gcd(Q,pb)=1$. Put $T_Q=\ord_Q(b)$, $\tau=\lcm(t,T_Q)$ and $G=b^\tau$. Then
$A(G)=a_0+v_p(T_Q)$, and for $T\ge1$ and $M_0=\lceil T/\tau\rceil$,
\[
\Bigl|\sum_{j<T}e(b^jy)\Bigr|\le\tau\max_{p\nmid c,\ M\le M_0}|S(G,e,c;M)|.
\]
If \Cref{lem:uniformorbit} applies to $(M_0,G,e)$ and $\tau\le T$, the sum is at most $18T(\tau/T)^\delta$.
\end{corollary}

\begin{proof}
By lifting the exponent, $v_p(b^{tm}-1)=a_0+v_p(m)$, and $\tau/t$ has the $p$-adic valuation of $T_Q$
because $p\nmid t$. By the Chinese remainder theorem $y\equiv a_1/p^e+a_2/Q\pmod1$ with $p\nmid a_1$.
For $j=\tau m+j_0$ with $0\le j_0<\tau$ we have $b^{\tau m}\equiv1\pmod Q$, so
$e(b^jy)=e(b^{j_0}a_2/Q)\,e(b^{j_0}a_1G^m/p^e)$ with $p\nmid b^{j_0}a_1$. Summing over the at most $M_0$
admissible $m$ gives the first bound, and the second is $9\tau M_0^{1-\delta}\le18T^{1-\delta}\tau^\delta$
(the ceiling in $M_0$ costs at most the factor $2$, and $Q=1$ is allowed, with order $1$).
\end{proof}

In this paper $Q$ always divides a fixed integer $M'$ prime to $pb$. Then $A(G)\le a_0+v_p(\ord_{M'}b)$, and
the standing bound $v_p(\ord_{M'}b)\le\log_pM'=O(1)$ holds because $\ord_{M'}b\le M'$; the weak inequality
is needed, since equality holds at $M'=1$, where both valuations vanish. Growing auxiliary moduli, for which
$A(G)$ must be controlled through the largest auxiliary prime, are treated in \cite{PaperXII}.

\section{A worked example: the binary cubic number}\label{sec:example}

Let
\[
\xi=\sum_{k\ge1}\frac{1}{3^{v_3(k)}\,2^{2k^3}}
=\frac1{2^{2}}+\frac1{2^{16}}+\frac1{3\cdot2^{54}}+\frac1{2^{128}}+\cdots .
\]
Thus the coefficient $1/3^{v_3(k)}$ is inserted at binary position $2k^3$. The companion number
with all coefficients equal to one is $\xi_0=\sum_k2^{-2k^3}$, whose binary expansion consists of
isolated ones.

\begin{theorem}\label{thm:example}\evidence{Paper proof}
For every $N\ge1$, $D_N(\xi;2)\le200N^{-1/3}$. In particular $\xi$ is normal to base $2$. On the
other hand $D_N(\xi_0;2)\to1$.
\end{theorem}

\begin{proof}
Write $n_k=2k^3$ and, for $n_k\le n<n_{k+1}$, let
\[
R_n=2^n\sum_{j\le k}\frac{1}{3^{v_3(j)}2^{n_j}}
\]
be the rational number seen after $n$ shifts. Let $J=\lfloor\log_3k\rfloor$.

\emph{Step 1: the reduced denominator of $R_n$ is exactly $3^J$.} Every $j\le k$ has
$v_3(j)\le J$, so $3^JR_n=\sum_{j\le k}3^{J-v_3(j)}2^{n-n_j}$ is an integer. Modulo $3$ only the
indices with $v_3(j)=J$ survive, that is $j=3^Jm$ with $3\nmid m$ and $m\le k/3^J<3$; there are one
or two of them. Since $2^2\equiv1\pmod3$ and $n_j$ is even, each contributes $2^{n-n_j}\equiv2^n$.
Hence $3^JR_n\equiv(1\text{ or }2)\cdot2^n\not\equiv0\pmod3$.

\emph{Step 2: orbits of the models.} On $[n_k,n_{k+1})$ we have $R_{n+1}=2R_n$, so the points
$y_n=\fr{R_n}$ are consecutive points $\fr{a2^j/3^J}$ with $3\nmid a$. Here $t=\ord_3(2)=2$ and
$v_3(2^2-1)=1$. By \Cref{cor:period} the $\ell$ points of any piece of this block meet every
half-open arc in $\ell|I|+\theta$ points with $|\theta|\le3\ell/3^J+3^J$.

\emph{Step 3: the real orbit.} Positivity and distinctness of the positions give
$0\le2^n\xi-R_n\le\sum_{m\ge n_{k+1}}2^{n-m}=2^{1-(n_{k+1}-n)}$.

\emph{Step 4: counting.} Let $N\ge2$ and choose $K\ge1$ with $2K^3\le N<2(K+1)^3$. Apply
\Cref{lem:blocks} with blocks $B_k=[n_k,\min(n_{k+1},N))$, $1\le k\le K$, and $U=n_1=2$. The
block $B_k$ has length at most $h_k=n_{k+1}-n_k=2(3k^2+3k+1)$, and $k/3<3^{J}\le k$, so Step~2 gives
$H_k\le9h_k/k+k\le127k$ and $\sum_{k\le K}H_k\le127K^2$. Take
$W=\lceil\log_2(K+1)\rceil+2$ and $\epsilon=1/(2(K+1))$. By Step~3 the positions with
$n_{k+1}-n\ge W$ have $\dist{2^n\xi-y_n}\le2^{1-W}\le\epsilon$; at most $W$ positions in each
block fail this, so $E\le KW$. \Cref{lem:blocks} gives
\[
D_N(\xi)\le\frac{2+KW+127K^2}{N}+\frac1{K+1}.
\]
Using $N\ge2K^3$ and $W\le\log_2(K+1)+3\le4K$, the right side is at most
$(1+2+63.5+1)/K=67.5/K$. Since $N<2(K+1)^3\le16K^3$, $1/K<16^{1/3}N^{-1/3}<2.52N^{-1/3}$, and the
bound follows. For $N=1$ it is trivial.

For $\xi_0$ every $R_n$ is an integer and $0<\fr{2^n\xi_0}\le2^{1-(n_{k+1}-n)}$. Fix $W\ge1$. At
most $W(K+1)=O(N^{1/3})$ indices $n<N$ have $n_{k+1}-n<W$, so the interval $[2^{1-W},1)$ contains
$o(N)$ of the first $N$ points while having length $1-2^{1-W}$. Hence
$\liminf D_N(\xi_0)\ge1-2^{1-W}$ for every $W$.
\end{proof}

The constant $200$ is not optimized. For $N=2^{10},2^{11},\ldots,2^{18}$, enclosures of $D_N(\xi)$
to within $2^{-62}$ (computed exactly from the binary expansion and transferred by \Cref{lem:compare};
the script is listed in \Cref{sec:formal}) give values of $N^{1/3}D_N(\xi)$ between $3.477$, at
$N=2^{12}$, and $7.024$, at $N=2^{15}$. The proof uses only the denominator identity of Step~1 and complete periods. \Cref{sec:switch} shows
that Step~1 is an instance of a general identity, and that normality persists when the
complete-period count no longer gives a rate.

\section{Valuation-weighted expansions}\label{sec:switch}

Fix a prime $p\nmid b$, $t=\ord_p(b)$, an integer $s\ge0$, and strictly increasing positive integers
$(n_k)_{k\ge1}$ with $t\mid n_k$. Put $n_0=0$, $w(k)=p^{-sv_p(k)}$,
\[
X=X_s=\sum_{k\ge1}w(k)b^{-n_k},\qquad U_K=\sum_{k\le K}w(k)b^{-n_k},
\]
and for $n\ge0$ let $K(n)=\#\{k\ge1:n_k\le n\}$ and $R_n=b^nU_{K(n)}$.

\subsection{Grouped and raw forms}

A sparse expansion $\sum_kc_kb^{-n_k}$ also has a dense representation, in which the coefficient
$c_k$ is spread over the whole gap $(n_{k-1},n_k]$. When the multipliers have numerator one and the
gaps are divisible by every positive integer, the reduced denominators of the dense coefficients
involve infinitely many primes, which disappear in the sparse form.

\begin{proposition}[Geometric grouping]\label{prop:group}\evidence{Kernel-verified}
For integers $m\ge0$, $h\ge1$ and every field element $b$ with $b\ne0$ and $b^h\ne1$,
\[
\frac{b-1}{b^h-1}\sum_{r=1}^{h}b^{-(m+r)}=b^{-(m+h)} .
\]
Consequently, for real $b>1$, every strictly increasing $(n_k)$ with $n_0=0$ and bounded real
multipliers $c_k$, the series whose coefficient at $b^{-m}$, $n_{k-1}<m\le n_k$, equals
$c_k(b-1)/(b^{n_k-n_{k-1}}-1)$ converges absolutely to $\sum_kc_kb^{-n_k}$. Write
$\rho_h=(b^h-1)/(b-1)$ for the repunit of length $h$. If $b\ge2$ is an integer and for every $M\ge1$
some gap $n_k-n_{k-1}$ is divisible by $M$ (for example $n_k=t\,k!$), then every prime not dividing $b$
divides $\rho_{n_k-n_{k-1}}$ for some $k$, and in particular infinitely many primes divide the gap
repunits $\rho_{n_k-n_{k-1}}$.
\end{proposition}

\begin{proof}
The first identity is the finite geometric sum
$\sum_{r=1}^hb^{-(m+r)}=b^{-(m+h)}(b^h-1)/(b-1)$. Absolute convergence follows because every raw
coefficient is bounded by $\sup|c_k|$, and grouping an absolutely convergent series preserves its
sum. Let $r\nmid b$ be prime. If $r\nmid b-1$, then $r\mid\rho_{r-1}$ by Fermat's theorem; if
$r\mid b-1$, then $\rho_r\equiv r\equiv0\pmod r$. In either case $r\mid\rho_{r(r-1)}$, and
$\rho_a\mid\rho_h$ whenever $a\mid h$. Lean declarations:
\lean{Phase18Synthetic}{76}{\ttt{geometric\_retirement\_weight}},
\lean{Phase18Synthetic}{146}{\ttt{grouped\_retirement\_hasSum}},
\lean{Phase18Synthetic}{290}{\ttt{raw\_prime\_support\_infinite}}.
\end{proof}

The Lean statement concerns the repunits, not the reduced raw coefficients, and the passage between
the two needs the form of the multipliers. If $c_k=1/m_k$ with a positive integer $m_k$, as for the
weights $p^{-sv_p(k)}$, then $b-1\mid b^h-1$ gives the raw coefficient $1/(m_k\rho_h)$, $h=n_k-n_{k-1}$,
whose reduced denominator $m_k\rho_h$ is divisible by $\rho_h$; so under the gap hypothesis every prime
not dividing $b$ occurs in some reduced raw denominator (a written step, not formalized). For general
multipliers this fails: with $c_k=\rho_h b^{-h}\in(0,1]$ every raw coefficient on the gap equals
$b^{-h}$, and only primes dividing $b$ occur. The gap hypothesis is also needed: a schedule with
finitely many distinct gaps involves only finitely many repunits.

\subsection{Exact denominators and the switch}

\begin{theorem}[Exact reduced denominators]\label{thm:den}\evidence{Kernel-verified}
For every $K\ge0$ and every $n\ge n_K$, the reduced denominator of $b^nU_K$ is
$p^{s\lfloor\log_pK\rfloor}$, where $\lfloor\log_p0\rfloor=0$. In particular the reduced denominator
of $R_n$ is $p^{s\lfloor\log_pK(n)\rfloor}$ for every $n\ge0$.
\end{theorem}

\begin{proof}
Let $J=\lfloor\log_pK\rfloor$ and $K\ge1$, $s\ge1$ (otherwise the value is an integer). Every
$k\le K$ has $v_p(k)\le J$, so $p^{sJ}b^nU_K=\sum_{k\le K}p^{s(J-v_p(k))}b^{n-n_k}$ is an integer.
Modulo $p$ only the indices $k=p^Jm$ with $p\nmid m$ survive, and $m\le K/p^J<p$, so there are
between $1$ and $p-1$ of them. Since $t\mid n_k$, each contributes $b^{n-n_k}\equiv b^n\not\equiv0$.
The count is a unit modulo $p$, so the numerator is prime to $p$. Lean declarations:
\lean{RetirementDenominator}{110}{\ttt{shiftedPrefix\_den}},
\lean{RetirementDenominator}{157}{\ttt{fullPrefix\_den}},
\lean{RetirementDenominator}{67}{\ttt{maximalLayer\_card}}.
\end{proof}

The qualitative analytic input is a fixed-support estimate for short orbits.

\begin{lemma}[Short orbits]\label{lem:short}\evidence{Kernel-verified}
Let $b,c\ge2$ be coprime integers. For every integer $d_0\ge0$ and every $\eta>0$ there is $M$ such
that for every rational $y$ whose reduced denominator $q$ divides a power of $c$ and satisfies
$M\le q\le T^{d_0}$,
\[
\Bigl|\sum_{j=0}^{T-1}e(b^jy)\Bigr|\le\eta T,\qquad e(\theta)=\exp(2\pi i\theta).
\]
\end{lemma}

This is \lean{SmoothOrbit}{207}{\ttt{fixed\_support\_orbit\_estimate}} (definition at
\lean{SmoothOrbit}{16}{\ttt{SmoothOrbitEstimate}}), proved there by elementary differencing. It is
also a qualitative consequence of Vandehey's quantitative bound \cite[Corollary~1.2]{Vandehey19}:
for a finite prime set $P$, $m$ supported on $P$, $\gcd(m,b)=\gcd(a,m)=1$, $\epsilon>0$ and
$T\ge m^\epsilon$, one has $|\sum_{j=1}^{T}e(ab^j/m)|\le C\,m^{-\delta}T(1+\log m)$ with
$\delta=\delta(\epsilon)>0$ and $C=C(\epsilon,b,P)$.


The short-orbit estimate gives normality from denominator information alone, at every prefix.

\begin{proposition}[All-prefix criterion]\label{prop:criterion}\evidence{Lean-supported paper proof}
Let $x=\sum_{m\ge0}a_mb^{-m}$ with rational $a_m$ satisfying $|a_m|\le C_0(m+1)^A$ for constants
$C_0,A\ge0$, put $R_n=\sum_{m\le n}a_mb^{n-m}$, and let $q_n$ be the reduced denominator of $R_n$.
Suppose that for some $\eta\in(0,1]$, some integer $c\ge2$ prime to $b$ and some $D_0>0$:
\begin{enumerate}[label=\textup{(\roman*)},nosep]
\item $\#\{m\le N:a_m\ne0\}=O(N^{1-\eta})$;
\item every $q_n$ divides a power of $c$, and $q_n\le C_0(n+1)^{D_0}$;
\item for every fixed $M$, $\#\{n<N:q_n<M\}=o(N)$.
\end{enumerate}
Then $x$ is normal to base $b$.
\end{proposition}

\begin{proof}
The Lean-supported step is \Cref{lem:short}. Fix an integer $h\ne0$ and $\eta_0>0$. For large $N$ put
$T=\lfloor N^{\eta/2}\rfloor$ and $W=\lceil(A+3)\log_b(2N)\rceil\le T$, and split $[0,N)$ into consecutive
blocks $[a,a+T)$, the last possibly shorter. Call a block \emph{clean} if no $m$ with $a_m\ne0$ lies in
$(a,a+T+W]$. Each such $m$ makes at most two blocks unclean, so by (i) the unclean blocks and the last
block cover $O(TN^{1-\eta})+T=o(N)$ indices.

On a clean block and for $0\le j<T$ we have $R_{a+j}=b^jR_a$, and the omitted terms of $b^{a+j}x$ have
positions $m>a+j+W$, so
\[
|b^{a+j}x-R_{a+j}|\le C_0\sum_{i>W}(2N+i)^Ab^{-i}\ll_{A,b,C_0}N^AW^Ab^{-W}=o(1).
\]
Since $\gcd(b,c)=1$, all $q_{a+j}$ equal $q_a$, and the reduced denominator $q'$ of $hR_a$ divides a power
of $c$ and satisfies $q_a/|h|\le q'\le q_a$. Choose an integer $d_0>2D_0/\eta+1$, so that
$q_a\le C_0(N+1)^{D_0}\le T^{d_0}$ for large $N$, and let $M_0$ be given by \Cref{lem:short} for $b$, $c$,
$d_0$ and $\eta_0$. On clean blocks with $q'\ge M_0$, $|\sum_{j<T}e(b^jhR_a)|\le\eta_0T$. Clean blocks with
$q'<M_0$ have $q_n<|h|M_0$ at each of their indices, so by (iii) they cover $o(N)$ indices. Hence
\[
\Bigl|\sum_{n<N}e(hb^nx)\Bigr|\le o(N)+\eta_0N+2\pi|h|N\cdot o(1),
\]
and $\limsup_N|N^{-1}\sum_{n<N}e(hb^nx)|\le\eta_0$ for every $\eta_0>0$. Weyl's criterion
\cite[Chapter~1, Theorem~2.1, p.~7]{KN} and \cite[Chapter~1, Theorem~8.1, p.~70]{KN} give normality.
\end{proof}

\begin{theorem}[Normality switch]\label{thm:switch}\evidence{Kernel-verified}
Assume in addition that $n_k\ge k^2$ for all $k$. Then $D_N(X_0)\to1$, and for every $s\ge1$,
$D_N(X_s)\to0$; equivalently $X_s$ is normal to base $b$ if and only if $s\ge1$. In particular this
holds for $n_k=tk^d$ with any integer $d\ge2$.
\end{theorem}

\begin{proof}
The statement is a Lean theorem (declarations listed at the end of this proof), with the definitions
of discrepancy and normality of \Cref{sec:intro}. We outline the argument.

Since $n_k\ge k^2$, $K(M)\le\sqrt M$. Positivity gives
$0\le b^nX-R_n\le\frac{b}{b-1}b^{-(n_{K(n)+1}-n)}$, and for fixed $W$ at most
$W(K(N+W)+1)=o(N)$ indices $n<N$ have an insertion in $(n,n+W]$. If $s=0$, every $R_n$ is an integer,
so all but $o(N)$ of the points $\fr{b^nX_0}$ lie in $[0,\frac{b}{b-1}b^{-W})$; as in the proof of
\Cref{thm:example}, $\liminf D_N(X_0)\ge1-\frac b{b-1}b^{-W}$ for every $W$. If $s\ge1$, apply
\Cref{prop:criterion} with $c=p$ and $\eta=1/2$: the coefficients are at most one, there are at most
$\sqrt N$ insertions below $N$, and by \Cref{thm:den} the reduced denominator of $R_n$ is
$p^{s\lfloor\log_pK(n)\rfloor}\le(n+1)^s$, which tends to infinity with $n$. For $n_k=tk^d$ the hypotheses
hold because $t\ge1$ and $d\ge2$. Lean declarations:
\lean{RetirementDenominator}{436}{\ttt{discrepancy\_switch}},
\lean{RetirementDenominator}{444}{\ttt{normal\_iff\_pos}} and, for $n_k=tk^d$,
\lean{RetirementDenominator}{456}{\ttt{polynomial\_discrepancy\_switch}}.
\end{proof}

\begin{corollary}[A normal Liouville family]\label{cor:liouville}\evidence{Lean-supported paper proof}
For every $s\ge1$, the number $L=\sum_{k\ge1}p^{-sv_p(k)}b^{-t(k+1)!}$ is normal to base $b$ and is a
Liouville number; in particular it is transcendental.
\end{corollary}

\begin{proof}
The schedule $n_k=t(k+1)!$ consists of strictly increasing multiples of $t$ with $n_k\ge k^2$, so
normality is the kernel-verified \Cref{thm:switch} (the verification of these two elementary
hypotheses is not formalized). The partial sum $U_k$ has reduced denominator
$Q_k=b^{n_k}p^{s\lfloor\log_pk\rfloor}\le b^{n_k}k^s$: indeed the newest term is the only one not
divisible by $b$ after clearing $b^{n_k}$, and \Cref{thm:den} gives the $p$-part. Moreover
$0<L-U_k\le\frac{b}{b-1}b^{-n_{k+1}}=\frac b{b-1}b^{-(k+2)n_k}$. For each fixed $\mu>0$ this is
smaller than $Q_k^{-\mu}$ for all large $k$, and the $U_k$ are distinct. Hence $L$ is a Liouville
number, and transcendental by Liouville's theorem.
\end{proof}

\subsection{The complete-period rate}

\begin{theorem}[Complete-period rate]\label{thm:rate}\evidence{Lean-supported paper proof}
Let $p\nmid b$ be an odd prime, $v=v_p(b^t-1)$, $s\ge1$, $d\ge s+2$, and $n_k=tk^d$. There is
$C=C(b,p,s,d)$ such that for all $N\ge2$,
\[
D_N(X_s)\le C\bigl(N^{-s/d}+N^{-(d-s-1)/d}+N^{-(d-1)/d}\log N\bigr).
\]
The exponent $\min(s,d-1-s)/d$ equals $1/3$ for $(s,d)=(1,3)$ and tends to $1/2$ along
$(s,d)=(s,2s+1)$.
\end{theorem}

\begin{proof}
The Lean-supported step is the denominator identity, \Cref{thm:den}; the rest follows the proof of
\Cref{thm:example}. Let $tK^d\le N<t(K+1)^d$ with $K\ge1$ (smaller $N$ are absorbed in $C$). On the
block $B_k=[n_k,\min(n_{k+1},N))$ the model $y_n=\fr{R_n}$ runs through consecutive points of
the orbit of a rational with reduced denominator $q_k=p^{s\lfloor\log_pk\rfloor}\in((k/p)^s,k^s]$.
The block has length at most $h_k=n_{k+1}-n_k\le t\,d\,2^{d-1}k^{d-1}$. By
\Cref{cor:period},
\[
\sum_{k\le K}H_k\le\sum_{k\le K}\Bigl(\frac{p^vh_k}{q_k}+q_k\Bigr)
\ll\sum_{k\le K}\bigl(k^{d-1-s}+k^s\bigr)\ll K^{d-s}+K^{s+1},
\]
using $d-1-s\ge1$. Take $W=\lceil2\log_b(2N)\rceil$ and $\epsilon=\frac b{b-1}b^{-W}\le2N^{-2}$; the
exceptional count is $E\le KW$, and $U=n_1=t$. \Cref{lem:blocks} and $N\asymp K^d$ give
$D_N\ll K^{-s}+K^{-(d-s-1)}+K^{1-d}\log N+N^{-2}$, which is the claim.
\end{proof}

\begin{remark}\label{rem:ratebeyond}
Complete periods give a rate only for $d\ge s+2$. For every $s\ge1$ and $d\ge2$ the short-orbit bound
of \Cref{lem:uniformorbit} gives $D_N(X_s)\ll N^{-\eta}$ with some $\eta=\eta(p,s,d)>0$; this is the case
$r=1$ of \Cref{thm:effective}. The exponent obtained there is very small (of order $2^{-2/\theta}$ with
$\theta\asymp1/s$), so \Cref{thm:rate} remains the better estimate when $d\ge s+2$. Bounding the final
incomplete period on each block by Gauss sums instead of trivially would improve the range of
\Cref{thm:rate}; we have not written that argument out.
\end{remark}

\section{Prescribing the discrepancy limit}\label{sec:program}

\begin{lemma}[Positive state programming]\label{lem:program}\evidence{Kernel-verified}
Let $0=n_0<n_1<\cdots$ be integers and $u_0=0,u_1,u_2,\ldots$ rationals in $[0,1)$. Define $R_0=0$,
$c_k\in(0,1]$ as the representative of $u_k-b^{n_k-n_{k-1}}R_{k-1}$ modulo $1$, and
$R_k=b^{n_k-n_{k-1}}R_{k-1}+c_k$. Then $R_k=b^{n_k}\sum_{i\le k}c_ib^{-n_i}$ and $\fr{R_k}=u_k$ for
every $k$. If every $u_k$ has reduced denominator a power of a prime $p$, so has every $c_k$.
\end{lemma}

\begin{proof}
By induction $R_k-u_k=b^{n_k-n_{k-1}}(R_{k-1}-u_{k-1})+(b^{n_k-n_{k-1}}u_{k-1}+c_k-u_k)\in\Z$. The
coefficient $c_k$ differs by an integer from $u_k-b^{n_k-n_{k-1}}u_{k-1}$, whose denominator divides
$\den(u_k)\den(u_{k-1})$. Lean declarations:
\lean{Phase18Synthetic}{383}{\ttt{programState\_fract}},
\lean{Phase18Synthetic}{438}{\ttt{programState\_eq\_prefix}},
\lean{Phase18Synthetic}{459}{\ttt{programmed\_prefix\_fract}},
\lean{Phase18Synthetic}{364}{\ttt{programCoeff\_mem}},
\lean{Phase18Synthetic}{426}{\ttt{programCoeffs\_prime\_power}}.
\end{proof}

For $x=\sum_kc_kb^{-n_k}$ as in \Cref{lem:program} and $n_k\le n<n_{k+1}$, the rational part after
$n$ shifts is $b^{n-n_k}R_k$, whose fractional part is $y_n=\fr{b^{n-n_k}u_k}$. Since
$0<c_i\le1$ and the positions are distinct,
\begin{equation}\label{eq:tail}
0<b^nx-b^{n-n_k}R_k\le\frac b{b-1}b^{-(n_{k+1}-n)} .
\end{equation}
Call block $k$ (the integers in $[n_k,n_{k+1})$) \emph{retired} if $u_k=0$ and \emph{active}
otherwise; block $0$ is retired.

\begin{proposition}[Modelled discrepancy]\label{prop:model}\evidence{Paper proof}
Let $n_K\le N<n_{K+1}$, let $Z$ be the number of indices $n<N$ in retired blocks, and suppose that
for every half-open arc $I$ the active indices $n<N$ satisfy
$|\#\{n:y_n\in I\}-(N-Z)|I||\le H$. Then for every integer $W\ge0$,
\[
\Bigl|D_N(x)-\frac ZN\Bigr|\le\frac{H+(K+1)W}{N}+\frac{2b}{b-1}b^{-W}.
\]
\end{proposition}

\begin{proof}
Retired models are $0$. For every half-open arc,
$\#\{n<N:y_n\in I\}-N|I|=(\text{active error})+Z(\mathbf 1_{0\in I}-|I|)$, which has modulus at most
$H+Z$, and equals at least $Z(1-|I|)-H$ for $I=[0,|I|)$. Letting $|I|\to0$ and using
\Cref{lem:arcs}, $|D(y)-Z/N|\le H/N$. By \eqref{eq:tail} at most $W$ indices of each block
$k\le K$ have $\dist{b^nx-y_n}>\frac b{b-1}b^{-W}$, and \Cref{lem:compare} completes the proof.
\end{proof}

\begin{theorem}[Every discrepancy limit]\label{thm:program}\evidence{Lean-supported paper proof}
Let $p\nmid b$ be an odd prime, $t=\ord_p(b)$ and $v=v_p(b^t-1)$. For $\delta\in[0,1]$ put $n_k=tk^3$,
$a_k=\lfloor k\delta\rfloor-\lfloor(k-1)\delta\rfloor\in\{0,1\}$, $q_k=p^{1+\lfloor\log_pk\rfloor}$, and
$u_k=0$ if $a_k=1$, $u_k=1/q_k$ if $a_k=0$. Let $x_\delta=\sum c_kb^{-n_k}$ be given by
\Cref{lem:program}. Then every $c_k\in(0,1]$ has reduced denominator dividing $q_k$, so a power of $p$
at most $pk$, and for all $N\ge1$,
\[
|D_N(x_\delta)-\delta|\le C_{b,p}N^{-1/3}
\]
with $C_{b,p}$ independent of $\delta$. Moreover the orbit measures
$N^{-1}\sum_{n<N}\delta_{\fr{b^nx_\delta}}$ converge weakly to $(1-\delta)\lambda+\delta\,\delta_0$, where
$\lambda$ is Lebesgue measure, and $x_\delta$ is irrational.
\end{theorem}

\begin{proof}
The Lean-supported step is \Cref{lem:program}. Since $q_{k-1}\mid q_k$, the coefficient $c_k$, which
differs by an integer from $u_k-b^{n_k-n_{k-1}}u_{k-1}$, has denominator dividing $q_k\le pk$. Let
$N\ge t$ and $tK^3\le N<t(K+1)^3$. On an active
block $k$ the models are consecutive points $\fr{b^j/q_k}$, with $k<q_k\le pk$, and
$h_k=n_{k+1}-n_k=t(3k^2+3k+1)\le7tk^2$. \Cref{cor:period} gives
$H\le\sum_{k\le K}(p^vh_k/q_k+q_k)\le(7tp^v+p)K(K+1)/2$, independently of the mask $(a_k)$.

The retired count is $Z=t+\sum_{1\le k<K}a_kh_k+a_K(N-n_K)$, and $N=t+\sum_{1\le k<K}h_k+(N-n_K)$. Put
$S_j=\sum_{i\le j}(a_i-\delta)$, so $|S_j|\le1$. Summation by parts with the increasing weights $h_k$
gives $|\sum_{1\le k<K}(a_k-\delta)h_k|\le2h_{K-1}$, and therefore
\begin{equation}\label{eq:Z}
|Z-\delta N|\le t+3h_K .
\end{equation}
\Cref{prop:model} with $W=\lceil\log_b(\frac b{b-1}(K+1))\rceil$ yields
\[
|D_N(x_\delta)-\delta|\le\frac{(7tp^v+p)K(K+1)/2+(K+1)W+t+3h_K}{N}+\frac2{K+1},
\]
which is $O_{b,p}(K^{-1})=O_{b,p}(N^{-1/3})$, uniformly in $\delta$. For $N<t$ use $D_N\le1$.

For weak convergence let $f$ be continuous on $\T$ and approximate it uniformly by step functions
on a fixed finite partition into arcs. The active indices have counts $(N-Z)|I|+O(H)$ on each arc,
the retired real points lie within $\frac b{b-1}b^{-W}$ of $0$ off $(K+1)W$ indices, and
$Z/N\to\delta$ by \eqref{eq:Z}. If $\delta<1$ the limit has a nonzero absolutely continuous part,
whereas the orbit measures of a rational number converge to the uniform measure on a finite cycle;
so $x_\delta$ is irrational. If $\delta=1$ every $u_k$ is $0$, every $c_k=1$, and the base-$b$ expansion
of $x_1$ has isolated digits $1$ separated by unbounded gaps, which is not eventually periodic.
\end{proof}

For example, with $b=2$, $p=3$ and $\delta=2/5$, the first five insertions are at
$2,16,54,128,250$, the targets are $1/3,1/3,0,1/9,0$, and the coefficients are
$1/3,1,2/3,1/9,5/9$.

\begin{remark}
The canonical numbers $x_\delta$ are irrational, but we do not know their algebraic nature. Altering
the mask at the sparse indices $k=2^j$ does not change the rate and produces continuum many
distinct realizations of each $\delta$; all but countably many are transcendental.
\end{remark}

\section{Lower and upper discrepancies and invariant measures}\label{sec:limitsets}

\begin{theorem}[All pairs of extreme limits]\label{thm:pairs}\evidence{Lean-supported paper proof}
Let $p,t,n_k,q_k$ be as in \Cref{thm:program} and $0\le\alpha\le\beta\le1$. There are rationals
$c_k\in(0,1]$ with reduced denominators dividing $q_k\le pk$, given by \Cref{lem:program} from targets
$u_k\in\{0,1/q_k\}$ and an explicit mask, such that $x=\sum_kc_kb^{-n_k}$ satisfies
$\liminf_ND_N(x)=\alpha$ and $\limsup_ND_N(x)=\beta$. More precisely $D_N(x)=\Lambda_N/N+O(N^{-1/3})$,
with $\Lambda_N$ defined in the proof.
\end{theorem}

\begin{proof}
Let $K_1=2$, $K_{j+1}=K_j^2$, and for $K_j\le k<K_{j+1}$ let $\theta_k=\alpha$ for odd $j$ and
$\theta_k=\beta$ for even $j$; put $\theta_1=\alpha$. Use the mask
$a_k=\lfloor\Theta_k\rfloor-\lfloor\Theta_{k-1}\rfloor$ with $\Theta_k=\sum_{i\le k}\theta_i$, and the
targets of \Cref{thm:program}. Then $|\sum_{i\le j}(a_i-\theta_i)|\le1$, and the argument for
\eqref{eq:Z} gives $|Z-\Lambda_N|\le t+3h_K$, where $\Lambda_N=\sum_{k\le K}\theta_k\ell_k$ and $\ell_k$
is the length of the part of block $k$ below $N$. Since the active-block bound is independent of
the mask, \Cref{prop:model} gives $D_N(x)=\Lambda_N/N+O(N^{-1/3})$. Each $\Lambda_N/N$ is a weighted
average of the $\theta_k$ up to $O(t/N)$, so lies in $[\alpha,\beta]+o(1)$. At $N=tK_{j+1}^3$ the
blocks before stage $j$ have total length at most $tK_j^3=o(N)$, so $\Lambda_N/N$ tends to $\alpha$ along
odd $j$ and to $\beta$ along even $j$.
\end{proof}

Let $\Prob$ be the space of Borel probability measures on $\T$ with the weak topology and $\Minv$
the closed subset of $T$-invariant measures. Fix nonzero trigonometric polynomials $f_1,f_2,\ldots$
with rational coefficients, dense in $C(\T)$, and put
$d(\mu,\nu)=\sum_i2^{-i}|\int f_id\mu-\int f_id\nu|/\|f_i\|_\infty$. Uniform approximation shows that
$d(\mu_n,\mu)\to0$ if and only if $\mu_n\to\mu$ weakly, and $\Prob$ is sequentially compact by a
diagonal argument and the Riesz representation theorem. Since each summand of $d$ is a convex
function of $\mu$, $d(\sum_r\gamma_r\mu_r,\nu)\le\max_rd(\mu_r,\nu)$ for convex weights $\gamma_r$. For $x\in\R$ let
$\mu_N^x=N^{-1}\sum_{n<N}\delta_{\fr{b^nx}}$ and $V(x)$ its set of weak limit points.

\begin{lemma}\label{lem:necessity}\evidence{Paper proof}
For every $x$, $V(x)$ is a nonempty compact connected subset of $\Minv$.
\end{lemma}

\begin{proof}
Compactness of $\Prob$ gives nonemptiness, and $V(x)$ is closed. For continuous $f$,
$\int f\circ T\,d\mu_N^x-\int fd\mu_N^x=(f(T^Nx)-f(x))/N\to0$, so limits are invariant. Since
$d(\mu_N^x,\mu_{N+1}^x)\le2/(N+1)$, connectedness follows from the general fact that the cluster set
of a sequence $z_N$ in a compact metric space with $d(z_N,z_{N+1})\to0$ is connected: if it were
the union of disjoint nonempty closed sets $A,B$ at distance $3r>0$, the sequence would eventually
avoid the set of points at distance at least $r$ from $A\cup B$ (a compact set containing no
cluster point), yet would pass infinitely often from the $r$-neighbourhood of $A$ to that of $B$ in
steps smaller than $r$, which is impossible.
\end{proof}

For $\mu\in\Prob$ let $D(\mu)=\sup_I|\mu(I)-|I||$ over half-open arcs; for an empirical measure this
is the discrepancy of \Cref{sec:prelim}.

\begin{lemma}\label{lem:Dcont}\evidence{Paper proof}
If $\mu_i\to\mu$ weakly in $\Prob$, then $D(\mu_i)\to D(\mu)$, even when $\mu$ has atoms.
\end{lemma}

\begin{proof}
As in \Cref{lem:arcs}, $D(\nu)$ is the supremum of the excess $\nu(I)-|I|$ over half-open arcs. Given
$\eta>0$, choose a rotated grid of $M\ge1/\eta$ equally spaced points that avoids the countably many
atoms of $\mu$. Every half-open arc is contained in an arc whose endpoints are grid points (or in the
whole circle), at a cost of at most $2\eta$ in length, so for every $\nu$ the maximum excess over the
finitely many grid arcs lies in $[D(\nu)-2\eta,D(\nu)]$. The masses of grid arcs converge, because
their endpoints are not atoms of $\mu$. Hence $\limsup|D(\mu_i)-D(\mu)|\le4\eta$.
\end{proof}

\begin{theorem}[Invariant-measure limit sets]\label{thm:limitsets}\evidence{Lean-supported paper proof}
A set $F\subset\Prob$ equals $V(x)$ for some real $x$ if and only if $F$ is a nonempty compact
connected subset of $\Minv$. For every integer $t\ge1$, every such $F$ is $V(x)$ for some
$x=\sum_{k\ge1}c_kb^{-tk^3}$ with rational $c_k\in(0,1]$ given by \Cref{lem:program}, whose state $u_k$
after the $k$th insertion is $0$ or a purely periodic rational with period at most $k$.
\end{theorem}

\begin{proof}
Necessity is \Cref{lem:necessity}. The Lean-supported step of the sufficiency proof is
\Cref{lem:program}, which programs the states.

\emph{Periodic approximation.} Let $\mu\in\Minv$, $\eta>0$, and let $g_1,\ldots,g_r$ be
Lipschitz functions with constant at most $\Lambda$ and $\|g_i\|_\infty\le1$. We find a rational $y$
with purely periodic base-$b$ expansion whose cycle measure $\pi$ satisfies
$|\int g_id\pi-\int g_id\mu|<\eta$. Fix integers $m$ and $L$ with $\Lambda b^{-m}+2m/L<\eta/3$, and let
$A_Lg(z)=L^{-1}\sum_{j<L}g(T^jz)$. By invariance $\int g_id\mu=\int A_Lg_i\,d\mu$. Partition the
range of $(A_Lg_1,\ldots,A_Lg_r)$ into Borel boxes of diameter $<\eta/3$, choose a point $z_\iota$ in
each preimage of positive measure, and replace the weights by rationals $M_\iota/M$ summing to one,
so that $|\sum_\iota(M_\iota/M)A_Lg_i(z_\iota)-\int g_id\mu|<2\eta/3$. Let $y$ be the purely periodic
number whose period is the concatenation, over $\iota$, of $M_\iota$ copies of the first $L$
base-$b$ digits of $z_\iota$. If a copy of the word of $z_\iota$ starts at offset $o$ of the period
and $0\le j\le L-m$, then $T^{o+j}y$ and $T^jz_\iota$ share their first $m$ digits, so
$\dist{T^{o+j}y-T^jz_\iota}\le b^{-m}$. The remaining $m$ offsets of each copy change $g_i$ by at
most $2$. Averaging over the period gives the claim; an
all-$(b-1)$ period represents $0\in\T$ and the same estimate holds on $\T$.

\emph{Chains.} Since $F$ is compact and connected, for every $\eta>0$ any two points of $F$ are
joined by a finite chain in $F$ with steps less than $\eta$ (the set of points reachable from a
given point is open and closed in $F$). Concatenating, for $r=1,2,\ldots$, chains with steps
$<1/r$ that start and end at a fixed point and visit a $1/r$-net of $F$, we obtain $\mu_j\in F$
with $d(\mu_j,\mu_{j+1})\to0$ whose cluster set is $F$.

\emph{Construction.} By the first step, applied to $g_i=f_i/\|f_i\|_\infty$ for $i\le j$, choose
periodic rationals $y_j\in[0,1)$ with cycle measures $\pi_j$, periods $P_j$, and
$d(\pi_j,\mu_j)<1/j+2^{1-j}$. Choose $K_1\ge\max(2,P_1)$ and $K_{j+1}\ge\max(K_j^2,P_{j+1})$. Let
$n_k=tk^3$, $u_k=0$ for $k<K_1$ and $u_k=y_j$ for $K_j\le k<K_{j+1}$, and let $x$ be given by
\Cref{lem:program}; then $u_k$ has period $P_j\le K_j\le k$. On block $k$ in stage $j$ the models are
consecutive points of the cycle of $y_j$; for a Lipschitz $g$ with $\|g\|_\infty\le1$, a piece of length
$\ell_k$ satisfies $|\sum g(y_n)-\ell_k\int gd\pi_j|\le2P_j\le2k$. Summing over $k\le K$, where
$tK^3\le N<t(K+1)^3$, gives $O(K^2)=O_t(N^{2/3})$. By
\eqref{eq:tail}, the real points are within $\frac b{b-1}b^{-W}$ of the models off $(K+1)W$ indices.
Hence $d(\mu_N^x,\nu_N)\to0$, where $\nu_N=N^{-1}\sum_k\ell_k\pi_{j(k)}$ (with $\pi=\delta_0$ before
stage $1$).

\emph{Limits.} Let $N_j=tK_{j+1}^3$ be the end of stage $j$. Stages before $j$ occupy the indices
below $tK_j^3$ and so have weight at most $(K_j/K_{j+1})^3\le K_j^{-3}$ in $\nu_{N_j}$, whence
$d(\nu_{N_j},\mu_j)\to0$. For $N_{j-1}<N\le N_j$,
$\nu_N$ is a convex combination of $\nu_{N_{j-1}}$ and $\pi_j$, both within $o(1)$ of $\mu_j$ because
$d(\mu_{j-1},\mu_j)\to0$; by convexity of balls $d(\nu_N,\mu_j)\to0$. Therefore every limit point of
$\mu_N^x$ lies in $F$, and every point of $F$, being a limit of a subsequence of $(\mu_j)$, is a limit
of the corresponding $\mu_{N_j}^x$.
\end{proof}

The symbolic realization of every nonempty compact connected set of invariant measures goes back to
Colebrook \cite{Colebrook} for base-$b$ expansions and to Sigmund \cite[Theorem~4]{Sigmund74} for systems
with the specification property; see also \cite{DGS}. The density of periodic orbit measures used in
the first step goes back to Parthasarathy \cite{Parthasarathy61}. The next remark shows that the
restriction to positive rational coefficients on a fixed schedule is, by itself, no restriction.

\begin{remark}[Qualitative content]\label{rem:qualitative}
Let $p$ be a prime and $0<n_1<n_2<\cdots$ any integers. For every $y\in[0,1)$ there are $m\ge0$ and
rationals $c_k\in(0,1)$ whose denominators are powers of $p$ such that
$x=\sum_kc_kb^{-n_k}$ satisfies $b^mx\equiv y\pmod1$. Indeed choose $m$ with
$x=(1+y)b^{-m}<b^{-n_1}$. If $0<r<b^{-n_k}$ is the residual before step $k$ (initially $r=x$), the
open interval of $c$ with $0<r-cb^{-n_k}<b^{-n_{k+1}}$ and $0<c<b^{n_k}r$ is nonempty and contains a
rational with $p$-power denominator; take it as $c_k$. Then $c_k<1$, the new residual lies in
$(0,b^{-n_{k+1}})$, and the residuals tend to $0$, so $x=\sum_kc_kb^{-n_k}$. The orbit of $x$ agrees
with that of $y$ after $m$ shifts, so $D_N(x)-D_N(y)=O(m/N)$ and $V(x)=V(y)$. Consequently, by the
classical realization above and \Cref{lem:Dcont}, every nonempty compact connected $F\subset\Minv$ is
$V(x)$ for such an $x$ on any schedule, even with a single denominator prime; taking
$F=\{(1-\theta)\lambda+\theta\delta_0:\alpha\le\theta\le\beta\}$, on which $D$ takes the values
$[\alpha,\beta]$, gives every pair of lower and upper limits, and $\alpha=\beta=\delta$ gives every limit.
\Cref{thm:program,thm:pairs,thm:limitsets} are therefore not new as existence statements. Their content
is the rate $O(N^{-1/3})$, uniform in $\delta$; the denominator bound $q_k\le pk$ with two explicit
targets in \Cref{thm:program,thm:pairs}; the bound $k$ on the period of the state in
\Cref{thm:limitsets}; and the explicit arithmetic schedule of targets.
\end{remark}

\section{Polynomial images below the schedule degree}\label{sec:quadratic}

Throughout this section $p\nmid b$ is an odd prime, $t=\ord_p(b)$, $v_0=v_p(b^t-1)\ge1$ and $d\ge2$ is an
integer; in \Cref{sec:quad-first,sec:quad-iso,sec:quad-null,sec:quad-proof}, which treat total degree two
as a model case, $d\ge3$.
For an integer $a\ge1$ put $w_a(i)=p^{-av_p(i)}$ and
\[
X_a=\sum_{i\ge1}w_a(i)\,b^{-ti^d},
\]
the number $X_a$ of \Cref{sec:switch} with $n_k=tk^d$. All the $X_a$ share the schedule $n_i=ti^d$. We
use the lifting of exponents in the form $v_p(b^{tm}-1)=v_0+v_p(m)$ for $m\ge1$ (\Cref{sec:prelim}).

\begin{theorem}[Quadratic images]\label{thm:quadratic}\evidence{Paper proof}
Let $S$ be a finite set of positive integers and let $P\in\Q[z_a:a\in S]$ be nonconstant of total
degree at most two. Then $P((X_a)_{a\in S})$ is normal to base $b$. Consequently:
\begin{enumerate}[label=\textup{(\roman*)},nosep]
\item every rational polynomial of degree one or two in $X_s$, $s\ge1$, is normal to base $b$;
\item the distinct nonconstant monomials of degree at most two in $(X_a)_{a\in S}$ are jointly normal:
the points $(\fr{b^nm_1},\ldots,\fr{b^nm_r})$, $n\ge0$, are uniformly distributed in $\T^r$ for the
values $m_1,\ldots,m_r$ of these monomials;
\item an algebraic $X_s$ would have degree at least three.
\end{enumerate}
\end{theorem}

\begin{corollary}\label{cor:flagship}\evidence{Paper proof}
The number $\xi^2$, with $\xi=\sum_{k\ge1}3^{-v_3(k)}2^{-2k^3}$, is normal to base $2$, and the binary
shifts of $(\xi,\xi^2)$ are uniformly distributed in $\T^2$.
\end{corollary}

\begin{proof}
Here $b=2$, $p=3$, $t=2$, $d=3$ and $\xi=X_1$; apply \Cref{thm:quadratic}(i) and (ii).
\end{proof}

\paragraph{Comparison with exponential schedules.} Mira Cult's note \cite[Theorem~1.1]{MiraPoly}
(a conventional proof, not refereed or formally verified) proves that every nonconstant rational
polynomial in $x=\sum_{k\ge1}q^{-k}b^{-c^k}$ is normal to base $b$ when $q>2$ and some prime divides $c$
and $q$ but not $b$; this includes the Stoneham numbers $\alpha_{b,c}$ with $c\ge3$, $\gcd(b,c)=1$.
There the schedule $c^k$ is exponential and the weight exponents increase with $k$, and the leading
$p$-adic part of a truncated power of $x$ is controlled by a nonintegral algebraic tail, which does not
vanish. On the polynomial schedule $ti^d$ with valuation weights the analogous leading residue of
$R_n(\xi^2)$ vanishes on a set of shifts of positive lower density (\Cref{sec:intro}), so no single
surviving pole is available; the proof below replaces it by the measure-theoretic statement that the
limiting mass has null level sets (\Cref{prop:nulllevel}). Neither result contains the other.

\begin{remark}[An earlier bound]\label{rem:history}
Before the proof of \Cref{thm:quadratic} was found, an exact enumeration of the cells of $u$ on which
the normalized prefix of $\xi^2$ (see \Cref{lem:freeze} below) vanishes modulo $3^{21}$ gave the bound
$\limsup_ND_N(\xi^2;2)<2.94\cdot10^{-10}$. \Cref{cor:flagship} supersedes that
bound. The theorem gives no rate.
\end{remark}

The proof occupies the rest of this section. Products of two series are expanded into ordered pairs of
indices, and we write $R_n(Z)$ for the shifted prefix of a series $Z=\sum_ma_mb^{-m}$, as in
\Cref{prop:criterion}. For a product the coefficient at $b^{-m}$ is the sum of the products of
coefficients over the ordered pairs $(i,j)$ with $n_i+n_j=m$.

\subsection{Normalization and the limiting mass}\label{sec:quad-first}

For $n\ge t$ let
\[
K(n)=\lfloor(n/t)^{1/d}\rfloor,\qquad J=J(n)=\lfloor\log_pK(n)\rfloor,\qquad u=u(n)=\frac{n}{tp^{dJ}} .
\]
Then $p^J\le K(n)<p^{J+1}$, so $u\in[1,p^d)$. Let $\Gamma$ be the set of positive rationals $a/p^h$ with
$h\ge0$ and $p\nmid a$, and call $\lambda(a/p^h)=h$ the \emph{loss}. For integers $a,c,D\ge1$ define, for
$u\in[1,p^d)$,
\[
M_{a,c}(u)=\sum_{\substack{(z,y)\in\Gamma^2\\ z^d+y^d\le u}}p^{a\lambda(z)+c\lambda(y)},\qquad
M^{\mathrm{lin}}_D(u)=\sum_{\substack{z\in\Gamma\\ z^d\le u}}p^{D\lambda(z)} .
\]
The conditions force $z,y<p$, so for each $L$ only finitely many terms are nonzero modulo $p^L$. The
sums therefore converge in $\Z_p$, and each reduction modulo $p^L$ is a right-continuous step function
of $u$ with finitely many jumps.

\begin{lemma}[Freezing]\label{lem:freeze}\evidence{Paper proof}
Let $a,c,D,L\ge1$ and let $n$ satisfy $J(n)\ge L+L/d$. Then
\[
p^{(a+c)J}b^{-n}R_n(X_aX_c)\equiv M_{a,c}(u)\pmod{p^L},\qquad
p^{DJ}b^{-n}R_n(X_D)\equiv M^{\mathrm{lin}}_D(u)\pmod{p^L}.
\]
Moreover $R_n(X_a)$ and $R_n(X_aX_c)$ are integral at every prime other than $p$, and
$v_p(R_n(X_a))\ge-aJ$, $v_p(R_n(X_aX_c))\ge-(a+c)J$.
\end{lemma}

\begin{proof}
By definition $b^{-n}R_n(X_aX_c)=\sum w_a(i)w_c(j)b^{-ti^d-tj^d}$ over the pairs with $t(i^d+j^d)\le n$.
Such indices satisfy $i,j\le K(n)<p^{J+1}$, so $\lambda_i=J-v_p(i)\ge0$ and
$p^{(a+c)J}w_a(i)w_c(j)=p^{a\lambda_i+c\lambda_j}$. This proves the integrality and the valuation bounds,
since $b^{n-n_i-n_j}$ is an integer. Write $z=i/p^J$, $y=j/p^J$; then $z,y\in\Gamma$ with losses
$\lambda_i,\lambda_j$, and $t(i^d+j^d)\le n$ if and only if $z^d+y^d\le u$. A term is nonzero modulo $p^L$
only if $a\lambda_i+c\lambda_j<L$, so $\lambda_i,\lambda_j<L\le J$; conversely every pair in $\Gamma^2$ with
losses less than $L$ arises from integers $i=p^Jz$, $j=p^Jy$. For such terms
$v_p(b^{ti^d}-1)=v_0+d(J-\lambda_i)\ge dL/d=L$, so $b^{-ti^d}\equiv1\pmod{p^L}$, and likewise for $j$.
The linear case is the same with one index.
\end{proof}

\subsection{Isolated jumps}\label{sec:quad-iso}

Fix an integer $D\ge2$ and integers $A_a$, $1\le a\le D/2$, not all zero, and $A_{\mathrm{lin}}$. Let
\[
M(u)=\sum_{1\le a\le D/2}A_aM_{a,D-a}(u)+A_{\mathrm{lin}}M^{\mathrm{lin}}_D(u).
\]
Let $a_*$ be the least $a$ with $A_a\ne0$ and $g=v_p(A_{a_*})$. It is convenient to use the variable
$x=(u-1)^{1/d}\in(0,X_{\max})$, $X_{\max}=(p^d-1)^{1/d}<p$; we write $M(x)$ for $M(1+x^d)$. A pair
$(z,y)\in\Gamma^2$ with $z^d+y^d>1$ contributes to $M(x)$ exactly for $x\ge x_{z,y}=(z^d+y^d-1)^{1/d}$, the
\emph{endpoint} of the pair; pairs with $z^d+y^d\le1$ contribute for all $x$. Similarly a linear term
$z>1$ contributes for $x\ge(z^d-1)^{1/d}$. A pair with $z=1$ has endpoint $y$.

For $H\ge1$ put $P_H=a_*H+g+1$. A \emph{primitive centre of level $H$} is a point $x_0=c/p^H\in(0,X_{\max})$
with $p\nmid c$.

\begin{lemma}[Primitive jumps]\label{lem:primjump}\evidence{Paper proof}
Let $H>g$. Modulo $p^{P_H}$, every jump of $M$ comes from one of the following: \textup{($\alpha$)} a pair
with a coordinate equal to $1$, whose other coordinate has loss at most $H$; \textup{($\beta$)} a pair
$(z,y)$ with $z\ne1\ne y$ and $\lambda(z)+\lambda(y)\le H$; \textup{($\gamma$)} a linear term with
$\lambda(z)\le H/2+g/D$. At a primitive centre $x_0$ of level $H$ the jumps of type $(\alpha)$ add up to
$A_{a_*}(p^{a_*H}+p^{(D-a_*)H})$, which is nonzero modulo $p^{P_H}$, and every other jump of type
$(\alpha)$ lies at distance at least $p^{-H}$ from $x_0$.
\end{lemma}

\begin{proof}
For the monomial of first order $a$, a pair of losses $(h,k)$ contributes $A_ap^{ah+(D-a)k}$. If $a=a_*$
this is nonzero modulo $p^{P_H}$ only if $a_*h+(D-a_*)k\le a_*H$, which forces $h+k\le H$ since
$D-a_*\ge a_*$. If $a>a_*$ and $A_a\ne0$, it is nonzero only if $ah+(D-a)k\le a_*H+g-v_p(A_a)$, which
forces $h+k\le(a_*H+g)/a<H$ for $H>g$. A linear term contributes $A_{\mathrm{lin}}p^{D\lambda(z)}$, nonzero
only if $D\lambda(z)\le a_*H+g$, and $a_*\le D/2$. Endpoints of pairs with a coordinate $1$ are the other
coordinates $y$, which lie in $p^{-H}\Z$ when $\lambda(y)\le H$; every point of $p^{-H}\Z$ other than
$x_0$ is at distance at least $p^{-H}$ from $x_0$, and the pair $(1,1)$ has endpoint $1\ne x_0$. At $x_0$
itself the pairs $(1,x_0)$ and $(x_0,1)$ are distinct; for the monomial $a_*$ they contribute
$A_{a_*}(p^{(D-a_*)H}+p^{a_*H})$, of valuation exactly $g+a_*H$ (if $a_*=D-a_*$ this is
$2A_{a_*}p^{a_*H}$, and $p$ is odd), and for $a>a_*$ their contributions have valuation at least
$v_p(A_a)+aH>a_*H+g$.
\end{proof}

Fix $\epsilon\in(0,1/4)$ and put $\delta_H=\epsilon/(Hp^H)$. A primitive centre $x_0$ of level $H$ is
\emph{isolated} if no endpoint of type $(\beta)$ or $(\gamma)$ lies within distance $\delta_H$ of $x_0$.
By \Cref{lem:primjump}, at an isolated centre $M\bmod p^{P_H}$ is constant on $(x_0-\delta_H,x_0)$ and on
$[x_0,x_0+\delta_H)$, and the two constants differ.

The equidistribution estimate behind the counting uses the third-derivative test.

\begin{lemma}[Third derivatives]\label{lem:third}\evidence{Paper proof}
There is an absolute constant $C_1$ with the following property. Let $\mathcal I$ be a set of at most $M$
consecutive integers, $M\ge1$, and let $g$ be real with three continuous derivatives on its convex hull,
$\lambda\le|g'''|\le2\lambda$ there for some $\lambda>0$. Then
$|\sum_{m\in\mathcal I}e(g(m))|\le C_1(M\lambda^{1/6}+M^{1/2}\lambda^{-1/6})$.
\end{lemma}

\begin{proof}
If $\lambda\ge1$ or $M\le\lambda^{-1/3}$ the bound is trivial, since then $M\le M\lambda^{1/6}$ or
$M\le M^{1/2}\lambda^{-1/6}$. Otherwise put $Q=\lfloor\lambda^{-1/3}\rfloor$, so $1\le Q\le M$. The
Weyl--van der Corput inequality \cite[Lemma~2.5 and (2.3.2), p.~10]{GK91} bounds $|\sum e(g(m))|^2$ by
$\frac{2M}{Q}\sum_{|q|<Q}|\sum_me(g(m+q)-g(m))|$, the inner sum over $m,m+q\in\mathcal I$. For $q\ne0$ the
function $g(m+q)-g(m)$ has second derivative $\int_m^{m+q}g'''$ of constant sign and modulus between
$|q|\lambda$ and $2|q|\lambda$, so the second-derivative test \cite[Theorem~2.2, p.~8]{GK91} bounds the inner sum
by $C(M(|q|\lambda)^{1/2}+(|q|\lambda)^{-1/2})$. Summing over $q$ gives
$|\sum e(g(m))|^2\ll M^2/Q+M^2(Q\lambda)^{1/2}+M(Q\lambda)^{-1/2}\ll M^2\lambda^{1/3}+M\lambda^{-1/3}$.
\end{proof}

\begin{lemma}[Isolated centres]\label{lem:isolated}\evidence{Paper proof}
Let $\mathcal K\subset(0,X_{\max})$ be a compact interval. There are $\epsilon\in(0,1/4)$ and $c_0>0$,
depending on $p,d,\mathcal K$, such that for every interval $I\subseteq\mathcal K$ of positive length and
every $H\ge H_0(I,M)$, at least $c_0|I|p^H$ isolated primitive centres of level $H$ lie in $I$.
\end{lemma}

\begin{proof}
There are at least $(1-1/p)|I|p^H-2$ primitive centres of level $H$ in $I$. Since $\delta_H<p^{-H}/2$,
each endpoint is within $\delta_H$ of at most one point of $p^{-H}\Z$, so it spoils at most one centre.
Endpoints of type $(\gamma)$ number $O(p^{H/2+g/D+1})=o(p^H)$. Let $I^+$ be the
$\delta_H$-neighbourhood of $I$, and fix a compact interval $\mathcal K'\subset(0,X_{\max})$ whose
interior contains $\mathcal K$, so that $I^+\subseteq\mathcal K'$ for large $H$. Let $B_H(I)$ be the number of pairs $(z,y)$ of type $(\beta)$
whose endpoint lies in $I^+$ and satisfies $\dist{p^Hx_{z,y}}\le\epsilon/H$, that is, lies within
$\delta_H$ of $p^{-H}\Z$. We show
\begin{equation}\label{eq:collision}
B_H(I)\le C_2\,\epsilon|I|p^H+o(p^H),
\end{equation}
with $C_2$ depending on $p,d,\mathcal K$ and the $o(p^H)$ uniform in $I\subseteq\mathcal K$. Choosing
$\epsilon=\min(1/8,(1-1/p)/(2C_2))$ then leaves at least $\frac12(1-1/p)|I|p^H-o(p^H)$ isolated centres,
which proves the lemma with $c_0=(1-1/p)/4$.

Write each pair as $(z,y)$ with $v=\lambda(z)\ge k=\lambda(y)$, counting both orders, so that $y$ is the
coarser coordinate; we let $z$ range over all of $p^{-v}\Z\cap(0,p)$, which only enlarges the count. For
fixed $y$ the endpoint is $f_y(z)=(z^d+B)^{1/d}$ with $B=y^d-1$; we call the set of $z$ for fixed $y$ a
\emph{row}. The rows of loss $k$ number at most $p^{k+1}$, and a row of loss $v$ has at most $p^{v+1}$
points. We use
\[
\sum_{\substack{k\le v\\v+k\le H}}p^{v+k}\ll_pHp^H,\qquad
\sum_{\substack{k\le v\\v+k\le H}}p^{v}\ll_pp^H,\qquad
\sum_{\substack{k\le v\\v+k\le H}}p^{k}\ll_pH^2p^{H/2}.
\]

\emph{Removals.} Charge every removed point as a collision. The rows with $v<2H/5$ contain
$O(H^2p^{4H/5})$ points in all. Put $\eta=H^{-6}$. Since $|y^d-1|\ge|y-1|$ for $y>0$, rows with
$|B|\le\eta$ have $0<|a-p^k|\le\eta p^k$ for $y=a/p^k$, hence there are at most $2\eta p^k$ of them for
each $k$, and they contain $O(\eta Hp^H)$ points. In each remaining row remove the points with
$z\le\eta$ or $|(d-2)B-(d+1)z^d|\le\eta^d$. Since $w^d-z^d\ge(w-z)^d$ for $w\ge z\ge0$, these conditions
define at most two intervals of total length $O_d(\eta)$, which contain $O(\eta p^v+1)$ grid points. The
total removal is $O(H^2p^{4H/5}+\eta Hp^H+H^2p^{H/2})=o(p^H)$.

\emph{Derivatives.} For $z>0$ with $f_y(z)>0$,
\[
f_y'=z^{d-1}f_y^{1-d},\qquad f_y''=(d-1)Bz^{d-2}f_y^{1-2d},\qquad
f_y'''=(d-1)Bz^{d-3}f_y^{1-3d}\bigl((d-2)B-(d+1)z^d\bigr).
\]
On the retained points with $f_y(z)\in\mathcal K'$ we have $z<p$, $|B|\le p^d$ and $f_y$ bounded above
and below, so $c_{\mathcal K}\eta^{2d-2}\le|f_y'''|\le C_{\mathcal K}$, using $d\ge3$. Moreover
$f_y''''=(d-1)Bz^{d-4}f_y^{1-4d}\,\Pi(z)$ with $\Pi$ a polynomial of degree $2d$ whose leading coefficient
is $(d+1)(d+2)\ne0$. Since $f_y$ is increasing, $\{z:f_y(z)\in I^+\}$ is an interval, and the retained
part of it splits into $O_d(1)$ intervals on which $f_y'''$ has constant sign and is monotone, hence into
$O(\log H)$ \emph{pieces} on which $\Lambda\le|f_y'''|\le2\Lambda$ for some $\Lambda\ge c_{\mathcal K}\eta^{2d-2}$.

\emph{Equidistribution in a piece.} Put $T=p^v$ and, for an integer $1\le l\le H^4$,
$g_l(m)=lp^Hf_y(m/T)$ on the integers $m$ of a piece, at most $pT$ of them. Then $\lambda_l\le|g_l'''|\le
2\lambda_l$ with $\lambda_l=lp^H\Lambda/T^3$. By \Cref{lem:third},
\begin{gather*}
\Bigl|\sum_me(g_l(m))\Bigr|\ll_pT\bigl(\lambda_l^{1/6}+T^{-1/2}\lambda_l^{-1/6}\bigr),\\
\lambda_l^{1/6}\ll H^{O(1)}p^{H/6-v/2}\le H^{O(1)}p^{-H/30},\qquad
T^{-1/2}\lambda_l^{-1/6}=(lp^H\Lambda)^{-1/6}\le H^{O(1)}p^{-H/6},
\end{gather*}
using $v\ge2H/5$, $l\le H^4$ and $\Lambda\ge H^{-O(1)}$. The Erdős--Turán inequality
\cite[Chapter~2, Theorem~2.5, p.~112]{KN} with $H^4$ frequencies bounds the number of $m$ in the piece
with $\dist{p^Hf_y(m/T)}\le\epsilon/H$ by
\[
\frac{2\epsilon}H\#\{m\}+O\Bigl(\frac{T}{H^4}+\sum_{l\le H^4}\frac1l\Bigl|\sum_me(g_l(m))\Bigr|\Bigr)
=\frac{2\epsilon}H\#\{m\}+O(TH^{-4}),
\]
since the arc $\{\theta\in\T:\dist\theta\le\epsilon/H\}$ has length $2\epsilon/H$. Summed over the
$O(\log H)$ pieces of a row, the error is $O(TH^{-3})$, and summed over the rows it is
$O(H^{-3}\sum p^{k+v+1})=O(p^H/H^2)$.

\emph{Main term.} For fixed losses $(v,k)$ the number of all grid points
$(z,y)\in(p^{-v}\Z\times p^{-k}\Z)\cap(0,p)^2$ with endpoint in $I^+=[x_1,x_2]$ is the number of lattice
points with $1+x_1^d<z^d+y^d\le1+x_2^d$, up to the endpoint convention. The area of
$\{z,y>0:z^d+y^d\le r^d\}$ is a constant times $r^2$, so the region has area
$O((1+x_2^d)^{2/d}-(1+x_1^d)^{2/d})=O_{\mathcal K}(|I^+|)$. Each sublevel set lies under a decreasing
graph, so counting by columns and comparing monotone sums with integrals gives
$p^{v+k}\cdot\mathrm{area}+O(p^{v+1}+p^{k+1})$ for its lattice count. Hence the main terms sum to at most
\[
\frac{4\epsilon}H\sum_{\substack{k\le v\\v+k\le H}}\Bigl(C|I^+|p^{v+k}+O(p^v+p^k)\Bigr)
\le C_2\epsilon|I|p^H+O(\epsilon p^H/H)+O(\delta_Hp^H).
\]
Together with the removals and the error terms this proves \eqref{eq:collision}.
\end{proof}

\subsection{Null level sets}\label{sec:quad-null}

\begin{proposition}[Null level sets]\label{prop:nulllevel}\evidence{Paper proof}
For every $w\in\Z_p$ the set $\{u\in[1,p^d):M(u)=w\}$ has Lebesgue measure zero.
\end{proposition}

\begin{proof}
Fix a compact interval $\mathcal K\subset(0,X_{\max})$, let $\epsilon$ and $c_0$ be as in
\Cref{lem:isolated}, and let $G_H$ be the union of the intervals $(x_0-\delta_H,x_0+\delta_H)$ over the
isolated primitive centres $x_0\in\mathcal K$ of level $H$.

\emph{Almost every point of $\mathcal K$ lies in infinitely many $G_H$.} Let $I\subseteq\mathcal K$ be an
interval. The intervals about isolated centres in the middle half of $I$ are disjoint and contained in
$I$ for large $H$, so by \Cref{lem:isolated}
\[
|I\cap G_H|\ge\tfrac12c_0|I|p^H\cdot2\delta_H=c_0\epsilon|I|/H\qquad(H\ge H_1(I)).
\]
Let $N_H$ be the $\delta_H$-neighbourhood of all of $p^{-H}\Z$. For $n<m$ with $|I|p^n\ge1$, at most
$|I|p^n+2\le3|I|p^n$ intervals of $N_n$ meet $I$, and each meets at most $2\delta_np^m+2$ intervals of
$N_m$, of length $2\delta_m$; hence
\[
|I\cap G_n\cap G_m|\le|I\cap N_n\cap N_m|\le3|I|p^n(2\delta_np^m+2)2\delta_m
\le12|I|\Bigl(\frac{\epsilon^2}{nm}+\frac{\epsilon p^{n-m}}m\Bigr),
\]
and $|I\cap G_n|\le6\epsilon|I|/n$. For $L\le H\le U$ with $L\ge H_1(I)$ put $S=\sum_{H=L}^U\mathbf1_{G_H}$
and $A=\sum_{H=L}^U1/H$. Then $\int_IS\ge c_0\epsilon|I|A$ and, using $\sum_{n<m}p^{n-m}/m\le A/(p-1)$,
$\int_IS^2\le|I|(12\epsilon^2A^2+C\epsilon A)$. By the Cauchy--Schwarz inequality
$|I\cap\{S>0\}|\ge(\int_IS)^2/\int_IS^2$, and letting $U\to\infty$, so that $A\to\infty$,
\[
\Bigl|I\cap\bigcup_{H\ge L}G_H\Bigr|\ge\frac{c_0^2}{12}|I| .
\]
The left side decreases in $L$, so this holds for every $L$, with a constant independent of $I$ and $L$.
Hence the set $\mathcal K\setminus\bigcup_{H\ge L}G_H$ has density at most $1-c_0^2/12$ in every interval
of $\mathcal K$, so it has no density point in the interior of $\mathcal K$, and by the Lebesgue density
theorem it is null. Taking the union over $L$ proves the claim.

\emph{Level sets.} The set $\Lambda_w=\{x\in\mathcal K:M(x)=w\}$ is the intersection over $L$ of the
finite unions of intervals $\{M\equiv w\bmod p^L\}$, hence Borel. Let $x\in\Lambda_w$ lie in $G_H$, within
$\delta_H$ of an isolated centre $x_0$ with $H>g$. The interval of radius $2\delta_H$ about $x$ contains
$(x_0-\delta_H,x_0+\delta_H)$. On the two halves of this interval $M\bmod p^{P_H}$ takes two different
constant values, so one whole half, of length $\delta_H$, misses $\Lambda_w$, and $\Lambda_w$ has relative
measure at most $3/4$ in the interval of radius $2\delta_H$ about $x$. If $x$ lies in infinitely many
$G_H$, it is therefore not a density point of $\Lambda_w$. Since almost every point of $\Lambda_w$ is a
density point and almost every point of $\mathcal K$ lies in infinitely many $G_H$, $\Lambda_w$ is null.
Exhausting $(0,X_{\max})$ by countably many compact intervals and returning to $u=1+x^d$, a $C^1$
diffeomorphism onto $(1,p^d)$, proves the proposition.
\end{proof}

\subsection{Proof of the quadratic theorem}\label{sec:quad-proof}

\begin{proof}[Proof of \Cref{thm:quadratic}]
\emph{Reduction.} Let $P$ have rational coefficients with common denominator $Q_0$, and write
$Q_0=Q_bQ_1$ with $\gcd(Q_1,b)=1$ and $Q_b\mid b^e$ for some $e$. Then $b^eP=Y/Q_1$ with $Y$ an integral
polynomial in the $X_a$, and $b^eP$ is normal if and only if $P$ is, since multiplying by $b^e$ shifts the
orbit by $e$ steps. If $r=\alpha/q$ in lowest terms, then $r/Q_1$ has reduced denominator between $q$ and
$qQ_1$, dividing $qQ_1$. Hence it suffices to verify the hypotheses of \Cref{prop:criterion} for the
series of $Y$, with $c=p$, and then for $Y/Q_1$ with $c=pQ_1$. We therefore assume that $P$ has integer
coefficients and verify (i)--(iii) for $Y=P((X_a))$.

\emph{Hypotheses (i) and (ii).} The positions of nonzero coefficients of $Y$ through $N$ lie in
$\{0\}\cup\{ti^d\le N\}\cup\{t(i^d+j^d)\le N\}$, which has $O(N^{2/d})$ elements; since $d\ge3$, (i) holds
with $\eta=1-2/d$. A coefficient at $b^{-m}$ is bounded by the sum of $|P|$'s coefficients times
$1+\#\{(i,j):t(i^d+j^d)=m\}\le1+m$. By \Cref{lem:freeze}, $q_n$ divides $p^{DJ}$, where $D$ is the largest
weight defined below, and $p^{DJ}\le K(n)^D\le(n+1)^D$.

\emph{Weights.} Give $X_a$ weight $a$ and $X_aX_c$ weight $a+c$, and let $D$ be the largest weight of a
monomial with nonzero coefficient in $P$. By \Cref{lem:freeze} every monomial of weight at most $D-1$
contributes to $p^{DJ}b^{-n}R_n(Y)$ a term of valuation at least $J$.

\emph{Case 1: the only monomial of weight $D$ is linear,} $A_{\mathrm{lin}}X_D$. By \Cref{thm:den}, $R_n(X_D)$
has $p$-adic valuation exactly $-DJ$, so $v_p(R_n(Y))=v_p(A_{\mathrm{lin}})-DJ$ for all large $n$. Hence
$q_n\to\infty$ and (iii) holds.

\emph{Case 2: some monomial $X_aX_{D-a}$ has a nonzero coefficient $A_a$} ($1\le a\le D/2$). Let $M$ be
the mass of the preceding subsection built from these coefficients and the coefficient
$A_{\mathrm{lin}}$ of $X_D$ (zero if absent). Fix $M_1\ge1$ and $L\ge1$. For $n$ with $J\ge L(1+1/d)$,
\Cref{lem:freeze} and the preceding remark give
\[
p^{DJ}b^{-n}R_n(Y)\equiv M(u(n))\pmod{p^{L}} .
\]
If $M(u(n))\not\equiv0\pmod{p^{L}}$, the left side has valuation at most $L-1$, so
$q_n\ge p^{DJ-L+1}$, which exceeds $M_1$ once $J$ is large. Therefore, for all large $n$, $q_n<M_1$ implies
$u(n)\in E_L=\{u:M(u)\equiv0\pmod{p^L}\}$. The set $E_L$ is a finite union of, say, $\kappa_L$ intervals.
The integers $n$ with $J(n)=j$ form the interval $[tp^{dj},tp^{d(j+1)})$, and those among them with
$u(n)\in E_L$ number at most $tp^{dj}|E_L|+\kappa_L$. If $tp^{dJ_0}\le N<tp^{d(J_0+1)}$, summing over
$j\le J_0$ gives
\[
\#\{n<N:q_n<M_1\}\le O_{L,M_1}(1)+\frac{tp^{d(J_0+1)}}{p^d-1}|E_L|+\kappa_L(J_0+1)
\le\Bigl(\frac{p^d}{p^d-1}|E_L|+o(1)\Bigr)N .
\]
The sets $E_L$ decrease to $\{M=0\}$, which is null by \Cref{prop:nulllevel}, so $|E_L|\to0$, and (iii)
follows. \Cref{prop:criterion} proves normality.

\emph{Consequences.} For (ii), every nonzero integer combination of distinct nonconstant monomials of
degree at most two is a nonconstant polynomial of total degree at most two, hence normal; uniform
distribution of the vector of shifts follows from Weyl's criterion in several dimensions \cite[Chapter~1,
Theorem~6.2, p.~49]{KN}. For (iii), if $X_s$ were algebraic of degree at most two, a rational polynomial
of degree at most two would vanish at $X_s$, and $0$ is not normal.
\end{proof}

\begin{remark}[Scope of the quadratic argument]\label{rem:quadscope}
The argument above gives no rate. For an $r$-fold product the number of grid tuples of total loss at
most $H$ is of order $H^{r-1}p^H$, and a single primitive family isolated at radius $1/(H^{r-1}p^H)$
covers measure of order $H^{1-r}$, which is summable when $r\ge3$; neither the second-moment step nor a
Borel--Cantelli argument then reaches full measure, and when $p\mid r$ the basic jump $rp^{sH}$ also
loses a factor $p$. The following subsections replace the single family by families built on dominant
anchors and prove normality in every degree $r<d$ together with an explicit rate. Normality of $\xi^3$
lies on the boundary $d=r$ (\Cref{sec:boundary}).
\end{remark}

\subsection{The \texorpdfstring{$r$}{r}-fold mass}\label{sec:rfold}

Fix integers $s\ge1$ and $r\ge1$ and write $C=X_s$. For an ordered tuple $z=(z_1,\ldots,z_r)\in\Gamma^r$ put
$S(z)=z_1^d+\cdots+z_r^d$ and $\lambda(z)=\lambda(z_1)+\cdots+\lambda(z_r)$, the \emph{level} of $z$. For
$u\in[1,p^d)$ let
\[
M_r(u)=\sum_{z\in\Gamma^r,\ S(z)\le u}p^{s\lambda(z)} .
\]
As for $M_{a,c}$, every coordinate is below $p$, only tuples of level at most $\lfloor(L-1)/s\rfloor$
matter modulo $p^L$, the sum converges in $\Z_p$, and each reduction modulo $p^L$ is a right-continuous
step function with finitely many jumps. The defining sum converges in $\Z_p$ for every real $u$ (for
each level only finitely many tuples have $S(z)\le u$), and we use it on all of $\R$, so that intervals
near $1$ or $p^d$ cause no boundary effects. Since the integers have loss $0$, $M_1(u)\equiv\lfloor
u^{1/d}\rfloor\not\equiv0\pmod p$ for $u\in[1,p^d)$.

\begin{lemma}[Freezing in degree $r$]\label{lem:freezer}\evidence{Paper proof}
Let $L\ge1$, let $P_0\in\Z[x]$ have degree $r$ and leading coefficient $A$, and let $n$ satisfy
$J(n)\ge J_0(L)=\lceil L/s+L/d\rceil$. Then $R_n(P_0(C))$ is integral at every prime other than $p$,
$v_p(R_n(P_0(C)))\ge-rsJ$, and
\[
p^{rsJ}b^{-n}R_n(P_0(C))\equiv A\,M_r(u)\pmod{p^L}.
\]
\end{lemma}

\begin{proof}
As in \Cref{lem:freeze}, $b^{-n}R_n(C^j)$ is the sum of $\prod_cp^{-sv_p(i_c)}b^{-ti_c^d}$ over ordered
$j$-tuples of indices with $t(i_1^d+\cdots+i_j^d)\le n$. These indices satisfy $i_c\le K(n)<p^{J+1}$, so
$p^{jsJ}\prod_cp^{-sv_p(i_c)}$ is an integer; this gives the integrality and $v_p(R_n(C^j))\ge-jsJ$. For
$j<r$ the term $p^{rsJ}b^{-n}R_n(C^j)$ therefore has valuation at least $sJ\ge L$. For $j=r$ write
$z_c=i_c/p^J\in\Gamma$, of loss $J-v_p(i_c)$; then $t\sum i_c^d\le n$ if and only if $S(z)\le u$. A term
is nonzero modulo $p^L$ only if $s\lambda(z)<L$, so every loss is below $L/s\le J$, and conversely every
tuple of level below $L/s$ arises from integers $i_c=p^Jz_c$. For these,
$v_p(b^{ti_c^d}-1)=v_0+d(J-\lambda(z_c))\ge d(J-L/s)\ge L$.
\end{proof}

\begin{lemma}[Reduction to the mass]\label{lem:reductionr}\evidence{Lean-supported paper proof}
Let $r\ge1$, $d\ge r+1$ and let $P\in\Q[x]$ have degree $r$. If $\{u\in[1,p^d):M_r(u)=0\}$ is
Lebesgue-null, then $P(C)$ is normal to base $b$.
\end{lemma}

\begin{proof}
Reduce to integer coefficients as in the proof of \Cref{thm:quadratic} and verify the hypotheses of
\Cref{prop:criterion} for $Y=P_0(C)$ with $c=p$. The insertions of $Y$ up to $N$ lie in the sets
$\{t(i_1^d+\cdots+i_j^d)\le N\}$, $j\le r$, which have $O(N^{r/d})=O(N^{1-1/d})$ elements; a coefficient at
$b^{-m}$ is at most a constant times $(m+1)^r$; and by \Cref{lem:freezer} the reduced denominator $q_n$
divides $p^{rsJ}\le(n+1)^{rs}$. Fix $M_1$ and $L$. For $J\ge J_0(L+v_p(A))$, if
$M_r(u(n))\not\equiv0\pmod{p^L}$ then $v_p(p^{rsJ}b^{-n}R_n(Y))\le L+v_p(A)-1$, so $q_n\ge p^{rsJ-L-v_p(A)+1}$,
which exceeds $M_1$ for large $J$. Hence, apart from finitely many $n$, $q_n<M_1$ forces
$u(n)\in E_L=\{M_r\equiv0\bmod p^L\}$, and the count in the proof of \Cref{thm:quadratic} bounds the upper
density of such $n$ by $p^d|E_L|/(p^d-1)$. The sets $E_L$ decrease to $\{M_r=0\}$. The Lean-supported
step is \Cref{lem:short}, through \Cref{prop:criterion}.
\end{proof}

\subsection{Anchors, centres and isolation}\label{sec:anchors}

In this and the next two subsections $r\ge2$ and $d\ge2$ are arbitrary. Put $\theta=(2r)^{-1/d}$,
$h_{\min}=\lceil\log_p(6/\theta)\rceil$ and $e_r=v_p(r!)$, and fix parameters $\kappa\in(0,1/(2r)]$ and
$\lambda=\kappa6^{-r}$, chosen in \Cref{lem:avgiso} in terms of $p,r,d,s$ only.

Let $H$ be large. For $r\ge3$ call a vector $(h_1,\ldots,h_{r-1})$ of positive integers \emph{dominant} if
\[
h_1\ge h_{\min},\qquad h_2\ge\max(2h_1+1,\lambda H),\qquad h_j\ge2(h_1+\cdots+h_{j-1})+1\ (3\le j\le r-1),
\]
and $k=h_1+\cdots+h_{r-1}\le\kappa H$; for $r=2$ a dominant vector is a single $h_1$ with
$h_{\min}\le h_1\le\kappa H$. The box $h_j\in[\beta_j,2\beta_j]$ with $\beta_j=\kappa H6^{j-r}/4$ consists of dominant
vectors for large $H$, so there are $\gg_r(\kappa H)^{r-1}$ of them. An \emph{anchor} is a tuple
$y=(y_1,\ldots,y_{r-1})$ with $y_j=a_j/p^{h_j}$, $p\nmid a_j$, $0<y_j\le\theta$, for a dominant vector. A
\emph{centre of level $H$} is a tuple $(y,z)$ with $y$ an anchor, $z=m/p^v$, $p\nmid m$, $v=H-k$, and value
$\xi=S(y,z)\in[1,p^d)$. Since $B_y=\sum y_j^d\le(r-1)\theta^d<1/2$, we have $z^d>1/2$ and $z<p$. The
losses $h_1<\cdots<h_{r-1}<v$ of a centre are distinct, so its $r!$ permutations are distinct ordered
tuples of level $H$.

Put $\pi_H=sH+e_r+1$. A tuple of level $\ell$ contributes $p^{s\ell}\equiv0\pmod{p^{\pi_H}}$ unless
$\ell\le H+\lfloor e_r/s\rfloor$; call the tuples with $\ell\le H+\lfloor e_r/s\rfloor$ \emph{relevant}. Put $\delta_H=\epsilon/(H^{r-1}p^H)$
with $\epsilon\in(0,1]$. A centre is \emph{isolated} if every relevant tuple $z'$ with
$|S(z')-\xi|<\delta_H$ is a permutation of the centre.

\begin{lemma}[Jumps at isolated centres]\label{lem:rjump}\evidence{Lean-supported paper proof}
If $\xi$ is the value of an isolated centre, then $M_r\bmod p^{\pi_H}$ is constant on $(\xi-\delta_H,\xi)$
and on $[\xi,\xi+\delta_H)$, and the two constants differ by $r!\,p^{sH}\not\equiv0\pmod{p^{\pi_H}}$. The
values of distinct isolated centres are at least $\delta_H$ apart.
\end{lemma}

\begin{proof}
Modulo $p^{\pi_H}$ only relevant tuples count, and by isolation the only ones with value in
$(\xi-\delta_H,\xi+\delta_H)$ are the $r!$ permutations of the centre, each of weight $p^{sH}$; and
$v_p(r!p^{sH})=sH+e_r<\pi_H$. A distinct centre is a relevant tuple that is not a permutation. The exact
jump of the truncated mass at a centre with distinct coordinates is the Lean declaration
\leanR{PowerSumPeeling}{483}{\ttt{truncMass\_jump\_of\_isolated}}, with
\leanR{PowerSumPeeling}{446}{\ttt{truncMass\_sub}}.
\end{proof}

\begin{lemma}[$q$-th derivative test {\cite[Theorem~5.13]{Titchmarsh86}}]\label{lem:qtest}\evidence{External input}
Let $q\ge2$, $K=2^{q-1}$, and let $g$ be real with $q$ continuous derivatives on an interval $I$ of length
at least $1$, with $\Lambda\le|g^{(q)}|\le2\Lambda$ on $I$. Then
$\sum_{m\in I\cap\Z}e(g(m))\ll_q|I|\Lambda^{1/(2K-2)}+|I|^{1-2/K}\Lambda^{-1/(2K-2)}$.
\end{lemma}

\begin{lemma}[Derivatives of the endpoint]\label{lem:fB}\evidence{Lean-supported paper proof}
For real $B$ and $x>0$ with $x^d+B>0$ let $f_B(x)=(x^d+B)^{1/d}$. For $q\ge2$,
\[
f_B^{(q)}(x)=(d-1)B\,x^{d-q}(x^d+B)^{1/d-q}\,\Phi_q(x^d,B),
\]
where $\Phi_2=1$ and $\Phi_{q+1}=[(d-q)(X+B)+(1-qd)X]\Phi_q+dX(X+B)\partial_X\Phi_q$. The polynomial
$\Phi_q$ is homogeneous of degree $q-2$ in $(X,B)$, and as a polynomial in $X$ it has degree $q-2$ and
leading coefficient $(-1)^q(d+1)(d+2)\cdots(d+q-2)\ne0$.
\end{lemma}

\begin{proof}
For $q=2$, $f_B''=(d-1)x^{d-2}(x^d+B)^{1/d-2}((x^d+B)-x^d)$. Differentiating the formula for $q$ and
factoring out $(d-1)Bx^{d-q-1}(x^d+B)^{1/d-q-1}$ leaves
$(d-q)(x^d+B)\Phi_q+(1-qd)x^d\Phi_q+dx^d(x^d+B)\partial_X\Phi_q$ at $X=x^d$, which is the recurrence;
homogeneity follows by induction. The degree and the leading coefficient are the Lean declarations
\leanR{PowerSumPeeling}{232}{\ttt{phiPoly\_natDegree}},
\leanR{PowerSumPeeling}{242}{\ttt{phiPoly\_leadingCoeff}} and
\leanR{PowerSumPeeling}{222}{\ttt{phiPoly\_ne\_zero}}, in $(\Z[B])[X]$.
\end{proof}

\begin{lemma}[Exponent margins]\label{lem:margins}\evidence{Paper proof}
For $q\ge3$ put $K_q=2^{q-1}$ and, for $0\le\alpha,\beta\le1$,
\[
e_1(q)=\beta+\frac{1-q\alpha}{2K_q-2},\qquad
e_2(q)=\beta+\alpha\Bigl(\frac{q}{2K_q-2}-\frac2{K_q}\Bigr)-\frac1{2K_q-2}.
\]
Let $\sigma_2=1/12$ and $\sigma_r=1/(r(2^{r+1}-2))$ for $r\ge3$. For every $(\alpha,\beta)$ with
$\beta/r\le\alpha\le\beta\le1$ there is $q\in\{3,\ldots,r+1\}$, with $q=3$ if $r=2$ and $q\ge4$ if $r\ge3$,
such that $\max(e_1(q),e_2(q))\le1-\sigma_r$.
\end{lemma}

\begin{proof}
For $r=2$ and $q=3$, $\alpha\ge\beta/2$ gives $e_1\le\frac16+\frac34\beta\le\frac{11}{12}$ and
$e_2=\beta-\frac16$. Let $r\ge3$. For $q\ge4$, $e_1$ decreases and $e_2$ increases in $\alpha$, since
$q2^q>4(2^q-2)$, and both increase with slope one in $\beta$; so for $\alpha\ge1/r$ it suffices to take
$\beta=1$. On $[1/3,1]$ take $q=4$: $1-e_1=1/42$ at $\alpha=1/3$ and $1-e_2=1/28$ at $\alpha=1$. For
$5\le q\le r+1$, on $[1/(q-1),1/(q-2)]$ take order $q$: $1-e_1=1/((q-1)(2K_q-2))$ at $\alpha=1/(q-1)$ and
$1-e_2=(2-4/K_q)/((q-2)(2K_q-2))$ at $\alpha=1/(q-2)$, both at least $\sigma_r$. These intervals cover
$[1/r,1]$. If $\alpha<1/r$, then $\beta\le r\alpha$; take $q=r+1$, so that
\[
e_1\le\frac1{2K_{r+1}-2}+\alpha\Bigl(r-\frac{r+1}{2K_{r+1}-2}\Bigr),\qquad
e_2\le\alpha\Bigl(r+\frac{r-3+4/K_{r+1}}{2K_{r+1}-2}\Bigr)-\frac1{2K_{r+1}-2}.
\]
Both coefficients of $\alpha$ are positive, and at $\alpha=1/r$ the margins are $\sigma_r$ and
$(3-4/K_{r+1})/(r(2^{r+1}-2))>\sigma_r$.
\end{proof}

\begin{lemma}[Averaged isolation]\label{lem:avgiso}\evidence{Lean-supported paper proof}
Let $r,d\ge2$. There are $\kappa\in(0,1/(2r)]$, $A\ge1$, $\mu\in(0,1/d)$ and $H_*$, depending only on
$p,r,d,s$, with the following property. Let $H\ge H_*$, $p^{-\mu H/8}\le\epsilon\le1$ and $0<\tau\le1$,
and let $X\subset[1,p^d)$ be a union of $N\ge1$ disjoint intervals. Then at least
$A^{-1}|X|H^{r-1}p^H-ANH^Ap^{\kappa H}$ centres of level $H$ have value in $X$, and at most
\[
A\epsilon\bigl(\tau^{1-d}|X|+\tau\bigr)H^{r-1}p^H+ANH^Ap^{(1-\mu)H}
\]
of them are not isolated.
\end{lemma}

\begin{proof}
Fix $c=\sigma_r/2$, then $\kappa\le\min(1/(2r),c/(2d))$ so small that the perturbations below keep every
exponent at most $(1-\sigma_r/2)H+O(1)$, and
$\mu=\frac12\min(\sigma_r/4,\lambda/3,c/d-\kappa,\kappa,1/(2r),1/d)$. Constants implied by $\ll$ and $O$
depend only on $p,r,d,s$.

\emph{Centres.} For an anchor $y$ and a component $I_i$ of $X$, the $z$ with $B_y+z^d\in I_i$ form an
interval of length at least $|I_i|/(dp^{d-1})$, which contains at least $(1-1/p)p^v|I_i|/(dp^{d-1})-2$
admissible $z$. Each coordinate of loss $h\ge h_{\min}$ has at least $(1-1/p)\theta p^h-2\ge(\theta/3)p^h$
admissible numerators, so each dominant vector carries at least $(\theta/3)^{r-1}p^k$ anchors; there are
$\gg(\kappa H)^{r-1}$ dominant vectors and at most $(\kappa H)^{r-1}p^{\kappa H+r}$ anchors in all. This
gives the first bound.

\emph{Rows.} A non-isolated centre has a relevant competitor $z'$ within $\delta_H$, not a permutation.
Write $z'=(w,x)$ with $x=m'/p^{v'}$ a coordinate of largest loss and $w$ the others; its level is
$H'\le H+\lfloor e_r/s\rfloor$, so $v'\ge H'/r$. Put $B=\sum w_i^d-B_y$ and $f_B$ as in \Cref{lem:fB}.
Then $x^d+B\in[1/4,p^d+1]$, $|B|\le rp^d$ and $|f_B(x)-z|\le\delta'=C_d\delta_H$. A \emph{row} is a pair
$(y,w)$; for fixed losses of $y$ and $w$ there are at most $p^{k+H'-v'+2r}$ rows, and summing over all
loss data multiplies counts by $H^{O(r)}$. In a row each $x$ meets at most one admissible $z$, and a
collision forces $\dist{p^vf_B(x)}\le\eta_0=p^v\delta'=C_d\epsilon p^{-k}H^{1-r}$.

\emph{Rows with $B=0$.} Here $|x-z|\le\delta'<p^{-H-\lfloor e_r/s\rfloor}\le p^{-\max(v,v')}$ for large
$H$, so $x=z$, hence $v'=v$ and the coarse loss of $w$ is $k'=k+(H'-H)\le k+\lfloor e_r/s\rfloor$, and
$\sum w_i^d=\sum y_i^d$. If $w=y$ as multisets the competitor is a permutation of the centre. For $r=2$,
$w_1^d=y_1^d$ forces $w=y$. For $r\ge3$ peel coordinates from $j=r-1$ down to $2$. Suppose those above $j$
have been matched; the remaining anchor sum $\sum_{i\le j}y_i^d$ has the unique top loss $h_j$, and the
remaining competitor coordinates have total loss $k'_j\le k_j+\lfloor e_r/s\rfloor$ with
$k_j=h_1+\cdots+h_j\le\frac32h_j-\frac12$, so $k'_j<2h_j$ for large $H$, because $h_j\ge\lambda H$. By the
peeling lemma \leanR{PowerSumPeeling}{372}{\ttt{dominant\_loss\_unique\_top\_coord}} the competitor has a
unique remaining coordinate $c'/p^{h_j}$ of loss $h_j$, all others being of smaller loss. If $c'=a_j$,
remove it and descend. Otherwise $a_j^d-c'^d=M\ne0$, where $M=p^{dh_j}(\sum'w_i^d-\sum_{i<j}y_i^d)$ is
fixed by the lower coordinates, $|M|\le rp^{d(h_j+1)}$, and by
\leanR{PowerSumPeeling}{79}{\ttt{card\_pow\_diff\_eq\_le\_divisors}} there are at most
$2\tau(|M|)\ll p^{h_j/6}$ pairs $(a_j,c')$. The matched coordinates contribute at most $p^{k-k_j+O(1)}$
choices, the lower ones $O(k^rp^{2k_{j-1}})$, and the fine coordinate $p^{v+1}$; the total is
$p^{H-h_j+k_{j-1}+h_j/6+O(1)}k^r\le p^{H-h_j/3+O(1)}k^r\le p^{(1-\lambda/3)H}H^{O(1)}$. If every step matches,
the last remaining coordinates satisfy $w_1^d=y_1^d$ and $w=y$. Summing over loss data, rows with
$B=0$ contribute $O(H^{O(r)}p^{(1-\lambda/3)H})$.

\emph{Rows with $0<|B|<p^{-cH}$.} The denominator of $B$ divides $p^{dg}$ with $g$ the largest coarse
loss, so $g>cH/d>\kappa H\ge k$, and $g$ is attained in $w$. With the other coordinates fixed, that
coordinate lies in a set of diameter at most $(2p^{-cH})^{1/d}$, which contains fewer than $3p^{g-cH/d}$
grid points. Charging every fine point gives $O(H^{O(r)}p^{k+H'-cH/d})=O(H^{O(r)}p^{(1-c/d+\kappa)H})$.

\emph{Removed points.} In rows with $|B|\ge p^{-cH}$ put $\eta=p^{-2\kappa H}$ and $\rho=\eta^{d(q-2)}$, with
$q$ chosen below, and remove the $x$ with $x\le\eta$ or $|\Phi_q(x^d,B)|\le\rho$. By \Cref{lem:fB} the
second set lies within $(\rho/|{\rm lead}|)^{1/(q-2)}=O(\eta^d)$ of the at most $q-2$ roots of
$X\mapsto\Phi_q(X,B)$, hence, since $|x^d-x'^d|\ge|x-x'|^d$, in at most $q-2$ intervals of length
$O(\eta)$. Since $k+H'-v'\le(\kappa+1-1/r)H+O(1)\le(1-\kappa)H+O(1)$, the removed points number
$O(H^{O(r)}(\eta p^{k+H'}+p^{k+H'-v'}))=O(H^{O(r)}p^{(1-\kappa)H})$.

\emph{Exponential sums.} On the retained points, $|f_B^{(q)}|\gg p^{-\gamma H}$ with
$\gamma=c+2\kappa(d(q-1)-q)$ if $q\le d$ and $\gamma=c+2\kappa d(q-2)$ if $q>d$, and $|f_B^{(q)}|\ll1$ if
$q\le d$ and $\ll p^{2\kappa(q-d)H}$ if $q>d$. Drop the condition $p\nmid m'$, so that $m'$ runs over
consecutive integers; this only enlarges the set counted, so it gives an upper bound for the collisions. Since $S(w,x)$ increases in $x$, the
$x$ of a row with $S(w,x)$ in the $\delta_H$-neighbourhood $X^+$ of $X$ form at most $N$ intervals; cutting
also where $f_B^{(q)}$ changes sign, where $|f_B^{(q)}|$ changes monotonicity or dyadic range, gives at most $N+O(qH)$ pieces on
which $|f_B^{(q)}|$ has constant sign and ratio at most $2$ (the equation $|f_B^{(q)}|=D$ becomes, after
raising to a power, a nonzero polynomial equation in $x$ of bounded degree). Choose $q$ for the class of
the competitor by \Cref{lem:margins}, with $\alpha=\min(1,v'/H)$ and $\beta=\min(1,H'/H)$; then $\beta/r\le\alpha\le\beta\le1$, and
$0\le v'-\alpha H\le\lfloor e_r/s\rfloor$ changes the exponents below by $O(1)$. Apply the
Erd\H{o}s--Tur\'an inequality \cite[Chapter~2, Theorem~2.5]{KN} with $L_{\rm ET}=\lceil p^kH^{r+1}/\epsilon\rceil$
frequencies to the points $p^vf_B(m'/p^{v'})$ of each piece, and \Cref{lem:qtest} to the phases
$lp^vf_B(m'/p^{v'})$, whose $q$-th derivative has size $lp^v|f_B^{(q)}|/p^{qv'}$. Summing $1/l$ over
$l\le L_{\rm ET}$, then over pieces (the first term of \Cref{lem:qtest} is linear in the piece length, and
the second grows by at most $(N+O(qH))^{2/K_q}\le NH^{O(1)}$ by concavity) and over rows, the error terms
total
\begin{gather*}
NH^{O(r)}\epsilon^{-1/(2K_q-2)}\bigl(p^{E_1}+p^{E_2}\bigr),\\
E_1=k+H'+\frac{H-qv'}{2K_q-2}+\frac{2\kappa(q-d)^+H}{2K_q-2},\qquad
E_2=k+H'-\frac{2v'}{K_q}+\frac{qv'-H+k+\gamma H}{2K_q-2},
\end{gather*}
together with at most $2(N+O(qH))$ points per row from pieces with fewer than two points. By
\Cref{lem:margins} and the choice of $c,\kappa$, $E_1,E_2\le(1-\sigma_r/2)H+O(1)$; since
$\epsilon^{-1}\le p^{\mu H/8}$ and $\mu\le\sigma_r/8$, these errors are $O(NH^{O(r)}p^{(1-\mu)H})$.

\emph{Main term.} The remaining contribution is
$\sum_y(2\eta_0+6/L_{\rm ET})\,\mathrm{Comp}(X^+)$, where $\mathrm{Comp}(X^+)$ counts relevant competitor
tuples with value in $X^+$ (with $p\nmid m'$ dropped); it does not depend on $y$. Since there are at most
$p^{k+r}$ anchors per dominant vector and at most $H^{r-1}$ vectors,
$\sum_y2\eta_0\le2C_dp^r\epsilon$ and $\sum_y6/L_{\rm ET}\le6p^r\epsilon H^{-2}$. Put
$\tilde\tau=\max(\tau,p^{-\mu H})$ and $\mu_0=1/(2r)$, and call a competitor \emph{regular} if some
coordinate $x_m\ge\tilde\tau$ has loss $g\ge\mu_0H$. Fixing the other coordinates (at most $p^{H'-g+r}$
choices) and using $\partial S/\partial x_m\ge d\tilde\tau^{d-1}$, at most $p^g|X^+|/(d\tilde\tau^{d-1})+N$
values of $x_m$ give $S\in X^+$; summing over loss vectors, regular competitors number
$O(\tilde\tau^{1-d}|X^+|H^{r-1}p^H)+O(NH^{r-1}p^{(1-\mu_0)H})$. Here $|X^+|\le|X|+2N\delta_H$ and
$\epsilon\cdot2N\delta_H\tilde\tau^{1-d}H^{r-1}p^H\le2Np^{(d-1)\mu H}\le2Np^{(1-\mu)H}$, because $\mu<1/d$. If an
irregular competitor has $H'\ge H/2$, some coordinate has loss at least $\mu_0H$, and all such
coordinates lie below $\tilde\tau$; so each loss vector carries at most $2\tilde\tau p^{H'+r}$ irregular
tuples, $O(\tilde\tau H^{r-1}p^H)$ in all, while loss vectors with $H'<H/2$ carry $O(H^rp^{H/2})$ tuples.
Since $\tilde\tau\le\tau+p^{-\mu H}$ and $\tilde\tau^{1-d}\le\tau^{1-d}$, the main term is at most
$A\epsilon(\tau^{1-d}|X|+\tau)H^{r-1}p^H+ANH^Ap^{(1-\mu)H}$ for a suitable $A$.
\end{proof}

The irregular competitors are bounded without any localization; their total mass is a fraction
$\tilde\tau$ of all competitors wherever $X$ lies. Localizing them is impossible at a fixed scale.

\begin{proposition}[Uniform isolation fails]\label{prop:nouniform}\evidence{Paper proof; finite certificate}
Let $p=3$, $d=4$, $r=3$, $\epsilon=\frac14$ and $\delta_H=\epsilon/(H^23^H)$. Put
$R_H=\lceil(8(32H^2)^4)^{1/3}\rceil+1$ and $\rho_H=1/R_H\asymp H^{-8/3}$. For every $H\ge6$, every point of
$(2+\rho_H,2+2\rho_H]$ lies within $\delta_H$ of the value $1+1+(a/3^H)^4$ of a relevant tuple with
$3\nmid a$, so no centre of level $H$ with value in the nonempty interval
$(2+\rho_H+\delta_H,2+2\rho_H-\delta_H)$ is isolated. For $H=4000$ one has $\rho_H>3^{-\sigma_3H/3}$.
\end{proposition}

\begin{proof}
Consecutive admissible numerators differ by at most $2$, and the gaps between the points $(a/3^H)^4$
increase with $a$. So all gaps up to the first point above $2\rho_H$ are at most
$(a_{\rm hi}^4-(a_{\rm hi}-2)^4)/3^{4H}<8a_{\rm hi}^3/3^{4H}$, where $a_{\rm hi}\le(2\rho_H)^{1/4}3^H+3$. The choice
of $R_H$ makes $8a_{\rm hi}^3\le2\delta_H3^{4H}$ for $H\ge5$. Since $\rho_H^{1/4}3^H\ge1$, the net starts below
$2+\rho_H$, and $\rho_H>2\delta_H$ for $H\ge6$, so the interval in the statement is nonempty (for $H=4,5$ the net
property still holds, but that interval is empty). A centre has three distinct positive losses,
so it is never one of the tuples $(1,1,a/3^H)$, which have level $H$. Finally $R_{4000}^{126}<3^{4000}$ is
checked exactly, and $\sigma_3/3=1/126$. The finite certificate (\Cref{sec:formal}) checks the net exactly for
$H\in\{6,10,100,1000,4000\}$, as an independent confirmation of the analytic step.
\end{proof}

The same mechanism works in general: next to a sum $W$ of $r-1$ coordinates of coarse loss $g$,
competitors with a small last coordinate accumulate with density of order $p^{H-g}\rho^{1/d-1}$ at
distance $\rho$, so isolation fails on intervals of length of order $(p^{-g}H^{1-r})^{d/(d-1)}$ just
above $W$. This is why the following estimate carries the exponent $d/(d-1)$.

\subsection{An effective small-ball estimate}\label{sec:smallball}

\begin{theorem}[Small balls]\label{thm:smallball}\evidence{Lean-supported paper proof}
Let $r,d\ge2$. There is $C_*$, depending on $p,r,d,s$, such that for every $w\in\Z_p$ and every $L\ge3$,
\[
\bigl|\{u\in[1,p^d):\ M_r(u)\equiv w\pmod{p^L}\}\bigr|\le C_*(\log L)^{-d/(d-1)}.
\]
In particular every level set of $M_r$ is Lebesgue-null.
\end{theorem}

\begin{proof}
Let $E=\{M_r\equiv w\bmod p^L\}$, take $\kappa,A,\mu,H_*$ from \Cref{lem:avgiso}, enlarge $A$ to
$A\ge\max(p^d,r-1)$, and put $\Gamma_0=\lceil4/\mu\rceil+1$. Choose levels
$H_0=H_*+\lceil C'\log\log L\rceil$ with $C'\ge24/(\mu\log p)$, and $H_{j+1}=\Gamma_0H_j$. We construct unions of
intervals $[1,p^d)=X_0\supset X_1\supset\cdots\supset E$ with $N_j\le AH_j^Ap^{H_j}(j+1)$ components, using
only levels with $\pi_{H_j}\le L$; this allows $m(L)\gg\log L$ steps.

Let $x=|X_j|$ and $H=H_{j+1}$. If $x\le(\log L)^{-d/(d-1)}$, stop. Otherwise apply \Cref{lem:avgiso} with
$\tau=(x/p^d)^{1/d}$ and $\epsilon=x^{(d-1)/d}/(8p^dA^2)$; for large $L$ both $\epsilon$ and $x$ exceed
$p^{-\mu H/8}$, because $H\ge\Gamma_0C'\log\log L$. Since $\tau^{1-d}x+\tau\le2p^{d-1}x^{1/d}$, the main
collision term is at most $xH^{r-1}p^H/(4A)$, and since $N_j\le AH_j^Ap^{H_j}(j+1)$ with $H_j\le\mu H/4$, every
error term is at most $p^{(1-\mu/2)H}\le xH^{r-1}p^H/(8A)$, since $x\ge p^{-\mu H/8}$ and $H$ is large. There are two
error terms, one in each bound of \Cref{lem:avgiso}, so at least $(1-\frac18-\frac14-\frac18)A^{-1}xH^{r-1}p^H=
xH^{r-1}p^H/(2A)$ centres with value in $X_j$ are isolated. By \Cref{lem:rjump} each has two flanks $(\xi-\delta_H/2,\xi)$ and
$(\xi,\xi+\delta_H/2)$ on which $M_r\bmod p^{\pi_H}$ takes different constant values, so one flank misses
$\{M_r\equiv w\bmod p^{\pi_H}\}\supset E$. These flanks are disjoint, and at most $2N_j$ centres lie within
$\delta_H/2$ of the boundary of $X_j$. Removing the missed flanks of the others gives $X_{j+1}\supset E$ with
$N_{j+1}\le N_j+pH^{r-1}p^H$ components and
\[
|X_{j+1}|\le x-\Bigl(\frac{xH^{r-1}p^H}{2A}-2N_j\Bigr)\frac{\delta_H}2\le x\bigl(1-c_1x^{(d-1)/d}\bigr),\qquad
c_1=\frac1{40p^dA^3}.
\]
Put $a=(d-1)/d$ and $y_j=x_j^{-a}$, where $x_j=|X_j|$. The Lean-supported steps are those of
\Cref{lem:rjump,lem:avgiso}. Bernoulli's inequality gives
$y_{j+1}\ge y_j(1+ac_1x_j^a)=y_j+ac_1$, so $x_j\le(ac_1j)^{-d/(d-1)}$ while the process runs. Taking
$j=m(L)\gg\log L$ proves the bound.
\end{proof}

The kill fraction per level is of order $|X_j|^{(d-1)/d}$ rather than a constant; this is where the
exponent $d/(d-1)$, rather than a power of $L$, comes from.

\subsection{Normality in every degree below \texorpdfstring{$d$}{d}}\label{sec:higherdegree}

\begin{theorem}[One order]\label{thm:higher}\evidence{Lean-supported paper proof}
Let $s,r\ge1$, $d\ge r+1$ and $C=C(b,p,s,d)$. Then $P(C)$ is normal to base $b$ for every nonconstant
$P\in\Q[x]$ of degree at most $r$. Consequently $C$ is not algebraic of degree less than $d$, and the
radix shifts of $(C,C^2,\ldots,C^{d-1})$ are uniformly distributed in $\T^{d-1}$.
\end{theorem}

\begin{proof}
For $\deg P=1$ apply \Cref{lem:reductionr}: $M_1\not\equiv0\pmod p$, so $\{M_1=0\}$ is empty. For $2\le\deg P\le r$, \Cref{thm:smallball} gives
$|\{M_{\deg P}=0\}|=0$, and \Cref{lem:reductionr} applies since $d\ge\deg P+1$. An algebraic $C$ of degree
less than $d$ would be a zero of a nonconstant rational polynomial of degree at most $d-1$, and $0$ is
not normal. Joint distribution follows from Weyl's criterion \cite[Chapter~1, Theorem~6.2]{KN}, since every
nonzero integer combination of $C,\ldots,C^{d-1}$ is a nonconstant polynomial of degree at most $d-1$.
The Lean-supported steps are those of \Cref{lem:rjump,lem:fB,lem:avgiso} and \Cref{lem:short}.
\end{proof}

For several orders the top-depth mass is a combination of masses with different weights, and the
centre must be chosen from the smallest order.

\begin{theorem}[Several orders]\label{thm:several}\evidence{Lean-supported paper proof}
Let $1\le s_1<\cdots<s_m$ be integers, $X_a=C(b,p,s_a,d)$, and let $P\in\Q[X_1,\ldots,X_m]$ be nonconstant of
total degree $r\le d-1$. Then $P(X_1,\ldots,X_m)$ is normal to base $b$. In particular the $X_a$ satisfy no
nonzero rational polynomial relation of total degree less than $d$, and distinct nonconstant monomials
of total degree less than $d$ have jointly uniformly distributed radix shifts.
\end{theorem}

\begin{proof}
\emph{Reduction.} Clear denominators, so that $P=\sum_\nu A_\nu X^\nu$ with integer $A_\nu$. A monomial
$X^\nu$ expands over \emph{slot tuples}: for each $a$ an ordered $\nu_a$-tuple of coordinates, and a slot
carrying variable $a$ and a coordinate of loss $h$ has weight $p^{s_ah}$. Put
$S^*=\max\{s\cdot\nu:A_\nu\ne0,\ \nu\ne0\}$ and $T=\{\nu:A_\nu\ne0,\ s\cdot\nu=S^*\}$. As in
\Cref{lem:freezer}, for
\begin{equation}\label{eq:freezeJ}
J\ge J_1(L)=\max\bigl(L,\ \lceil L/s_1\rceil+\lceil L/d\rceil\bigr),
\end{equation}
\[
p^{S^*J}b^{-n}R_n(P)\equiv M_P(u)=\sum_{\nu\in T}A_\nu M_\nu(u)\pmod{p^L},
\]
where $M_\nu(u)$ sums the slot weights over slot tuples with $S\le u$, and monomials outside $T$ carry a
factor $p^{(S^*-s\cdot\nu)J}$ with $S^*-s\cdot\nu\ge1$, which vanishes modulo $p^L$ because $J\ge L$; a
slot tuple of a monomial in $T$ is nonzero modulo $p^L$ only if every loss $h_i$ satisfies $h_i<L/s_1$, and
then $v_p(b^{ti^d}-1)=v_0+d(J-h_i)\ge L$ by \eqref{eq:freezeJ}. The condition $J\ge L$ cannot be dropped:
with only $J\ge\lceil L/s_1+L/d\rceil$, a weight gap of one survives (for example $L=2$, $J=1$, $u=1$ gives
the residue $1+p$ instead of $1$ in the fixture of \Cref{sec:formal}); this is harmless for normality,
where $L$ is fixed and $J\to\infty$, but matters for the rate in \Cref{thm:severalrate}. The insertions number $O(N^{r/d})$, so as in \Cref{lem:reductionr} it suffices that
$\{M_P=0\}$ be null. Only slot tuples of weighted level below $L$ matter modulo $p^L$.

\emph{The target monomial.} For $\nu\in T$ list its slot orders in nondecreasing order as
$\mathbf o(\nu)$; distinct $\nu$ give distinct lists, none a proper prefix of another (their entries sum
to $S^*$). Let $\nu^*$ have the lexicographically least list, $s_*=s_{a_*}$ its first entry,
$n_*=\nu^*_{a_*}$, $r_*=\deg\nu^*$, $m_*=r_*-n_*$, and $u_1\le\cdots\le u_{m_*}$ its remaining orders, all
larger than $s_*$. Lexicographic minimality minimizes $s_*$ and maximizes $\nu_{a_*}$ over $T$. If
$r_*=1$ then $T=\{X_{a_*}\}$ and $M_P\equiv A_{\nu^*}\lfloor u^{1/d}\rfloor\not\equiv0\pmod{p^{v_p(A_{\nu^*})+1}}$, so
$\{M_P=0\}$ is empty. Assume $r_*\ge2$.

\emph{The centre.} Fix $G_0>\max_{\nu\in T}v_p(A_\nu\prod_a\nu_a!)+2$ and put
$e_{\rm prec}=v_p(A_{\nu^*}\prod_a\nu^*_a!)+1$. A centre consists of a \emph{bounded part}
$y_j=a_j/p^{g_j}\le\theta$, $1\le j\le m_*$, a \emph{free part} $b_j/p^{h_j}\le\theta$, $1\le j<n_*$, and a last
coordinate $z=m/p^v$, all numerators prime to $p$. The bounded part does not depend on $H$; it is chosen
sequentially: with $B_{j-1}=\frac{s_m}{s_*}\sum_{i<j}g_i+e_{\rm prec}$, take $g_j$ larger than $2B_{j-1}$,
than $s_m\sum_{i<j}g_i+G_0$ and than $h_{\min}$, and so large that the number of tuples $(w_l)$ of total loss at
most $B_{j-1}$, times $2\max_{0<|M|\le rp^{d(g_j+1)}}\tau(|M|)$, is less than the number of admissible
$a_j\le\theta p^{g_j}$; then choose $a_j$ outside all solutions of $c^d-a_j^d=M\ne0$ with
$M=p^{dg_j}(\sum_{i<j}y_i^d-\sum_lw_l^d)$. Put $M_s=\max(2,s_m)$, choose $\kappa<\min(1/(2r),1/(M_s+1))$ and
$\lambda\le\kappa/(3M_s)^{n_*-1}$, and require $h_1\ge\max(\lambda H,M_s\sum g_i+G_0)$,
$h_{j+1}\ge M_s(\sum g_i+h_1+\cdots+h_j)+G_0$, $k=\sum g_i+\sum h_j\le\kappa H$ and $v=H-k$. All losses are
distinct, and $v-s_mk\ge G_0$ for large $H$.

\emph{The jump.} By the rearrangement inequality the least weight of an assignment of the centre
coordinates to the slots of $\nu^*$ puts the ascending orders on the descending losses:
$w^*=s_*(v+h_1+\cdots+h_{n_*-1})+\sum_iu_ig_{m_*+1-i}$, attained by exactly $\prod_a\nu^*_a!$ slot tuples;
every other assignment has weight at least $w^*+G_0$ by the growth conditions on the losses. At precision
$\pi=w^*+e_{\rm prec}$ the jump at an isolated centre is $A_{\nu^*}\prod_a\nu^*_a!\,p^{w^*}\not\equiv0$.

\emph{Competitors.} A slot tuple of $\nu'\in T$ of weighted level below $\pi$ has unweighted level
$H'<H+C_g$ with $C_g=O(1)$. Such tuples number $o(p^H)$ if $\nu'_{a_*}=0$, $O(H^{\nu'_{a_*}-1}p^H)$ if
$\nu'_{a_*}\ge1$, and $o(H^{n_*-1}p^H)$ unless $\nu'_{a_*}=n_*$; centres in an interval $I$ number
$\asymp(\kappa H)^{n_*-1}p^{H-\sum g_i}|I|$. Put $\delta_H=\epsilon/(H^{n_*-1}p^H)$ with
$\epsilon\le\epsilon_0=(2C_d)^{-1}p^{-C_g}$.

\emph{Exact coincidences.} As in \Cref{lem:avgiso}, $\epsilon\le\epsilon_0$ forces $x=z$ and $v'=v$; the
coordinate $x$ sits in an order-$s_*$ slot, since otherwise its weight alone exceeds $\pi$. Cancelling $s_*v$,
the remaining weights obey a budget $(*)$: total weight at most
$s_*(h_1+\cdots+h_{n_*-1})+\sum_iu_ig_{m_*+1-i}+e_{\rm prec}$. Peel the free coordinates downwards: at
$h_j\ge\lambda H$ the budget bounds the remaining unweighted loss by less than $2h_j$, so by the peeling lemma
there is a unique competitor coordinate of loss $h_j$; if it differs from the anchor coordinate the divisor
bound saves $p^{-h_j/3}$, and if it matches it must occupy an order-$s_*$ slot. After the free part, $\nu'$
has used $n_*$ order-$s_*$ slots, so $\nu'_{a_*}=n_*$, and its remaining orders satisfy
$\sum_lu_l(\nu')e_l\le\sum_iu_ig_{m_*+1-i}+e_{\rm prec}$. Now peel the bounded coordinates for
$j=m_*,\ldots,1$: inductively the lower orders agree, lexicographic minimality gives
$u_{m_*+1-j}(\nu')\ge u_{m_*+1-j}$, the budget forces a unique competitor coordinate $c/p^{g_j}$ in a slot of
order exactly $u_{m_*+1-j}$, the remaining coordinates have total loss at most $B_{j-1}$, and the choice of
$a_j$ forces $c=a_j$. At the end $\nu'=\nu^*$ and the competitor is one of the slot tuples of the jump.

\emph{Other collisions.} Degree-one competitors $\nu'\in T$ have order above $s_*$ and number
$O(p^{s_*H/o_1})=o(p^H)$; they are discarded. Competitors of degree $2\le j\le r$ satisfy $\alpha\ge\beta/r$,
and rows with small $|B|$, removed points and exponential sums are treated exactly as in
\Cref{lem:avgiso}, with $q\le r+1\le d$, $L_{\rm ET}=\lceil p^kH^{n_*+1}\rceil$ and an error
$o(H^{n_*-1}p^H)$.

\emph{Main term.} Fixed higher-slot coordinates $w_{\rm hi}$ with $u_0=S(w_{\rm hi})$ leave $n_*<d$
order-$s_*$ slots and produce integrable singularities $(u-u_0)_+^{n_*/d-1}$ in the limiting competitor
density, so isolation cannot hold uniformly in $I$. Discard competitors with $\nu'_{a_*}<n_*$, and split the
others at a weighted higher-slot level $E_0(|I|)$: those above it number
$O(H^{n_*-1}p^Hp^{-\eta_1E_0})$ by convergence of $\prod_j\sum_{e\ge0}p^{e(1-o_j/s_*)}$, with $\eta_1>0$; for
the finitely many others, the lattice count
$p^{H'}\mathrm{Vol}(R_u)-(p+1)^{n_*-1}\sum_jp^{H'-h'_j}\le\#(R_u\cap\Lambda)\le p^{H'}\mathrm{Vol}(R_u)$ for the grid
$\Lambda=\prod_jp^{-h'_j}\Z_{>0}$ and $R_u=\{x\in(0,p)^{n_*}:\sum x_j^d\le u\}$ (pack interior-disjoint lower and
upper boxes), applied on $I^+-S(w_{\rm hi})$, has error $o(H^{n_*-1}p^H)$. Hence for every interval $I$,
\[
\limsup_{H\to\infty}\frac{\mathrm{Comp}(I^+)}{H^{n_*-1}p^H}\le C_P\int_If,\qquad
f(u)=1+\sum_{\nu',w_{\rm hi}}p^{-\frac1{s_*}\sum_jo_je_j}\,\frac{d}{du}\mathrm{Vol}_{n_*}\Bigl(\sum_{i\le n_*}x_i^d\le
(u-S(w_{\rm hi}))_+\Bigr),
\]
and $f\in L^1$ because $n_*/d>0$. For $M\ge1$ take $\epsilon_M=\min(\epsilon_0,c_1/M)$ with $c_1$ small; on
every interval with $|I|^{-1}\int_If\le M$, isolation of half the centres holds for $H\ge H_0(I,M)$, and the
reverse Fatou lemma gives $|I\cap\limsup_HG_H(\epsilon_M)|\ge c_2\epsilon_M|I|$, where $G_H$ is the union of the
intervals of radius $\delta_H/2$ about isolated centres. As in the proof of \Cref{prop:nulllevel}, no point of
$\limsup_HG_H$ is a density point of a level set $E=\{M_P=w\}$, so $|I\cap E|\le(1-c_2\epsilon_M)|I|$ on such
intervals. If $|E|>0$, a point of $E$ that is both a density point of $E$ and a Lebesgue point of $f$, with
$M=\lceil f(x_0)\rceil+1$, gives a contradiction. So every level set of $M_P$ is null.

The consequences follow as in \Cref{thm:higher}, with Weyl's criterion in several dimensions. The
Lean-supported steps are the peeling lemma, the divisor bound and \Cref{lem:fB}.
\end{proof}

\subsection{Rates}\label{sec:rates}

\begin{theorem}[Effective rates]\label{thm:effective}\evidence{Paper proof}
Let $s,r\ge1$, $d\ge r+1$, $C=C(b,p,s,d)$, and let $P\in\Q[x]$ have degree $r$. Then
\[
D_N(P(C))\ll_P(\log\log N)^{-d/(d-1)}\quad(r\ge2),\qquad D_N(P(C))\ll_PN^{-\eta}\quad(r=1),
\]
with $\eta=\eta(p,s,d)>0$.
\end{theorem}

\begin{proof}
Clear denominators: by the shift used in the proof of \Cref{thm:quadratic}, $b^eP(C)=P_0(C)/Q_1$ with
$P_0\in\Z[x]$ of leading coefficient $A$ and $\gcd(Q_1,b)=1$, and $D_N(P(C))\le D_{N-e}(b^eP(C))+e/N$. The
shifted prefixes then have reduced denominators $q_n\mid p^{rsJ+O_P(1)}M'$ with $M'$ fixed and prime to $pb$.
Take $L=\lfloor c_0\log N\rfloor$ with $c_0\log p\le rs/(8d)$, $(c_0\log p)/s<r/d$ and
$c_0(1+d/s)\log p<1/2$; the last condition puts every shift of the shells $J<J_0(L)$ below $\sqrt N$. Fix
$\epsilon_1=1/(6d)$ and $\eta_0=\min(\frac12,\frac{rs}{8d})$.

\emph{Runs.} Cut $[1,N]$ at the insertion positions, at the shell boundaries $tp^{dJ}$, at the
$O(\log N\cdot L^{r-1}p^{L/s})=O(N^{r/d})$ cell boundaries of $M_r\bmod p^L$ and at $\lceil\sqrt N\rceil$; this
gives $O(N^{r/d})$ maximal half-open runs $[n_0,n_0+T_0)$, each lying entirely below $\lceil\sqrt N\rceil$ or
entirely at or above it. On a run, $R_{n_0+j}\equiv b^jR_{n_0}\pmod1$, so $q_n$ and the residue of $M_r(u)$ are
constant. Put $B=\lceil2\log_bN\rceil$; the \emph{trimmed run} is $[n_0,n_0+T)$ with $T=\max(0,T_0-B)$. Discard:
the runs on which $M_r(u)\equiv0\pmod{p^{L-v_p(A)}}$ ($O(N(\log L)^{-d/(d-1)})$ shifts, by \Cref{thm:smallball},
the sampling of shells and the cell count); the runs with $T_0<N^{1-r/d-\epsilon_1}$ ($O(N^{1-\epsilon_1})$
shifts); the runs below $\lceil\sqrt N\rceil$; and the final $B$ positions of every remaining run
($O(N^{r/d}\log N)$ shifts). The coefficients of $P_0(C)$ at $b^{-m}$ are $O(m^{(r-1)/d})$, so at trimmed
positions $|b^nb^eP(C)-R_n|\le N^{-1}$; this pointwise bound is not used at the discarded final positions. On a
retained run $T_0\ge2B$ for large $N$, so $T\ge T_0/2$, and the slack below absorbs this factor.

\emph{A retained run.} By \Cref{lem:freezer}, $v_p(q_{n_0})\ge rsJ-L+1-O(1)$ with
$rsJ\ge\frac{rs}{2d}\log_pN-O(1)$. Apply the Erd\H{o}s--Tur\'an inequality with $H_{\rm ET}=\lceil N^{\eta_0}\rceil$
frequencies to the model points $\fr{b^jR_{n_0}}$, $j<T$, with $T$ the trimmed length. For $h\le H_{\rm ET}$ the reduced form of $hR_{n_0}$
is $a/(p^eQ')$ with $Q'\mid M'$ and $e=v_p(q_{n_0})-v_p(h)\ge\frac{rs}{4d}\log_pN-O(1)$, which exceeds
$A(G)\le a_0+v_p(\ord_{M'}b)$ for $G=b^\tau$, $\tau=\lcm(t,\ord_{Q'}b)=O(1)$. Also
$e\le\frac{rs}d\log_pN+O(1)$ and $M_0\ge T/\tau\gg N^{1-r/d-\epsilon_1}$, so $p^{\theta e}\le M_0$ with
\[
\theta=\min\Bigl(1,\frac{d-r-1/2}{rs}\Bigr)\ge\frac1{2rs};
\]
this is where $d\ge r+1$ enters. \Cref{cor:split,lem:uniformorbit} give
$|\sum_{j<T}e(hb^jR_{n_0})|\ll T^{1-\delta}$ uniformly in $h\le H_{\rm ET}$, so the run has discrepancy
$\ll N^{-\eta_0}+N^{-\delta/(2d)}\log N$, using $T\ge N^{1/(2d)}$.

\emph{Summation.} \Cref{lem:blocks} with the retained runs as blocks, the discarded shifts as uncovered,
and $\epsilon=N^{-1}$ gives the bound for $r\ge2$, the discarded failed runs dominating. For $r=1$ the frozen
residue is $A\lfloor u^{1/d}\rfloor\not\equiv0\pmod{p^{v_p(A)+1}}$, so no run fails, and every term is a
negative power of $N$.
\end{proof}

The exponents of the power savings are poor: $\delta=2^{-\lceil2/\theta\rceil-3}$, for example
$2^{-19}$ for $(r,s,d)=(2,2,3)$. They affect only the terms other than the failed runs. The rate for
$r\ge2$ comes from \Cref{thm:smallball} alone; a power of $\log N$ would follow from a local form of the
competitor count in \Cref{lem:avgiso} near sums of fewer coordinates, which we have not proved.

\begin{corollary}\label{cor:flagship2}\evidence{Paper proof}
$D_N(\xi^2;2)\ll(\log\log N)^{-3/2}$. The numbers $C(2,3,1,4)$, $C(2,3,1,4)^2$ and $C(2,3,1,4)^3$ are jointly
normal to base $2$, and $D_N(C(2,3,1,4)^3;2)\ll(\log\log N)^{-4/3}$.
\end{corollary}

\begin{proof}
Apply \Cref{thm:effective} with $(b,p,s,d,r)=(2,3,1,3,2)$ and $(2,3,1,4,3)$, and \Cref{thm:higher}.
\end{proof}

\paragraph{Several orders.} We keep the notation of the proof of \Cref{thm:several}: $P=\sum_\nu A_\nu X^\nu$
has integer coefficients and total degree $r\le d-1$; $S^*$ and $T$ are the largest weighted degree and the
set of monomials attaining it; $\nu^*$ is the monomial of $T$ with the lexicographically least list of
slot orders, with smallest order $s_*$ of multiplicity $n_*$, degree $r_*$ and remaining orders
$u_1\le\cdots\le u_{m_*}$; $M_P=\sum_{\nu\in T}A_\nu M_\nu$ is the weighted mass. Every list of a monomial in
$T$ has entries at least $s_*$ and at most $n_*$ entries equal to $s_*$, and no list is a proper prefix of
another. If $r_*=1$, then $T=\{X_{a_*}\}$ and $M_P\equiv A_{\nu^*}\lfloor u^{1/d}\rfloor$; this is the case of a
\emph{single linear maximizer}. Assume now $r_*\ge2$, and fix the bounded part $(y_j)_{j\le m_*}$ of the centres
and the constants $G_0$ and $e_{\rm prec}$ as in that proof.

For the rate the free losses range over boxes of length comparable to $H$. Put $\beta_j=\lambda H(24s_m)^{j-1}$
for $1\le j<n_*$, and let a \emph{centre of level $H$} be a tuple consisting of the bounded part, free
coordinates $b_j/p^{h_j}\le\theta$ with $p\nmid b_j$ and $h_j\in[\beta_j,2\beta_j]$, and a last coordinate $z=m/p^v$,
$p\nmid m$, with $v=H-k$, $k=\sum_ig_i+\sum_jh_j$, and value $\xi\in[1,p^d)$. Then every free loss exceeds $6s_m$
times the sum of the preceding losses plus $G_0$, and $k\le\kappa H$ once $\lambda\le\kappa(24s_m)^{-n_*}/4$ and $H$
is large; for $n_*=1$ there is no free part and $k$ is fixed. The minimal weight of an assignment of the
centre to the slots of $\nu^*$ is $w^*=s_*H+c_g$ with $c_g=\sum_i(u_i-s_*)g_{m_*+1-i}$ fixed, attained by exactly
$\prod_a\nu^*_a!$ slot tuples, and every other assignment has weight at least $w^*+G_0$. Put
\[
\pi_H=s_*H+c_g+e_{\rm prec},\qquad V_H=H^{n_*-1}p^H,\qquad \delta_H=\epsilon/V_H .
\]
Call a slot tuple of a monomial in $T$ \emph{relevant} if its weighted level is below $\pi_H$; every relevant
tuple has unweighted level $H'\le H+C$ for a constant $C$. A centre is \emph{isolated} if every relevant tuple
whose value lies within $\delta_H$ of $\xi$ is one of the $\prod_a\nu^*_a!$ minimal assignments of the centre. At an
isolated centre $M_P\bmod p^{\pi_H}$ is constant on $(\xi-\delta_H,\xi)$ and on $[\xi,\xi+\delta_H)$, and jumps by
$A_{\nu^*}\prod_a\nu^*_a!\,p^{w^*}\not\equiv0\pmod{p^{\pi_H}}$ at $\xi$.

\begin{lemma}[Weighted averaged isolation]\label{lem:wiso}\evidence{Lean-supported paper proof}
Assume $r_*\ge2$. Choose, in this order: $c>0$ small; $\kappa\in\bigl(0,\min(c/d,\,1/(2(s_*+1)),\,1/(s_m+2))\bigr)$;
$\lambda>0$ as above; $\mu_0\in(0,\min(1/d,1/(2r)))$ smaller than $\lambda/3$, than $1/(2(s_*+1))-\kappa$ and than
the margin $\sigma_r$ of \Cref{lem:margins}; and then $A,H_*\ge1$ and $\epsilon_0>0$. For $H\ge H_*$,
$p^{-\mu_0H/8}\le\epsilon\le\epsilon_0$, $0<\tau\le1$, and every union $E\subset[1,p^d)$ of $J_E\ge1$ intervals, the
centres with value in $E$ number at least $A^{-1}|E|V_H-AJ_EH^Ap^{\kappa H}$, and those of them that are not
isolated number at most
\[
A\epsilon\bigl(\tau^{1-d}|E|+\tau\bigr)V_H+AJ_EH^Ap^{(1-\mu_0)H}.
\]
\end{lemma}

\begin{proof}
\emph{Centres.} Each free coordinate and the last coordinate range over unit grids; the derivative of $z\mapsto
z^d$ is at most $dp^{d-1}$, so for each choice of the other coordinates the last one gives $\gg|E|p^v-O(J_E)$
values in $E$. There are $\asymp H^{n_*-1}$ free loss vectors and $\asymp p^{\sum_jh_j}$ free numerators for each, and
the bounded part is fixed; this gives the first bound, including $n_*=1$.

\emph{Exact coincidences.} Let a relevant competitor with largest loss $v'$ lie within $\delta_H$ of a centre,
and let $B$ be the difference of the remaining power sums. If $B=0$, rational spacing forces equality of the
largest coordinates once $C_d\delta_H<p^{-H-C}$, which holds for $\epsilon_0$ small. The common coordinate sits in a
slot of order $s_*$, since $v>(s_m+1)k+O(1)$ for small $\kappa$ and a larger order would exceed the weight
budget. Remove its weight. Peel the free coordinates downwards as in \Cref{thm:several}: at a free loss $h$ the
remaining weight is at most $s_*h+s_m\sum_{\rm earlier}h_i+O(1)$, so the remaining unweighted loss is below $3h/2$,
and \Cref{lem:avgiso}'s peeling step (the Lean lemma \leanR{PowerSumPeeling}{372}{\ttt{dominant\_loss\_unique\_top\_coord}})
gives a unique competitor coordinate of loss $h$; a larger order on it exceeds the budget. If its numerator
differs from the centre's, fixing the smaller coordinates leaves a nonzero equation $a^d-c^d=M$ with
$O(p^{\epsilon_{\rm div}dh})$ solutions (\leanR{PowerSumPeeling}{79}{\ttt{card\_pow\_diff\_eq\_le\_divisors}} and the
divisor bound), while the competitor's smaller losses cost at most $p^{h/2+O(1)}H^{O(r)}$ choices relative to the
$p^h$ numerators; with $\epsilon_{\rm div}\le1/(6d)$ this saves $p^{h/3}$. A matching numerator peels. After the $n_*$
slots of order $s_*$ have matched, lexicographic minimality forces each further order to be at least the
corresponding target order, and the bounded coordinates are forced as in \Cref{thm:several} by the sequential
choice of their numerators. Positivity excludes unpaired leftover slots. Hence exact coincidences off the minimal
assignments number $O(H^{O(r)}p^{H-\lambda H/3})$, or none if $n_*=1$. A degree-one competitor in $T$ has order
above $s_*$ and costs $O(H^{O(r)}p^{(\kappa+s_*/(s_*+1))H})$ even when charged once per anchor; this is admissible
because $\kappa<1/(2(s_*+1))$.

\emph{Rows.} For a competitor of degree $j\in[2,r]$ with unweighted level $H'\le H+C$, its largest loss satisfies
$H'/r\le v'\le H'$, and there are $O(H^{O(r)}p^{k+H'-v'})$ rows for given losses. A nonzero $|B|<p^{-cH}$ forces a
competitor coordinate of loss above $cH/d$, and varying it over the admissible interval saves
$p^{(c/d-\kappa)H}$. For $|B|\ge p^{-cH}$ the rows are those of \Cref{lem:avgiso}, with the derivative test of
order $q\le r+1\le d$ (\Cref{lem:qtest,lem:fB}), row cutoff $\lceil p^kH^{n_*+1}/\epsilon\rceil$, and the margins of
\Cref{lem:margins}, which apply because $H'\le H+O(1)$ and $v'\ge H'/r$. Intersecting each row with the
$O(\delta_H)$-enlargement $E^+$ of $E$ gives at most $J_E+O_r(H)$ pieces by monotonicity; bounding each piece by the
full row bound costs at most $J_EH^{O(r)}$. These errors fit the last term of the claim.

\emph{Main term.} What remains is $O(\epsilon)$ times the weighted count of competitor points in $E^+$, allowing
non-unit numerators for the fine row grid: each anchor contributes $O(\epsilon p^{-k}H^{1-n_*}+1/L_{\rm ET})$, which
sums to $O(\epsilon)$. For a competitor with $j_0\ge1$ slots of order $s_*$, fix its larger orders $o_i>s_*$ and their
losses $e_i$; summing over its $j_0$ minimal losses costs $O(H^{j_0-1}p^{H+C-\sum_io_ie_i/s_*})$, and multiplying by
$p^{\sum e_i}$ and summing gives a convergent product. If $j_0=0$ the sum is exponentially smaller. So the weighted
count over all competitors is $O(V_H)$, for every multiplicity $j_0$.

\emph{Regular and irregular competitors.} Put $\tau'=\max(\tau,p^{-\mu_0H})$ and $g_0=H/(2r)$. Loss vectors of total
loss below $H/2$ cost $O(H^{O(r)}p^{H/2})$. Otherwise some coordinate has loss $g\ge g_0$. If one such coordinate is
at least $\tau'$, fix the others and vary it: at most $p^g|E^+|/(d\tau'^{d-1})+O(J_E)$ grid points qualify, and
summing as in the main term bounds these competitors by $O(\tau'^{1-d}|E^+|V_H)$ with error $O(J_EH^{O(r)}p^{H-g_0})$.
If all coordinates of loss at least $g_0$ lie below $\tau'$, choose one; its positive grid has at most $\tau'p^g$
such points, which gives $O(\tau'V_H)$ in all. This holds even when fixed higher-order slots translate the power sum
into $E$ and the density of the remaining coordinates is singular. Finally $|E^+|\le|E|+O(J_E\delta_H)$, so the
enlargement costs $O(J_Ep^{(d-1)\mu_0H})$ after multiplication by $\epsilon$, and $\mu_0<1/(2r)$ absorbs all errors.
\end{proof}

The regular/irregular split replaces the Lebesgue-point localization of \Cref{thm:several}; this is what makes
the estimate uniform over unions of intervals.

\begin{theorem}[Rates for several orders]\label{thm:severalrate}\evidence{Paper proof}
Let $1\le s_1<\cdots<s_m$ be integers, $X_a=C(b,p,s_a,d)$, and let $P\in\Q[X_1,\ldots,X_m]$ be nonconstant of total
degree $r\le d-1$. Then
\[
D_N(P(X_1,\ldots,X_m))\ll_{P,b,p,d,\mathbf s}(\log\log N)^{-d/(d-1)} .
\]
If the largest weighted degree $\max\{\sum_as_a\nu_a\}$ over the monomials of $P$ is attained by a single monomial
of degree one, then $D_N(P(X_1,\ldots,X_m))\ll_PN^{-\eta_P}$ for some $\eta_P>0$; this applies to every nonconstant
linear form.
\end{theorem}

\begin{proof}
\emph{Small balls.} Assume first $r_*\ge2$. We show that for every $w\in\Z_p$ and $L\ge3$,
\begin{equation}\label{eq:wsmallball}
\bigl|\{u\in[1,p^d):M_P(u)\equiv w\pmod{p^L}\}\bigr|\ll(\log L)^{-d/(d-1)} .
\end{equation}
Iterate \Cref{lem:wiso} on an outer cover of this set by unions of intervals, starting from $[1,p^d)$, exactly as in
the proof of \Cref{thm:smallball}. At a step with current measure $x>(\log L)^{-d/(d-1)}$ take
$\tau=(x/p^d)^{1/d}$ and $\epsilon=c_Px^{(d-1)/d}$ with $c_P$ small, so that the first term of the collision bound
costs at most a quarter of the main term. Use levels starting at $H_0\asymp1+\log\log L$ with a sufficiently large
constant and multiplied each time by an integer $\Gamma>4/\mu_0$; the previous deletions have
$O((j+1)H_{\rm prev}^{O(r)}p^{H_{\rm prev}})$ components, so the second error term and the boundary loss are $o(xV_H)$, and
$\epsilon\ge p^{-\mu_0H/8}$. At each isolated centre one flank of length $\delta_H/2$ misses the target residue as long as
$\pi_H\le L$, and the flanks are disjoint up to $O(J_E)$ boundary flanks. The new measure is at most
$x(1-c'_Px^{(d-1)/d})$, so $x^{-(d-1)/d}$ increases by a fixed amount at each of the $\gg\log L$ available levels, which
proves \eqref{eq:wsmallball}. If $r_*=1$, the residue $A_{\nu^*}\lfloor u^{1/d}\rfloor$ is nonzero modulo
$p^{v_p(A_{\nu^*})+1}$, because $\lfloor u^{1/d}\rfloor\in[1,p-1]$, and the set in \eqref{eq:wsmallball} is empty for
$L>v_p(A_{\nu^*})$.

\emph{Digits.} Write $b^eP=Y/D$ with $Y=\sum_\nu A_\nu X^\nu$ integral and $D>0$ prime to $b$; the shift by $e$ costs $e/N$.
Take $L=\lfloor c_0\log N\rfloor$ with $c_0>0$ so small that $tp^{dJ_1(L)}\le\sqrt N$ for large $N$ (recall
$J_1(L)\le2L+2$), that $L\le(S^*/(8d))\log_pN$, and that the number $O(\log N\cdot L^rp^{L/s_1})$ of cells of
$M_P\bmod p^L$ in all shells through $N$ is $o(N^{r/d})$. Then every shift $n\ge\sqrt N$ is frozen, and by
\eqref{eq:freezeJ} its normalized prefix is $M_P(u(n))$ modulo $p^L$; division by $D$ changes the depth by a fixed
amount. Cut $[1,N]$ at the $O(N^{r/d})$ insertion positions, the shell boundaries, the cell boundaries and
$\lceil\sqrt N\rceil$, trim the final $B=\lceil2\log_bN\rceil$ positions of each run, and discard the runs below
$\lceil\sqrt N\rceil$, the runs with untrimmed length below $N^{1-r/d-\epsilon_1}$, $\epsilon_1=1/(6d)$, and the runs whose
residue is $0$ modulo $p^{L-v_p(A_{\nu^*})-O(1)}$. In a full or partial shell the shift grid has mesh
$1/(tp^{dJ})$, so counting on cells costs $O(1)$ per cell beyond the failed measure times $tp^{dJ}$; summing the
geometric shell sizes, \eqref{eq:wsmallball} discards $O(N(\log L)^{-d/(d-1)})$ shifts, and nothing in the case of a
single linear maximizer. On a retained run the reduced denominator is $p^{e'}Q'$ with $Q'$ dividing a fixed integer
prime to $pb$ and $S^*J-L-O(1)\le e'\le S^*J+O(1)$, so after multiplication by $h\le N^{\eta_0}$ the pole exponent
$e$ satisfies $c_1\log N\le e\le c_2\log N$ with constants $c_1,c_2>0$, for small $\eta_0$. The stride
$\tau=\lcm(t,\ord_{Q'}b)$ is bounded, $A(b^\tau)=O(1)$, and $M_0\ge T/\tau\gg N^{1-r/d-\epsilon_1}$, so
$p^{\theta e}\le M_0$ with $\theta=\min(1,(d-r-1/2)/S^*)>0$; here $d\ge r+1$ is used. The Erd\H{o}s--Tur\'an
inequality with $\lceil N^{\eta_0}\rceil$ frequencies, \Cref{cor:split} and \Cref{lem:uniformorbit} give a power
saving on each retained run, and \Cref{lem:blocks} with $\epsilon=N^{-1}$ completes the proof.
\end{proof}

With \Cref{thm:higher,thm:several,thm:effective,thm:severalrate} the proof of \Cref{thm:D} is complete.

\section{Infinite sums}\label{sec:sums}

Let $s_3=2$ and $s_\ell=1$ for primes $\ell\ge5$, and for every odd prime $\ell$ put
\[
Y_\ell=\sum_{j\ge1}\frac{1}{\ell^{s_\ell j}\,2^{16^{\ell^j}}},\qquad
Y=\sum_{\ell\ \text{odd prime}}Y_\ell .
\]
The insertions are indexed by the odd prime powers $m=\ell^j$ (the \emph{labels}); the label $m$ sits
at position $16^m$, the positions are distinct, and every coefficient is at most $1/5$. For $n\ge1$ let
$L=L(n)=\lfloor\log_{16}n\rfloor$, so that the terms inserted by shift $n$ are those with label at most
$L$, and for $\ell\le L$ let $J_\ell=\lfloor\log_\ell L\rfloor$.

\begin{lemma}\label{lem:Yden}\evidence{Paper proof}
Let $r_\ell$ $(\ell$ odd prime$)$ be finitely supported rationals with odd denominators, and let $R_n$ be
$2^n$ times the sum of the terms of $\sum_\ell r_\ell Y_\ell$ with label at most $L(n)$. For every $\ell$
with $r_\ell\ne0$, $v_\ell(R_n)=v_\ell(r_\ell)-s_\ell J_\ell$ for all sufficiently large $n$, and $R_n$ is
integral at $2$ and at every prime outside the support of $(r_\ell)$ and the denominators of the
$r_\ell$. With $r_\ell=1$ for all $\ell$, the reduced denominator of $R_n$ is exactly
$q(n)=3^{2J_3}\prod_{5\le\ell\le L}\ell^{J_\ell}$ for every $n$ with $L(n)\ge3$.
\end{lemma}

\begin{proof}
The terms $r_\ell\ell^{-s_\ell j}2^{n-16^{\ell^j}}$ with $\ell^j\le L$ are $2$-integral. At a prime $\ell$,
the terms from other axes have valuation at least $\min_{\ell'}v_\ell(r_{\ell'})$, a constant; on the axis
$\ell$ the term with $j=J_\ell$ has valuation $v_\ell(r_\ell)-s_\ell J_\ell$ and the others have strictly
larger valuation. Since $J_\ell\to\infty$, the top term has the unique least valuation for large $n$.
When all $r_\ell=1$ the other axes are $\ell$-integral and the conclusion holds as soon as $\ell\le L$.
\end{proof}

We use the elementary bound $\prod_{\ell\le x}\ell<4^x$ over primes: the primes in $(m+1,2m+1]$ divide
$\binom{2m+1}{m}\le4^m$, and induction on $x$ completes the proof. It gives
$q(n)\le3^{J_3}\prod_{\ell\le L}\ell^{J_\ell}\le L^{\sqrt L+1}4^L$, since only primes $\ell\le\sqrt L$
have $J_\ell\ge2$.

\begin{theorem}[All odd primes]\label{thm:Y}\evidence{Lean-supported paper proof}
The number $Y$ is normal to base $2$ and $D_N(Y;2)\ll1/\log(2N)$ for all $N\ge1$. Moreover $Y$ is a
Liouville number.
\end{theorem}

\begin{proof}
Let $L(n)\ge3$, write $q=q(n)=3^eQ$ with $e=2J_3$ and $Q$ odd and prime to $3$, and let
$\tau=\lcm(2,\ord_Q(2))$, with the conventions $\ord_1(2)=1$ and that an empty maximum is zero (both
apply when $L\le4$, where $Q=1$). Then $2^\tau\equiv1\pmod{3Q}$ and, lifting the exponent from $v_3(4-1)=1$,
\[
c=v_3(2^\tau-1)=1+v_3(\ord_Q(2))\le1+\max_{5\le\ell\le L}v_3(\ell-1)\le1+J_3,
\]
because the order of $2$ modulo $\ell^f$ divides $\ell^{f-1}(\ell-1)$ and $3^{v_3(\ell-1)}\le\ell-1<L$.
Hence $e-c\ge J_3-1\ge0$, and \Cref{lem:grid} shows that any $L'$ consecutive points $\fr{a2^j/q}$
with $3\nmid a$ (in particular with $a$ prime to $q$, as for the reduced prefixes) have count error at most $L'3^{1-J_3}+\tau3^{e-c}\le9L'/L+q$ on every half-open arc, using
$3^{J_3+1}>L$ and $\tau\le2Q$. The doubled order at $3$ is what makes the grid grow: the loss $c$ can be
as large as $J_3$, but $e=2J_3$. The Lean-supported steps are the order and count statements cited in
\Cref{lem:grid}.

Let $N$ be large, $N_0=\lceil\sqrt N\rceil$, and apply \Cref{lem:blocks}: leave $[0,N_0)$ uncovered and
split $[N_0,N)$ at the insertion positions. Fewer than $\log_{16}(2N)$ labels have position below $2N$,
so there are $O(\log N)$ blocks; on each, $R_{n+1}=2R_n$ and the models are consecutive points
$\fr{a2^j/q}$ with $q$ constant, and $L(n)\ge\frac12\log_{16}N-1$. With $w=\lceil\log_2N\rceil+1$, a
shift $n<N$ with no insertion in $(n,n+w]$ satisfies $0\le2^nY-R_n<1/N$, and at most $w\log_{16}(2N)$
shifts fail this. Since $L(N)\le\log_{16}N$, $q\le L^{\sqrt L+1}4^L\le N^{1/2+o(1)}$, so
\[
D_N(Y)\le\frac{N_0+w\log_{16}(2N)+O(N/\log N)+O(\log N)N^{1/2+o(1)}}{N}+\frac2N\ll\frac1{\log N}.
\]
Small $N$ are absorbed in the constant.

For the Liouville property, given $h\ge1$ choose distinct primes $r_1,\ldots,r_{2h}$ and, by the Chinese
remainder theorem, arbitrarily large $z$ with $r_{2i-1}r_{2i}\mid z+i$ for $1\le i\le h$. None of
$z+1,\ldots,z+h$ is a prime power, so there are consecutive labels $u<v$ with $v-u>h$ and $u$
arbitrarily large. The truncation $P$ through label $u$ has denominator at most $2^{16^u}q$ with
$\log_2q=O(u)$, while $0<Y-P\le2^{1-16^v}\le2^{1-16^{h}\cdot16^u}$. As $h$ is arbitrary and $Y$ is
irrational, $Y$ is a Liouville number.
\end{proof}

If only $3$ and the primes $\ell\equiv2\pmod3$ are used, $3\nmid\ord_Q(2)$, so $c=1$ and the weight
$3^{-j}$ already gives a growing grid at $3$; with positions $2^{\ell^{2j}}$ this yields the rate
$O((\log N)^{-1/2})$ by the same argument.

\begin{theorem}[Finite combinations]\label{thm:combos}\evidence{Lean-supported paper proof}
For every finite set $S$ of odd primes and rationals $r_\ell$ $(\ell\in S)$, not all zero, the number
$\sum_{\ell\in S}r_\ell Y_\ell$ is normal to base $2$. Consequently $(Y_\ell)_{\ell\in S}$ is jointly
normal: the points $(\fr{2^nY_\ell})_{\ell\in S}$ are uniformly distributed in $\T^S$.
\end{theorem}

\begin{proof}
Multiplying a number by $2^a$ changes $D_N$ by at most $a/N$, so we may assume that all $r_\ell$ have odd
denominators. Choose $\ell_0\in S$ with $r_{\ell_0}\ne0$, let $F$ be the finite set of odd primes dividing
$\prod_{\ell\in S}\ell$ or a denominator of some $r_\ell$, and let
$\beta=\max_{\ell\in F\setminus\{\ell_0\}}v_{\ell_0}(\ell-1)$ (zero if empty). By \Cref{lem:Yden}, for large
$n$ the reduced denominator of $R_n$ is $q=\ell_0^eQ$ with $e=s_{\ell_0}J_{\ell_0}-v_{\ell_0}(r_{\ell_0})\to\infty$
and $Q$ supported on $F\setminus\{\ell_0\}$. The order $\ord_Q(2)$ divides $\lcm_\ell\ell^{f_\ell-1}(\ell-1)$,
so its $\ell_0$-valuation is at most $\beta$. With $t_0=\ord_{\ell_0}(2)$ and $\tau=\lcm(t_0,\ord_Q(2))$,
lifting the exponent gives $c=v_{\ell_0}(2^\tau-1)=v_{\ell_0}(2^{t_0}-1)+v_{\ell_0}(\tau/t_0)\le
v_{\ell_0}(2^{t_0}-1)+\beta=:c^*$, independent of $n$. \Cref{lem:grid} gives count error at most
$L'\ell_0^{c^*-e}+\tau\ell_0^{e-c}\le L'\ell_0^{c^*-e}+q$, and $q\le(\log N)^{O(1)}$ for $n<N$. The block and tail
argument of \Cref{thm:Y}, with $O(\log N)$ insertions and the tail multiplied by $\sum|r_\ell|$, gives
$D_N\to0$. Joint normality follows by applying this to the integer combinations $\sum h_\ell Y_\ell$,
$h\ne0$, together with Weyl's criterion and its multidimensional form
\cite[Chapter~1, Theorems~2.1 and~6.2, pp.~7 and~49]{KN}.
\end{proof}

Only a fixed finite set of primes enters $Q$ here, so the loss $c$ is bounded. For the infinite sum
$Y$ the loss can grow, and the doubled order at $3$ compensates for it.

\begin{theorem}[Rational totals]\label{thm:rationaltotal}\evidence{Lean-supported paper proof}
For every real $R>0$ there are rationals $c_\ell>0$, indexed by the odd primes, with
$\sum_\ell c_\ell Y_\ell=R$. For rational $R$ the total is not normal, although every finite subfamily
of the positive terms $c_\ell Y_\ell$ is jointly normal.
\end{theorem}

\begin{proof}
For positive generators $g_i$ and $R>0$, choose inductively rationals with $e_{i-1}/2<c_ig_i<e_{i-1}$,
where $e_0=R$ and $e_i=e_{i-1}-c_ig_i$; then $0<e_i<R2^{-i}$. This is the kernel-verified
\lean{PositiveSummation}{15}{\ttt{hasSum\_positive\_rational\_multiples}}, applied to an enumeration of
the $Y_\ell$. Joint normality of finite subfamilies is \Cref{thm:combos}, since nonzero rational scaling
preserves its hypothesis. A rational number has an eventually periodic expansion and is not normal.
\end{proof}

The tails of a series of normal numbers cannot all be negligible. For real $z$ let
$M_N(z)=N^{-1}\sum_{n<N}\dist{b^nz}$ and $\overline M(z)=\limsup_NM_N(z)$.

\begin{proposition}[Consecutive tails]\label{prop:tails}\evidence{Lean-supported paper proof}
Let $\sum_iS_i$ be a convergent series of numbers normal to base $b$, with tails $E_k=\sum_{i>k}S_i$.
Then for every $k$,
\[
\overline M(E_k)+\overline M(E_{k+1})\ge\frac14\qquad\text{and}\qquad
\limsup_N\bigl(D_N(E_k)+D_N(E_{k+1})\bigr)\le1 .
\]
In particular at most one of any two consecutive tails $z$ satisfies $D_N(z)\to1$, that is, by
\Cref{thm:sharpcited}, $\sup_x|D_N(x+z)-D_N(x)|\to0$.
\end{proposition}

\begin{proof}
Since $S_{k+1}=E_k-E_{k+1}$, $\dist{b^nS_{k+1}}\le\dist{b^nE_k}+\dist{b^nE_{k+1}}$. By uniform distribution
\cite[Chapter~1, Corollary~1.1, p.~3]{KN} the averages of the left side tend to $\int_0^1\dist{y}\,dy=1/4$.
For the second inequality apply \Cref{thm:sharpcited} with $x=E_k$ and $z=-E_{k+1}$: since
$D_N(-E_{k+1})=D_N(E_{k+1})$, it gives $D_N(S_{k+1})\ge D_N(E_k)+D_N(E_{k+1})-1$, and $D_N(S_{k+1})\to0$.
The Lean-supported step is \Cref{thm:sharpcited}.
\end{proof}

In \Cref{thm:Y} normality of the total therefore comes from the arithmetic grid at $3$, which survives
every new prime, and not from smallness of the tails. The translations $z$ with $D_N(z)\to1$ are
classified in \cite[Theorem~3.3]{CompanionH}; in base $2$ they are exactly the $z$ with $M_N(z)\to0$.

\section{The boundary \texorpdfstring{$d=r$}{d = r}}\label{sec:boundary}

Throughout this section $p\nmid b$ is an odd prime, $s\ge1$, $d\ge2$ and $C=C(b,p,s,d)$. For an ordered
$r$-tuple $(i_1,\ldots,i_r)$ of positive integers the \emph{level} is $v_p(i_1)+\cdots+v_p(i_r)$ and the
\emph{position} is $i_1^d+\cdots+i_r^d$, so that $C^r=\sum_mc(m)b^{-tm}$ with $c(m)$ the sum of $p^{-s\ell}$ over the
tuples of level $\ell$ at position $m$. At $d=r$ the ordered tuples with position at most $X$ number
$\sim V_{r,d}X$, with $V_{r,d}=\Gamma(1+1/d)^r/\Gamma(1+r/d)$, so the sparsity hypothesis (i) of
\Cref{prop:criterion} fails.

\subsection{What survives, and a barrier}

\begin{proposition}[Depth is tight in density]\label{prop:depthtight}\evidence{Paper proof}
Let $d\ge2$ and let $P\in\Q[x]$ have degree $r\ge2$; $d\le r$ is allowed. There is $C_P$ such that for every
$L\ge3$ the shifts $n$ with $v_p(q_n)<rsJ(n)-L$ have upper density at most $C_P(\log L)^{-d/(d-1)}$.
Consequently, for every function $\omega(n)\to\infty$, $v_p(q_n)\ge rsJ(n)-\omega(n)$ on a set of shifts of
density one.
\end{proposition}

\begin{proof}
\Cref{lem:freezer} uses no sparsity, and \Cref{thm:smallball} holds for every $d\ge2$. The finitely many
$L$ with $3\le L\le v_p(A)+2$ are absorbed into $C_P$; for larger $L$ the integer sampling of the cells of
$M_r\bmod p^{L-v_p(A)}$ is as in \Cref{lem:reductionr}. For the last claim, the exceptional
set for $\omega$ is contained in $\{n:\omega(n)<L\}\cup\{n:v_p(q_n)<rsJ(n)-L\}$ for every fixed $L$.
\end{proof}

So the $p$-adic half of the method survives at $d=r$. The real half does not: the runs between insertions
are too short for their denominators, and no estimate that sees only the denominator can help.

\begin{lemma}[Denominators do not determine digits]\label{lem:barrier}\evidence{Paper proof}
Let $q=p^e$ and let $T\ge1$ with $b^T\le q/2$. Then every word of $T$ base-$b$ digits is the leading word of
$a/q$ for some $a$ prime to $p$.
\end{lemma}

\begin{proof}
Consecutive integers prime to $p$ differ by at most $2$, so as $a$ runs through them from $1$ to $q-1$ the
integer $\lfloor ab^T/q\rfloor$ advances by at most $1$ at each step and takes every value in $[0,b^T)$.
\end{proof}

\begin{lemma}[No uniform cancellation on short orbits]\label{lem:nocancel}\evidence{Paper proof}
Let $q=p^e$, $0<\epsilon<1$ and $T\ge1$ with $b^T\le\epsilon q$. Then for every nonzero integer $h$,
\[
\Bigl|\sum_{j<T}e(hb^j/q)-T\Bigr|\le2\pi|h|\frac{b^T-1}{(b-1)q}<\frac{2\pi|h|\epsilon}{b-1}.
\]
\end{lemma}

\begin{proof}
$|e(\theta)-1|\le2\pi|\theta|$ and $\sum_{j<T}b^j=(b^T-1)/(b-1)$.
\end{proof}

Thus a bound for $\sup_{p\nmid a}|\sum_{j<T}e(hab^j/q)|$ uniform in the numerator cannot improve on the
trivial bound by more than $O_h(1)$ when $T\le\log_b(\epsilon q)$, whatever method produces it. At $d=r$ the
shifted prefixes have $\log q_n\asymp\log N$ on a set of density close to one (\Cref{prop:depthtight}), while
the gaps between insertions are short: for $d=2$ the distinct positions are the sums of two positive
squares, of which there are $\asymp X/\sqrt{\log X}$ up to $X$ by Landau's theorem \cite{Landau08}, and for
$d=3$ they are the sums of three positive cubes, of which there are at least $X^{0.917}$ by Wooley's theorem
\cite{Wooley15} and at most $(V_{3,3}/6+o(1))X$, conjecturally a positive proportion \cite{DHL06}.

\subsection{Retirement of the integer part at \texorpdfstring{$d=2$}{d = 2}}

\begin{proposition}[Level filtration]\label{prop:filtration}\evidence{Paper proof}
For all $r\ge1$, $d\ge2$ and $H\ge0$,
\[
p^{sH}C^r=X_H+Z_H,\qquad X_H=\sum_ma_H(m)b^{-tm},\qquad Z_H=\sum_mz_H(m)b^{-tm},
\]
where $a_H(m)=\sum p^{s(H-\ell)}$ over the tuples at $m$ of level $\ell\le H$ is a nonnegative integer and
$z_H(m)=\sum p^{-s(\ell-H)}$ over those of level $\ell>H$. Moreover $\sum_{m\le Y}a_H(m)\le p^{sH}V_{r,d}Y^{r/d}$
for all $Y\ge0$, the shifted prefixes satisfy $\fr{R_n(Z_H)}=\fr{p^{sH}R_n(C^r)}$, and their reduced
denominators satisfy $q_n(Z_H)=q_n(C^r)/p^{\min(sH,v_p(q_n(C^r)))}$.
\end{proposition}

\begin{proof}
Split the absolutely convergent tuple expansion by level. Each lattice point $(i_1,\ldots,i_r)$ owns the unit
cube $\prod_c[i_c-1,i_c]$, which lies in $\{x\ge0:\sum x_c^d\le Y\}$ when $\sum i_c^d\le Y$, so the tuples number at
most $V_{r,d}Y^{r/d}$. The prefix of $X_H$ is an integer after every shift.
\end{proof}

For an integer series a self-contained transfer suffices.

\begin{lemma}[Sparse integer series]\label{lem:sparseint}\evidence{Paper proof}
Let $F=\sum_{m\ge1}a_mb^{-m}$ with integers $a_m\ge0$ satisfying $\sum_{m\le X}a_m\le C_FX$ for all $X\ge1$, and
let $S(X)=\#\{m\le X:a_m\ne0\}$. For integers $1\le L\le N$ and every $\epsilon\in(0,1)$,
\[
1-D_N(F)\le\epsilon+\frac{L\,S(N+L)}N+\frac{6C_Fb^{-L}}{(b-1)\epsilon}.
\]
\end{lemma}

\begin{proof}
For $n<N$, $\fr{b^nF}$ is the fractional part of $T_n=\sum_{m>n}a_mb^{n-m}\ge0$. At most $LS(N+L)$ shifts have an
insertion in $(n,n+L]$; at the others $T_n=\sum_{m>n+L}a_mb^{n-m}$. Summing over $n<N$, the terms with
$m\le N+L$ contribute at most $C_F(N+L)b^{-L}/(b-1)\le2C_FNb^{-L}/(b-1)$, and by partial summation the terms
with $m>N+L$ contribute at most $(b-1)^{-1}\sum_{j>L}a_{N+j}b^{-j}\le C_F(N+L+2)b^{-L}/(b-1)\le4C_FNb^{-L}/(b-1)$. By
Markov's inequality at most $6C_FNb^{-L}/((b-1)\epsilon)$ of the remaining shifts have $T_n\ge\epsilon$. All other
points lie in $[0,\epsilon)$, and \Cref{lem:arcs} gives the claim.
\end{proof}

\begin{theorem}[Absorption of the integer part]\label{thm:absorb}\evidence{Paper proof}
Let $d=r=2$. Uniformly in $H\ge0$, $1-D_N(X_H)\ll_{b,p,t}(\log\log N+sH)/\sqrt{\log N}$ for $N\ge3$. In particular
$D_N(X_H)\to1$ for every fixed $H$, and also for $H=H(N)=o(\sqrt{\log N})$.
\end{theorem}

\begin{proof}
As a series in $b^{-n}$, $X_H$ has coefficient $a_H(n/t)$ at $t\mid n$ and $0$ otherwise. Its support up to $X$
is contained in $t$ times the set of sums of two positive squares, which by Landau's theorem has
$\ll X/\sqrt{\log X}$ elements up to $X$, and by \Cref{prop:filtration} its mass up to $X$ is at most
$C_FX$ with $C_F=\pi p^{sH}/(4t)$. Put $\rho=(1+S(2N))/N\ll1/\sqrt{\log N}$, $\epsilon=\rho$ and
$L=\max\bigl(1,\lceil\log_b(6C_F/((b-1)\rho^2))\rceil\bigr)$, so that $L\le2\log_b(1/\rho)+sH\log_bp+O_{b,t}(1)$. If
$\rho\ge1$ (which happens only for bounded $N$) or $L>N$, then $\rho(L+2)>1\ge1-D_N(X_H)$. Otherwise
$\epsilon\in(0,1)$ and $1\le L\le N$, and \Cref{lem:sparseint} gives $1-D_N(X_H)\le\rho(L+2)$ in every case.
\end{proof}

\begin{lemma}[Integer dilation]\label{lem:intmult}\evidence{Paper proof}
For every finite sequence $(z_n)$ in $\T$ and every integer $M\ge1$, $D(Mz)\le M\,D(z)$ and
$1-D(Mz)\le M(1-D(z))$, where $Mz=(Mz_n)$.
\end{lemma}

\begin{proof}
The preimage of an arc $J$ under $z\mapsto Mz$ consists of $M$ arcs of length $|J|/M$, each with error at most
$D(z)$. For the second inequality use $1-D=\inf_I(1-\mu(I)+|I|)$ (\Cref{lem:arcs}). If $|I|<1/M$, the image $MI$
is an arc of length $M|I|$ with $\mu_{Mz}(MI)\ge\mu_z(I)$, so $1-D(Mz)\le1-\mu_z(I)+M|I|\le M(1-\mu_z(I)+|I|)$; if
$|I|\ge1/M$, then $1-\mu_z(I)+|I|\ge1/M\ge(1-D(Mz))/M$.
\end{proof}

The converse of the second inequality fails: for $b=2$, $D_N(1/3)=2/3$ for $N\ge2$ while $D_N(1)=1$. For
normality the converse holds (a case of Wall's theorem on rational multiples): if $Mx$ is normal, every
weak limit $\nu$ of the orbit measures of $x$ is $T$-invariant (\Cref{lem:necessity}) and satisfies
$\nu(I)\le\nu(M^{-1}(MI))=M|I|$ for closed arcs $I$ of length below $1/M$, so $\nu\le M\lambda$, and ergodicity of
Lebesgue measure $\lambda$ under $T$ forces $\nu=\lambda$.

\begin{corollary}[Fixed-depth retirement]\label{cor:retire}\evidence{Lean-supported paper proof}
Let $d=r=2$ and fix $H\ge0$. Then $C^2$ is normal if and only if $Z_H$ is normal;
$\limsup_ND_N(p^{sH}C^2)=\limsup_ND_N(Z_H)$; and $1-D_N(C^2)\ge p^{-sH}(1-D_N(Z_H)-\delta_N(H))$ with
$\delta_N(H)=1-D_N(X_H)\to0$. For $H=0$, $|D_N(C^2)-D_N(Z_0)|\le\delta_N(0)$.
\end{corollary}

\begin{proof}
\Cref{thm:sharpcited} with $x=Z_H$ and $z=X_H$ gives $|D_N(p^{sH}C^2)-D_N(Z_H)|\le\delta_N(H)$, and
\Cref{thm:absorb} gives $\delta_N(H)\to0$. Normality of $p^{sH}C^2$ and of $C^2$ are equivalent by
\Cref{lem:intmult} and the converse just proved, and the inequality for $1-D_N(C^2)$ is the second part of
\Cref{lem:intmult} with $M=p^{sH}$. The Lean-supported step is \Cref{thm:sharpcited}.
\end{proof}

So the discrepancy problem for $C^2$ is exactly the problem for the pairs of level at least one. Only level
zero is invisible to every frequency: a tuple of level $\ell$ contributes $p^{-s\ell}b^{n-tm}$, which lies in
$h^{-1}\Z$ for all $n\ge tm$ if and only if $s\ell\le v_p(h)$. Letting $H$ grow controls $C^2$ only at the
frequencies $p^{sH}h$, or through \Cref{lem:intmult} at the cost of the factor $p^{sH}$.

\subsection{The runs after retirement}

Let $d=r=2$, let $\mathcal S^+$ be the set of sums of two positive squares, and for $H\ge0$ let $\mathcal P_H$ be
the set of positions of $Z_H$, namely the $m=i^2+j^2$ with $i,j\ge1$ and $v_p(i)+v_p(j)>H$. A \emph{run} of $Z_H$
is a maximal shift interval $[tm,tm')$ between consecutive elements $m<m'$ of $\mathcal P_H$. Put
$k_H=\lceil(H+1)/2\rceil$.

\begin{theorem}[Runs after retirement]\label{thm:runs}\evidence{Paper proof}
Let $d=r=2$ and $H\ge0$.
\begin{enumerate}[label=\textup{(\roman*)},nosep]
\item $p^{2k_H}\mathcal S^+\subseteq\mathcal P_H\subseteq\mathcal S^+$, so $|\mathcal P_H\cap[1,Y]|\asymp_HY/\sqrt{\log Y}$ and the
mean run length up to $N$ is $\asymp_H\sqrt{\log N}$.
\item For $\gamma\in(1,5/3)$ and $\epsilon>0$ the proportion of shifts $n<N$ lying in runs of length at least
$\epsilon\log N$ is $\ll_\gamma(tp^{2k_H}/\epsilon)^{\gamma-1}(\log N)^{-(\gamma-1)/2}$, uniformly for
$p^{2k_H}\le N^{1/2}$.
\item For $L\ge3$, outside a set of upper density at most $C(\log L)^{-2}$,
$\log_bq_n(Z_H)\ge s\log_b(n/t)-(2s+L+sH)\log_bp$.
\end{enumerate}
Consequently, for $\epsilon<s/(2\ln b)$, almost every shift lies in a run of length $T<\log_b(q_n(Z_H)/2)$, on
which \Cref{lem:barrier,lem:nocancel} apply.
\end{theorem}

\begin{proof}
(i) If $m'=u^2+v^2$ with $u,v\ge1$, then $p^{2k_H}m'=(p^{k_H}u)^2+(p^{k_H}v)^2$ has a representation of level at
least $2k_H>H$; the counts follow from Landau's theorem, since the sums of two squares that are not sums of
two positive squares are perfect squares.

(ii) A run of length at least $G$ lies in a gap of $Q\mathcal S^+$, $Q=p^{2k_H}$, of length at least $G/t$. After
rescaling by $Q$ it suffices to bound the proportion of $[1,Y]$, $Y=N/(tQ)$, covered by gaps of $\mathcal S^+$ of
length at least $g=G/(tQ)$, plus the final gap. Since $k^2<k^2+1<(k+1)^2$ with $k^2+1\in\mathcal S^+$, a gap of
$\mathcal S^+$ contains at most one sum of two squares in its interior, so it is a gap of the set $\mathcal S$ of sums
of two squares or the union of two adjacent ones; hence the total length of gaps of $\mathcal S^+$ of length at
least $g$ is at most twice that of gaps of $\mathcal S$ of length at least $g/2$. Hooley's theorem
\cite{Hooley71} gives $\sum_{s_{n+1}\le Y}(s_{n+1}-s_n)^\gamma\ll_\gamma Y(\log Y)^{(\gamma-1)/2}$ for $0<\gamma<5/3$, and
Markov's inequality with $G=\epsilon\log N$ gives the bound. The final gap has length $O(tQY^{1/4})$ by the
Bambah--Chowla theorem \cite{BambahChowla47}, a proportion $O(Y^{-3/4})$, which is dominated by the main
term whenever it matters.

(iii) By \Cref{prop:filtration}, $v_p(q_n(Z_H))\ge v_p(q_n(C^2))-sH$, and \Cref{prop:depthtight} with $P(x)=x^2$
gives $v_p(q_n(C^2))\ge2sJ(n)-L$ off the stated set; also $2J>\log_p(n/t)-2$.

For the last claim restrict to $n\ge\sqrt N$, where $\log_b(q_n/2)\ge\frac s2\log_bN-O_{H,L}(1)>\epsilon\log N$,
intersect the sets of (ii) and (iii), and let $L\to\infty$.
\end{proof}

Integer jumps $z_H(m)\in\Z$ do not move the rational orbit, so the orbit of $\fr{R_n(Z_H)}$ restarts only at
the positions $\mathcal P_H^{\rm frac}=\{m\in\mathcal P_H:z_H(m)\notin\Z\}$, and an argument using the rational orbit alone may
merge runs across integer jumps. Such jumps occur: $16810=13^2+129^2=41^2+123^2=67^2+111^2$ has exactly six
ordered representations, all of level one for $p=3$, so $z_0(16810)=2$. For the merged runs we can prove only
the mean.

\begin{proposition}[Fractional jumps]\label{prop:fractional}\evidence{Paper proof}
Let $p\equiv3\pmod4$, $H\ge0$, and let $\mathcal F_p$ be the set of squarefree $n\in\mathcal S^+$ with $(n/p)=-1$. Then
$p^{2k_H}\mathcal F_p\subseteq\mathcal P_H^{\rm frac}$; more precisely $z_H(p^{2k_H}n)=2^{\omega_1(n)}p^{-s(2k_H-H)}$ for
$n\in\mathcal F_p$, where $\omega_1(n)$ is the number of prime factors $\equiv1\pmod4$ of $n$. Consequently
$|\mathcal P_H^{\rm frac}\cap[1,Y]|\asymp_HY/\sqrt{\log Y}$, and the mean length of the runs between consecutive
elements of $\mathcal P_H^{\rm frac}$ up to $N$ is $\asymp_H\sqrt{\log N}$.
\end{proposition}

\begin{proof}
For $p\equiv3\pmod4$, $v_p(i^2+j^2)=2\min(v_p(i),v_p(j))$, so every representation of $m=p^{2k}n$, $k=k_H$,
$p\nmid n$, has $i=p^ki'$, $j=p^kj'$ with $i'^2+j'^2=n$; and $p\mid i'$ would make $n$ a square modulo $p$. So every
representation has level exactly $2k>H$. A squarefree $n>1$ in $\mathcal S$ has no prime factor $\equiv3\pmod4$ and
is not a square, so its ordered positive representations number $r_2(n)/4=\sum_{e\mid n}\chi_{-4}(e)=2^{\omega_1(n)}$,
prime to $p$. For the count, the multiplicative function $\mu^2\mathbf1_{\mathcal S}$ has partial sums
$\sim c\,Y/\sqrt{\log Y}$, and its twist by $(\cdot/p)$ contributes $o(Y/\sqrt{\log Y})$, both by the
Selberg--Delange method \cite[Chapter~II.5]{Tenenbaum15}; the upper bound is $\mathcal P_H^{\rm frac}\subseteq\mathcal S^+$.
\end{proof}

A density-one statement for these merged runs would follow from a Hooley-type moment bound for the gaps
between squarefree sums of two squares in a residue class modulo $p$; we have not established one. For
$p\equiv1\pmod4$ the count of fractional jumps follows from \Cref{thm:jumplaw} below. Letting $H$ grow does
not help: runs of length at least $\epsilon\log N$ on a
positive proportion of shifts force $p^{H(N)}\gg_\epsilon\sqrt{\log N}$ by \Cref{thm:runs}(ii), and then the
factor in \Cref{cor:retire} is $\gg_{\epsilon,p,s}(\log N)^{s/2}$. Along the runs the arrival phases obey the affine
recursion $y_{k+1}=\fr{b^{\ell_k}y_k+z_0(m_{k+1})}$, with $\ell_k$ the run lengths, and $C^2$ is normal if and only
if the concatenated initial segments $(\fr{b^jy_k})_{j<\ell_k}$ are uniformly distributed. That is the missing
input; it concerns numerators at resolution $b^{-\ell_k}$, which \Cref{lem:barrier} shows cannot be replaced by
denominator information.

\subsection{Irrationality of the square}

Let $\mathcal S$ be the set of positive integers that are sums of two squares, zero coordinates allowed, listed as
$\sigma_1<\sigma_2<\cdots$. The following theorem is used as an external input.

\begin{theorem}[Long gaps between sums of two squares {\cite[Theorem~1]{DEKKM}}]\label{thm:dekkm}\evidence{External input}
$\displaystyle\limsup_{n\to\infty}\frac{\sigma_{n+1}-\sigma_n}{\log\sigma_n}\ge\frac{390}{449}$, with the natural logarithm.
\end{theorem}

This is Theorem~1 of the five-author paper of Dietmann, Elsholtz, Kalmynin, Konyagin and Maynard (we read the
accepted manuscript; its theorem numbering should be confirmed against the typeset version). The general
discriminant result of the same paper gives a smaller constant at discriminant $-4$, which would not suffice for
the next theorem.

\begin{theorem}[Irrational squares]\label{thm:squareirr}\evidence{Paper proof}
Let $p\nmid b$ be an odd prime, $t=\ord_p(b)$, and let $s\ge1$ satisfy $s<(390/449)\,t\log b$. Then $C(b,p,s,2)^2$ is
irrational, so its base-$b$ expansion is not eventually periodic and has infinitely many digit transitions. The
condition holds for $s=1$ and every $b$ and $p$; in particular $C(2,3,1,2)^2$ is irrational.
\end{theorem}

\begin{proof}
Write $C^2=\sum_mc(m)b^{-tm}$ with $c(m)=\sum p^{-s(v_p(i)+v_p(j))}$ over $i,j\ge1$ with $i^2+j^2=m$. Every
coordinate at a position $m'\le M$ has $v_p\le\lfloor\log_p\sqrt M\rfloor$, so the shifted prefix $R$ at the shift
$tM$ has reduced denominator dividing $p^{2s\lfloor\log_p\sqrt M\rfloor}\le M^s$. Suppose $C^2=A/Q$. The tail
$T_M=\sum_{m>M}c(m)b^{-t(m-M)}$ is positive and equals $b^{tM}A/Q-R$, a rational whose denominator divides
$QM^s$; hence $T_M\ge1/(QM^s)$, also when $R$ is an integer. Now let $M<M+g$ be consecutive elements of $\mathcal S$.
Then $c(m)=0$ for $M<m<M+g$, and $c(m)\le r_2(m)\le4\tau(m)\ll_\epsilon m^\epsilon$ (this also covers positions that have
only representations with a zero coordinate). Since consecutive squares give $g\le2\sqrt M+1$,
\[
T_M\ll_{\epsilon,b,t}b^{-tg}(M+g)^\epsilon\ll M^\epsilon b^{-tg}.
\]
By \Cref{thm:dekkm} there are infinitely many such $M$ with $g\ge(390/449-o(1))\log M$, so
$1/(QM^s)\ll M^{\epsilon-(390/449-o(1))t\log b}$. Choosing $\epsilon$ below the margin $(390/449)t\log b-s$ gives a
contradiction for large $M$. For $s=1$: $b^t\equiv1\pmod p$ and $b^t>1$ give $b^t\ge p+1\ge4$, and
$(390/449)\log4>1$ because $\log2>2/3$ (the midpoint bound for $\int_1^2dx/x$, valid by convexity). An eventually
periodic expansion would be rational, and an expansion with finitely many digit transitions is eventually
constant.
\end{proof}

The proof uses rare long gaps; no positive density of digit transitions follows. Irrationality of
$C(b,p,s,2)^2$ for $s\ge(390/449)t\log b$, and of $\xi^3$, remains open.

\subsection{The innovation law and packets of arrivals}

Let $d=r=2$ and $H\ge0$. At a position $m\in\mathcal P_H^{\rm frac}$ the shifted prefix of $Z_H$ receives the fractional
\emph{innovation} $\fr{z_H(m)}$: if $y_i$ is the phase after the $i$-th fractional arrival and $g_i$ the number of
shifts to the next one, then $y_{i+1}=\fr{b^{g_i}y_i+z_H(m_{i+1})}$. (The ordinary difference $y_{i+1}-y_i$ also contains
$(b^{g_i}-1)y_i$.) The next theorem extends \Cref{prop:fractional} to both residue classes of $p$ modulo $4$ and adds
the limiting law.

\begin{theorem}[Innovation law]\label{thm:jumplaw}\evidence{Paper proof}
Let $p\nmid b$ be an odd prime, $H\ge0$, $k=\lceil(H+1)/2\rceil$, $D=p^{s(2k-H)}$ and $a=\ord_D(2)$. Let $\mathcal F_p$ be
the set of squarefree $n\in\mathcal S^+$ with $(n/p)=-1$. Then:
\begin{enumerate}[label=\textup{(\roman*)},nosep]
\item for $n\in\mathcal F_p$, $z_H(p^{2k}n)=2^{\omega_1(n)}/D$, where $\omega_1(n)$ is the number of prime factors
$\ell\equiv1\pmod4$ of $n$; in particular $p^{2k}\mathcal F_p\subseteq\mathcal P_H^{\rm frac}$;
\item $|\mathcal F_p\cap[1,Y]|\asymp Y/\sqrt{\log Y}$, and as $n$ runs through $\mathcal F_p\cap[1,Y]$ the innovations
$\fr{z_H(p^{2k}n)}$ are asymptotically uniformly distributed on the $a$ residues $\fr{2^j/D}$, $0\le j<a$;
\item $|\mathcal P_H^{\rm frac}\cap[1,Y]|\asymp_HY/\sqrt{\log Y}$.
\end{enumerate}
\end{theorem}

\begin{proof}
(i) Let $i^2+j^2=p^{2k}n$ with $i,j\ge1$, and $u=\min(v_p(i),v_p(j))$. Since $p\nmid n$, $u\le k$. If $u<k$, then
$(i/p^u)^2+(j/p^u)^2\equiv0\pmod p$ with one of the two units, so both are units and the level is $2u\le2k-2\le H$;
such representations belong to $X_H$, for both residue classes of $p$ modulo $4$ (for $p\equiv3\pmod4$ there are
none). If $u=k$, then $i'^2+j'^2=n$ with $i=p^ki'$, $j=p^kj'$, and $p\mid i'$ would make $n$ a square modulo $p$; so the
level is exactly $2k>H$ and the contribution is $p^{-s(2k-H)}=1/D$. The ordered positive representations of the
squarefree non-square $n>1$ number $r_2(n)/4=2^{\omega_1(n)}$, which is prime to $p$.

(ii) For an $a$-th root of unity $\zeta$ let $f_\zeta(n)=\zeta^{\omega_1(n)}$ on the squarefree $n\in\mathcal S$ prime to $p$,
and $f_\zeta(n)=0$ otherwise. Its Dirichlet series has Euler factors $1+\zeta\ell^{-\sigma}$ at primes
$\ell\equiv1\pmod4$, $\ell\ne p$, a finite factor at $2$, and $1$ elsewhere, so it equals
$(\zeta(\sigma)L(\sigma,\chi_{-4}))^{\zeta/2}G_\zeta(\sigma)$ with $G_\zeta$ given by an Euler product converging absolutely for
$\Re\sigma>1/2$. The twist by $\chi_p=(\cdot/p)$ equals $(L(\sigma,\chi_p)L(\sigma,\chi_p\chi_{-4}))^{\zeta/2}G'_\zeta(\sigma)$ with both
characters nonprincipal. By the Selberg--Delange method \cite[Chapter~II.5]{Tenenbaum15}, $\sum_{n\le Y}f_1(n)\sim
c\,Y/\sqrt{\log Y}$ with $c>0$, $\sum_{n\le Y}f_\zeta(n)=O(Y(\log Y)^{\Re\zeta/2-1})=o(Y/\sqrt{\log Y})$ for $\zeta\ne1$, and every
twisted sum is $o(Y/\sqrt{\log Y})$, since its zeta exponent is $0$. For $p\nmid n$ the indicator of $(n/p)=-1$ is
$(1-\chi_p(n))/2$, so by finite Fourier inversion
\[
\#\{n\in\mathcal F_p\cap[1,Y]:\omega_1(n)\equiv j\pmod a\}=\frac1{2a}\sum_{n\le Y}f_1(n)+o\Bigl(\frac Y{\sqrt{\log Y}}\Bigr)
\]
uniformly in $j$. Since $j\mapsto2^j\bmod D$ is injective on $\Z/a\Z$, (i) gives the claim.

(iii) The lower bound follows from (i) and (ii); the upper bound is $\mathcal P_H^{\rm frac}\subseteq\mathcal S^+$ and
Landau's theorem \cite{Landau08}.
\end{proof}

The denominators of the actual phases at arrivals are determined in relative density.

\begin{proposition}[Aggregate denominator law]\label{prop:aggregate}\evidence{Paper proof}
Let $d=r=2$, fix $H\ge0$, let $m_1<m_2<\cdots$ enumerate $\mathcal P_H^{\rm frac}$, let $y_i$ be the phase of the shifted
prefix of $Z_H$ at the shift $tm_i$, and $e_i=v_p(\den y_i)$. Let $M=M(Y)\ge1$ satisfy $\log M=o(\log Y)$. For a
proportion tending to one of the indices $i$ with $Y/2\le m_i\le Y$, simultaneously for $1\le j\le M$,
\[
v_p\bigl(\den(y_{i+j}-y_i)\bigr)=e_i-v_p(b^t-1)-v_p(m_{i+j}-m_i).
\]
In particular $\log\den(y_{i+j}-y_i)/\log Y\to s$ in relative density, uniformly for $1\le j\le M$.
\end{proposition}

\begin{proof}
By \Cref{thm:jumplaw}, differencing its asymptotics at the endpoints, there are $\gg_HY/\sqrt{\log Y}$ arrivals with
$Y/2\le m_i\le Y$. Fix $L$. By \Cref{thm:smallball} with $d=r=2$, the frozen set on which $M_2\equiv0\pmod{p^L}$ has
finitely many component intervals in each normalized shell and measure $O((\log L)^{-2})$; on every other cell,
freezing (\Cref{lem:freezer}) and \Cref{prop:filtration} give $e_i\ge s\log_pY-O_{H,L}(1)$. The relevant shifts meet
$O_p(1)$ shells. In the position variable the bad cells have total length $O(Y(\log L)^{-2})$, and differencing
Landau's asymptotic at their endpoints bounds the sums of two squares in them by
$O((\log L)^{-2}Y/\sqrt{\log Y})+o_L(Y/\sqrt{\log Y})$. Letting $Y\to\infty$ and then $L\to\infty$, and using the
elementary bound $\den y_i\le Y^s$, gives $e_i/\log_pY\to s$ in relative density at the arrivals.

Fix $\delta>0$ with $3\delta<s$ and omit the last $M$ indices before $Y$. Starts with $m_{i+M}-m_i>Y^\delta$ number at
most $MY^{1-\delta}$ by telescoping. A position whose innovation has denominator at least $Y^\delta$ carries a pair of
level $h\ge(\delta/s)\log_pY$; a fixed pair of valuations occurs at $O(Yp^{-h})$ positions up to $Y$, so there are
$O(Y^{1-\delta/s}(1+\log Y))$ such positions, each spoiling at most $M$ starts. Both costs are negligible on the scale
$Y/\sqrt{\log Y}$. For the other starts the propagated innovation $y_{i+j}-b^{t(m_{i+j}-m_i)}y_i$ has denominator below
$Y^\delta$, since the denominators are nested powers of $p$, while by lifting the exponent the term
$(b^{t(m_{i+j}-m_i)}-1)y_i$ has $p$-exponent exactly $e_i-v_p(b^t-1)-v_p(m_{i+j}-m_i)$, which exceeds $\delta\log_pY$
with relative probability tending to one. So no cancellation occurs. Letting $\delta\to0$ gives the limit.
\end{proof}

A denominator tending to infinity does not bound the distance of a phase change from $0$ on the circle, and
\Cref{prop:aggregate} says nothing about digits. The next two statements show that exact jumps over a
logarithmic packet cannot say more. For a series $\sum_mc_mb^{-m}$ with rational coefficients write
$R_n=\sum_{m\le n}c_mb^{n-m}$, $T_n=\sum_{m>n}c_mb^{n-m}$ and $V_j=\sum_{i=1}^jb^{j-i}c_{n+i}$, so that
$R_{n+j}=b^jR_n+V_j$ and $T_{n+j}=b^jT_n-V_j$.

\begin{lemma}[Packets]\label{lem:packet}\evidence{Paper proof}
Let $p$ be an odd prime, suppose the reduced denominator of $R_n$ is $q=p^e>1$, and suppose that each coefficient
$c_{n+i}$, $1\le i\le B$, has $p$-power denominator and lies in $p^{1-e}\Z_{(p)}$. Replacing $\fr{R_n}$ by any $a/q$ with
$p\nmid a$, and keeping the coefficients $c_{n+i}$ and the future tails $T_{n+j}$ fixed, preserves every reduced
denominator of $R_{n+j}$, $j\le B$. If $2b^B\le q$, every base-$b$ word of length $B$ is the word of digits
$n+1,\ldots,n+B$ of some member of this family.
\end{lemma}

\begin{proof}
Each $V_j$ has $p$-adic valuation at least $1-e$, so $b^ja/q+V_j$ has reduced denominator exactly $p^e$, as does
$R_{n+j}$. The modified real number satisfies $b^ja/q+V_j+T_{n+j}=b^j(a/q+T_n)$, so its digits $n+1,\ldots,n+B$ are the
leading word of $\fr{a/q+T_n}$. The residues $a/q$ with $p\nmid a$ have maximal spacing $2/q\le b^{-B}$ on the circle,
since $p\ge3$, so every translated half-open cylinder of length $b^{-B}$ contains one.
\end{proof}

\begin{proposition}[Eligible packets]\label{prop:packets}\evidence{Paper proof}
Let $d=r\ge2$, $0<\beta<s$ and $B=\lfloor\beta\log_bN\rfloor$, and apply \Cref{lem:packet} to the shifted prefixes of $C^d$.
Then all but $o(N)$ starts $n\in[N/2,N]$ satisfy its hypotheses with $2b^B\le q$. The same holds for $Z_H$ for every
fixed $H$.
\end{proposition}

\begin{proof}
Fix $\beta<\beta'<s$. On $[N/2,N]$ one has $dsJ(n)=s\log_pN+O(1)$, so \Cref{prop:depthtight} (first with a fixed depth
defect, then letting it grow) gives $q\ge N^{\beta'}\ge2b^B$ outside $o(N)$ shifts; for $Z_H$ the loss of at most
$p^{sH}$ is absorbed by a slightly larger intermediate exponent. A coefficient with pole order at least
$\beta'\log_pN$ requires a tuple of total level $h\ge(\beta'/s)\log_pN$. Through $2N$ a fixed vector of valuations
$(h_1,\ldots,h_d)$ occurs at most $\prod_c(2N/t)^{1/d}p^{-h_c}=O(Np^{-h})$ times, and there are $O((1+h)^{d-1})$ such
vectors, so these coefficients number $O(N^{1-\beta'/s}(1+\log N)^{d-1})$; each spoils at most $B$ starts, which costs
$o(N)$.
\end{proof}

Thus, at density one of starts, the exact jumps and every reduced prefix denominator throughout a packet of
length $\lfloor\beta\log_bN\rfloor$ admit every word on that packet when the incoming numerator varies, for squares and
for the cube $\xi^3$ alike, with arbitrarily many intervening arrivals. The statement concerns hypothetical
continuations: it does not assert that the alternative numerators occur in the given series. It explains why an
input on the \emph{actual} incoming numerators is needed. For instance, if for some fixed $\epsilon,\eta>0$ every
long prefix, divided into packets with starting phases $\lambda_n=\fr{R_n+T_n}$ and lengths $m_n$, satisfied
$\sum_n\#\{0\le j<m_n:\fr{b^j\lambda_n}\notin I\}\ge\eta N$ for every arc $I$ of length $\epsilon$, then
$\limsup_ND_N\le1-\min(\epsilon,\eta)$.

\subsection{The cubic boundary}

For $d=r=3$ the identity of \Cref{prop:filtration} holds, with $\sum_{m\le Y}a_H(m)\le3^HV_{3,3}Y$ for
$C(2,3,1,3)^3=\xi^3$. \Cref{lem:sparseint} would give $D_N(X_H)\to1$ if the sums of three positive cubes had
density zero. They number at least $Y^{0.91709477}$ up to $Y$ \cite{Wooley15} and at most $(V_{3,3}/6+o(1))Y\approx
0.1187Y$, and the model of Deshouillers, Hennecart and Landreau \cite{DHL06} predicts density about $0.0999$. If
that density is positive, the support of $X_H$ has positive lower density for large $H$, the transfer says
nothing, and the mean run length of $Z_H$ is $O_H(1)$ while $\log_2q_n\asymp\log N$ on a set of density close to
one (\Cref{prop:depthtight}). At $d=2$ the integer part is provably negligible and the obstruction is the
length of the runs; at $d=3$ even the integer part is not known to be negligible. Residues modulo $9$ give exact
information on the cube.

\begin{proposition}[The cube modulo $9$]\label{prop:cubic9}\evidence{Paper proof}
Let $p=3$, $d=r=3$, $s\ge1$ and $H\ge0$, with $z_H$ as in \Cref{prop:filtration}.
\begin{enumerate}[label=\textup{(\roman*)},nosep]
\item If $m\equiv\pm3\pmod9$, every triple at $m$ consists of units with cubes of the same sign modulo $9$, and $z_H(m)=0$.
\item If $m\equiv\pm4\pmod9$, there are no triples at $m$.
\item If $m\equiv\pm2\pmod9$, every triple at $m$ has exactly one coordinate divisible by $3$, and for each level the
number of triples at $m$ is divisible by $3$.
\item If $3a>H$, then $3^{s(3a-H)}z_H(3^{3a+1})\equiv1\pmod9$; the reduced denominator of $z_H(3^{3a+1})$ is exactly
$3^{s(3a-H)}$.
\end{enumerate}
\end{proposition}

\begin{proof}
Cubes modulo $9$ lie in $\{0,\pm1\}$, and $3\mid i$ exactly when $i^3\equiv0\pmod9$. This gives (i), where all levels are
$0\le H$, and (ii). In (iii) the residues are $\{1,1,0\}$ or $\{-1,-1,0\}$; cyclic permutation of the slots preserves
the level and has no fixed triple. In (iv), $i^3+j^3+k^3=3^{3a+1}$ forces $ijk\le3^{3a}$ by the arithmetic--geometric
mean inequality, with equality only for $i=j=k=3^a$, of level $3a>H$. Every other triple has $v_3(ijk)\le3a-1$, and
these fall into cyclic orbits of size three; after multiplication by $3^{s(3a-H)}$ an orbit of level $\ell\le3a-1$
contributes $3\cdot3^{s(3a-\ell)}\equiv0\pmod9$.
\end{proof}

So forced innovations of unbounded depth occur at sparse positions, alongside the packet obstruction of
\Cref{prop:packets}. A long gap would again give irrationality.

\begin{proposition}[A conditional irrationality criterion]\label{prop:cubiccond}\evidence{Paper proof; conditional}
If consecutive sums $M<M+g$ of three positive cubes satisfy $\limsup g/\log M>2/(3\log2)$, then $\xi^3$ is irrational.
\end{proposition}

\begin{proof}
Here $b=2$, $p=3$, $t=2$, $s=1$. The shifted prefix of $\xi^3$ at the shift $2M$ has reduced denominator dividing
$3^{3\lfloor\log_3M^{1/3}\rfloor}\le M$. To bound a coefficient at $m$, fix the first of the $O(m^{1/3})$ coordinates; for
$j^3+k^3=B>0$ the divisor $j+k$ of $B$ determines $jk$, which gives at most $2\tau(B)$ ordered pairs. So the
coefficients are $O_\epsilon(m^{1/3+\epsilon})$ and the tail after the shift $2M$ is $O_\epsilon(2^{-2g}(M+g)^{1/3+\epsilon})$, with
$M+g\ll M$ because of the values $i^3+1+1$. As in \Cref{thm:squareirr}, a rational $\xi^3=A/Q$ would make the tail at
least $1/(QM)$, which fails along the indicated subsequence because $(1+1/3)/(2\log2)=2/(3\log2)$.
\end{proof}

The hypothesis is open. Ghidelli's theorem \cite[Theorem~1.1]{Ghidelli20} gives, for every positive-definite
diagonal cubic form in three variables, runs of non-values of length of order $\sqrt{\log N}/(\log\log N)^2$ below $N$
with nonnegative variables, hence also for sums of three positive cubes; its stated bound does not reach the
logarithmic constant required here.

\begin{remark}[Finite measurements]\label{rem:finitedata}
In the first $2^{16}$ binary digits of $C(2,3,1,2)^2$, $C(2,3,2,2)^2$ and $\xi^3=C(2,3,1,3)^3$, computed with
certified interval arithmetic, every block of length at most four has frequency within $0.0021$ of its
expected value, and $D_{2^{16}}\le0.0037$. With $H=0$, $b=2$, $s=1$ and unscaled positions $m\le4096$, the square
$C(2,p,1,2)^2$ has $643$ fractional arrivals for $p=3$ and $479$ for $p=5$, and $\xi^3$ has $164$; for $n\le400$ in
$\mathcal F_p$ with $H=1$ the innovation counts by numerator modulo $D$ are $(1{:}9,2{:}29)$ for $p=3$ and
$(1{:}1,2{:}30,4{:}1)$ for $p=5$, far from the limiting law of \Cref{thm:jumplaw}, as expected from convergence
at scale $\log\log Y$. These are finite measurements; they are consistent with normality, are not evidence for
the limiting laws, and prove nothing asymptotic. The scripts are listed in \Cref{sec:formal}.
\end{remark}

\section{Questions}\label{sec:questions}

\begin{question}\label{q:boundary}
Is $\xi^3=C(2,3,1,3)^3$ normal to base $2$? Is $C(b,p,s,2)^2$ normal to base $b$? For either number, is
$\limsup_ND_N<1$, or at least $D_N\not\to1$? Are their expansions disjunctive, and what are their asymptotic
digit-transition frequencies? (\Cref{thm:D} gives the rate $(\log\log N)^{-3/2}$ for $\xi^2$;
\Cref{sec:boundary} explains why denominator information, even exact jumps over logarithmic packets, cannot
settle these questions.)
\end{question}

\begin{question}\label{q:irrational}
Is $\xi^3$ irrational? Is $C(b,p,s,2)^2$ irrational when $s\ge(390/449)\,t\log b$? (\Cref{thm:squareirr} covers
smaller $s$, and \Cref{prop:cubiccond} reduces the first question to a logarithmic gap between sums of three
positive cubes with constant above $2/(3\log2)$.)
\end{question}

\begin{question}
Is $\xi$ transcendental? Is it a Liouville number? (Normality excludes rationality, and \Cref{thm:D}
excludes algebraic degree at most two, that is, less than $d=3$.)
\end{question}

\begin{question}
What is the true order of $D_N(P(C(b,p,s,d)))$ for $\deg P\ge2$, and of the several-order discrepancies of
\Cref{thm:severalrate}? Is it a power of $N$? A local form of the competitor counts of
\Cref{lem:avgiso,lem:wiso} near sums of fewer coordinates would give a power of $\log N$. Is the exponent $1/3$
in \Cref{thm:B,thm:C} optimal for these particular numbers?
\end{question}

\section{Formal verification and data availability}\label{sec:formal}

The Lean~4 \cite{Lean4} declarations cited above, built with Mathlib \cite{Mathlib}, are in the
repository \url{https://github.com/eldernos/MathematicsResearch}. Links to the modules
\ttt{RetirementDenominator}, \ttt{Phase18Synthetic}, \ttt{SharpTranslation}, \ttt{SmoothOrbit} and
\ttt{PositiveSummation} are pinned to the commit
\begin{center}\ttt{\leancommit},\end{center}
links to \ttt{SurvivingGrid} and \ttt{PrimePowerGrid} to the later commit
\begin{center}\ttt{\phaseXIXcommit},\end{center}
and links to \ttt{PowerSumPeeling} to the commit
\begin{center}\ttt{\phaseXXIcommit}.\end{center}
At each commit the full build completed and the axiom audit accepted only \ttt{propext},
\ttt{Classical.choice} and \ttt{Quot.sound}, and an independent kernel replay covered the modules used
here. The theorem verification map, the novelty ledger, the list of changes from version~1.0 and the
reproduction commands are in the files \ttt{VERIFICATION.md}, \ttt{LEDGER.md} and \ttt{CHANGES.md} next to
the source of this paper (directory \ttt{papers/07-synthetic-expansions-v2} of the repository); version~1.0
remains in \ttt{papers/07-synthetic-expansions}.

The kernel-verified statements are
\Cref{prop:group,thm:den,lem:short,thm:switch,lem:program,thm:sharpcited}; the last is quoted from
\cite{CompanionH}. The Lean declarations used in
\Cref{lem:rjump,lem:fB,lem:avgiso,lem:wiso} are finite arithmetic and combinatorial cores:
the exact jump of a truncated mass at an isolated centre, the degree and leading coefficient of $\Phi_q$, the
dominant-loss peeling lemma and the divisor bound for differences of $d$-th powers. None of them states an exponential-sum estimate,
a measure-theoretic statement, a normality conclusion, a rate or a Diophantine property. All rates
(\Cref{thm:example,thm:rate,thm:program,thm:Y,thm:effective,thm:severalrate}),
the limit-set theorems, \Cref{prop:criterion} beyond the short-orbit estimate, the passage from repunits to
reduced raw coefficients after \Cref{prop:group}, and \Cref{sec:quadratic,sec:sums,sec:boundary}
beyond the named Lean steps have written proofs. In particular \Cref{thm:quadratic,thm:effective,thm:severalrate}
and \Cref{cor:flagship,cor:flagship2} are paper proofs, and \Cref{thm:smallball,thm:higher,thm:several,lem:wiso} are
Lean-supported paper proofs whose Lean steps are only the finite cores named above; none of these statements,
and none of their analytic or measure-theoretic steps, is formalized. Exact rational computations exercising these statements are
\repo{scripts/phase18_principal.py}{\ttt{scripts/phase18\_principal.py}} and
\repo{scripts/phase18_programming.py}{\ttt{scripts/phase18\_programming.py}} at the first commit, and
\href{https://github.com/eldernos/MathematicsResearch/tree/\phaseXIXcommit/scripts/phase19_grids.py}{\ttt{scripts/phase19\_grids.py}}
and
\href{https://github.com/eldernos/MathematicsResearch/tree/\phaseXIXcommit/scripts/phase19_nonlinear.py}{\ttt{scripts/phase19\_nonlinear.py}}
at the second, and
\href{https://github.com/eldernos/MathematicsResearch/tree/\phaseXXIcommit/scripts/phase20_normality.py}{\ttt{scripts/phase20\_normality.py}}
and
\href{https://github.com/eldernos/MathematicsResearch/tree/\phaseXXIcommit/scripts/phase21_nonlinear.py}{\ttt{scripts/phase21\_nonlinear.py}}
at the third. The scripts \ttt{scripts/phase22\_short\_orbits.py} (for \Cref{sec:uniformorbit}),
\ttt{scripts/phase22\_nonlinear\_boundary.py} (for \Cref{sec:boundary} and \Cref{rem:finitedata}) and
\ttt{scripts/phase23\_nonlinear\_short\_orbits.py} (for the freezing fixture after \eqref{eq:freezeJ}, tied
multivariate residues, \Cref{prop:aggregate,lem:packet,prop:cubic9}, the order computations and
\Cref{rem:finitedata}) accompany this version in the same repository. All of them
are finite checks and illustrations, not proofs of the asymptotic statements; the exception is the finite
certificate in \Cref{prop:nouniform}: \ttt{phase21\_nonlinear.py} certifies the net exactly for
$H\in\{6,10,100,1000,4000\}$ together with $R_{4000}^{126}<3^{4000}$. The enclosures of $N^{1/3}D_N(\xi)$ quoted after \Cref{thm:example}
are computed by \ttt{checks/xi\_discrepancy.py} in the directory of version~1.0 (standard library only).

\section{Use of AI tools}

This manuscript was produced in an AI-assisted research project. Large language models proposed
conjectures and proof strategies, wrote the proofs and the text, implemented the exact computations
and Lean~4 formalizations, searched the literature and reviewed drafts across several model
providers. The repository owner directed the project and is responsible for its release. The Lean
kernel checked the declarations listed in the verification map. The written proofs, literature
comparisons and priority statements have been checked only by AI models from different providers.
They have not been checked by an independent human expert, and no human specialist has read the proofs
new in version~2.0. The material added after the first review of this draft (\Cref{lem:wiso},
\Cref{thm:severalrate,thm:squareirr,thm:jumplaw}, \Cref{prop:aggregate,prop:packets,prop:cubic9,prop:cubiccond}
and \Cref{lem:packet}) rests on a source note that received one cross-provider model review, repaired and
re-read; the assembled manuscript has not been re-reviewed as a whole. Cross-model review is not independent
human peer review.

\bibliographystyle{plainurl}
\bibliography{references}
\end{document}
